% Lesson 21 — Grammar: Adjectives as nouns; "So do I", "Neither did he" — Anglais 1ère A — Cours signé OnBuch+
% Programme officiel MINESEC. Compiler avec : tectonic ENA21-grammar-adjectives-as-nouns-so-do-i-neither-did-he.tex
\newcommand{\DOCMATIERE}{Anglais}
\newcommand{\DOCNIVEAU}{1ère A}
\newcommand{\DOCMODULE}{Chapter 2 : Economic Life and Occupations}
\newcommand{\DOCLECON}{ENA21}
\newcommand{\DOCTITRE}{Lesson 21 — Grammar: Adjectives as nouns; "So do I", "Neither did he"}
% OnBuch+ — préambule « cours signés » ENRICHI (compilé avec tectonic / XeLaTeX)
% Système visuel « Pop » d'OnBuch+ : fond crème (jamais blanc), bordures encre
% épaisses, ombres « dures » décalées, titres Archivo Black en capitales.
% Version MATHÉMATIQUES : graphes (pgfplots/tikz), boîtes pédagogiques
% (définition, propriété, méthode, attention, le savais-tu, exercice résolu,
% exercice, TP, bilan), page de garde et signature OnBuch+ ancrée en bas.
%
% Macros attendues dans main.tex AVANT \input{preamble} :
%   \DOCMATIERE  \DOCNIVEAU  \DOCMODULE  \DOCLECON (numéro)  \DOCTITRE
\documentclass[11pt]{article}

\usepackage{fontspec}
\usepackage[english,french]{babel} % français = langue principale ; l'anglais via \eng{..} / textebox
\usepackage[a4paper,margin=2cm,top=3.4cm,bottom=2.8cm,headheight=1.4cm,footskip=1.3cm]{geometry}
\usepackage{amsmath,amssymb,amsfonts}
\usepackage{mathtools} % \xmapsto, \coloneqq... souvent utilisés par les modèles
\usepackage{cancel} % simplifications barrées (bilans de forces)
% \abs{z} (module, valeur absolue) et \norm{u} (norme), utilisés par les modèles
\providecommand{\abs}{}\let\abs\relax\DeclarePairedDelimiter{\abs}{\lvert}{\rvert}
\providecommand{\norm}{}\let\norm\relax\DeclarePairedDelimiter{\norm}{\lVert}{\rVert}
\usepackage[table]{xcolor} % [table] : fournit aussi \rowcolors (alternance de lignes)
\usepackage{pagecolor}
\usepackage{tikz}
\usetikzlibrary{arrows.meta,positioning,calc,shapes.geometric,shapes.misc,shapes.arrows,decorations.pathmorphing,decorations.pathreplacing,decorations.markings,patterns,shadows,mindmap,trees,fit,backgrounds,matrix,intersections,angles,quotes}
\usepackage{pgfplots}
\usepackage[version=4]{mhchem} % équations nucléaires \ce{^{235}_{92}U}
\usepackage[european,siunitx]{circuitikz} % schémas électriques
\pgfplotsset{compat=1.18}
\usepgfplotslibrary{fillbetween,groupplots} % aires entre courbes (calcul intégral)
\usepackage{enumitem}
\usepackage{fancyhdr}
% [most] charge toutes les bibliothèques tcolorbox (listings, minted, raster,
% documentation...) et gonfle inutilement la mémoire TeX sur un document long
% (~300 boîtes) : on ne charge que ce qui est réellement utilisé.
\usepackage[skins,breakable]{tcolorbox}
\usepackage{graphicx}
\usepackage{colortbl}
\usepackage{booktabs}
\usepackage{tabularx}
\usepackage{diagbox} % case d'en-tête diagonale des tableaux à double entrée
% Colonnes extensibles centrée / gauche / droite (C, L, R), souvent utilisées par les modèles
\newcolumntype{C}{>{\centering\arraybackslash}X}
\newcolumntype{L}{>{\raggedright\arraybackslash}X}
\newcolumntype{R}{>{\raggedleft\arraybackslash}X}
\usepackage{array}
\usepackage{multicol}
\usepackage{multirow}
\usepackage{siunitx}
% parse-numbers=false : accepte aussi \SI{2,5 \times 10^{-3}}{...}
\sisetup{locale=FR,output-decimal-marker={,},group-minimum-digits=4,parse-numbers=false}
\DeclareSIUnit\ohm{\ensuremath{\symup{\Omega}}} % U+2126 absent de la police
\DeclareSIUnit\division{div} % réglages d'oscilloscope (ms/div, V/div)
\DeclareSIUnit\tr{tr} % tours (vitesse de rotation, ex. \SI{1500}{\tr\per\minute})
\DeclareSIUnit\molar{M} % concentration molaire (ex. \SI{3}{\molar}, \SI{1}{\micro\molar})
\DeclareSIUnit\an{an} % année (le modèle confond parfois avec \year, un compteur TeX) — ex. \SI{5}{\kilo\meter\per\an}
\DeclareSIUnit\FCFA{FCFA} % franc CFA (ex. \SI{500}{\FCFA\per\kilo\gram})
\DeclareSIUnit\degreeBrix{\textdegree Brix} % teneur en sucre d'un sirop/jus (réfractométrie)
\DeclareSIUnit\month{mois} % le modèle utilise parfois \month, un compteur TeX, comme unité de durée
\usepackage{qrcode}
% \degree : commande fréquemment utilisée par les modèles pour un angle ou une
% température en dehors de \SI{}{} (non fournie par les paquets chargés).
\providecommand{\degree}{^{\circ}}
% \qed (fin de démonstration), utilisé par les modèles sans amsthm
\providecommand{\qed}{\hfill$\square$}
% \barre : double barre des valeurs interdites dans un tableau de variations (tkz-tab)
\providecommand{\barre}{\ensuremath{\|}}
% environnement proof (amsthm non chargé) utilisé par les modèles
\ifdefined\proof\else\newenvironment{proof}[1][Démonstration]{\par\noindent\textit{#1.}\ }{\qed\par}\fi
\usepackage{lastpage}
\usepackage{hyperref}
\hypersetup{hidelinks}

% ============================================================
% Palette « Pop » OnBuch+ — mêmes hex que la classe P (plus_ui.dart)
% ============================================================
\definecolor{poporange}{HTML}{F59321}
\definecolor{popdark}{HTML}{E07A0C}
\definecolor{popcream}{HTML}{FFF6E8}
\definecolor{popink}{HTML}{1C1714}
\definecolor{poppurple}{HTML}{7A5AE0}
\definecolor{popgreen}{HTML}{2E9E6A}
\definecolor{poppink}{HTML}{E0457B}
\definecolor{popblue}{HTML}{2F7DE1}
\definecolor{popgold}{HTML}{C98A12}
\definecolor{popmuted}{HTML}{8A7F76}
\definecolor{popline}{HTML}{EADFD2}
\definecolor{popgray}{HTML}{8A7F76}
\colorlet{popgrey}{popgray}
% Alias tolérants (le modèle nomme parfois les couleurs différemment)
\colorlet{popwhite}{white}
\colorlet{popblack}{popink}
\colorlet{popred}{poppink}
\colorlet{popyellow}{popgold}
% Teintes claires (fonds de tableaux/encadrés) — le modèle nomme parfois
% les couleurs avec un suffixe L (ex. popblueL) sans les avoir définies.
\colorlet{poporangeL}{poporange!20}
\colorlet{popdarkL}{popdark!20}
\colorlet{poppurpleL}{poppurple!20}
\colorlet{popgreenL}{popgreen!20}
\colorlet{poppinkL}{poppink!20}
\colorlet{popblueL}{popblue!20}
\colorlet{popgoldL}{popgold!20}
\colorlet{popredL}{popred!20}
\colorlet{popgrayL}{popgray!20}
% Teintes claires pour fonds de tableaux / zones
\colorlet{popblueL}{popblue!10!white}
\colorlet{poporangeL}{poporange!14!white}
\colorlet{popgreenL}{popgreen!12!white}
\colorlet{poppurpleL}{poppurple!12!white}
\colorlet{poppinkL}{poppink!12!white}
\colorlet{popgoldL}{popgold!14!white}

\pagecolor{popcream}

% ============================================================
% Typographie
% ============================================================
\defaultfontfeatures{Ligatures=TeX}
\setmainfont{PlusJakartaSans-Medium.ttf}[Path=../../../fonts/,BoldFont=PlusJakartaSans-Bold.ttf]
% Symboles, API et emojis absents de la police principale : repli sur DejaVu Sans (caractères actifs)
\newfontface\symfont{DejaVuSans.ttf}[Path=../../../fonts/]
\makeatletter
\newcommand\symactive[1]{\catcode#1=13 \begingroup\lccode`\~=#1 \lowercase{\endgroup\def~}{{\symfont\char#1}}}
\newcommand\symblank[1]{\catcode#1=13 \begingroup\lccode`\~=#1 \lowercase{\endgroup\def~}{}}
\newcommand\symhyph[1]{\catcode#1=13 \begingroup\lccode`\~=#1 \lowercase{\endgroup\def~}{\mbox{-}}}
\newcount\symc
\newcommand\symrange[2]{\symc=#1 \@whilenum\symc<#2 \do{\expandafter\symactive\expandafter{\the\symc}\advance\symc by 1 }}
\newcommand\symblankrange[2]{\symc=#1 \@whilenum\symc<#2 \do{\expandafter\symblank\expandafter{\the\symc}\advance\symc by 1 }}
\makeatother
\symrange{"0250}{"02FF}\symactive{"2713}\symactive{"2714}\symactive{"2717}\symactive{"2718}\symactive{"2610}\symactive{"2611}\symactive{"2612}
\symhyph{"2011}\symhyph{"2E1A}\symblankrange{"1F300}{"1FAFF}\symblank{"FE0F}
% Maths en sans-serif (Fira Math) avec chiffres et lettres droites de Plus
% Jakarta, pour que les formules se fondent dans le texte.
\usepackage[math-style=ISO,bold-style=ISO]{unicode-math}
\setmathfont{FiraMath-Regular.otf}[Path=../../../fonts/]
\setmathfont{PlusJakartaSans-Medium.ttf}[Path=../../../fonts/,range={up/num,up/{Latin,latin},\mathup}]
\usepackage{icomma} % virgule décimale sans espace en mode math
% Caractères mathématiques Unicode absents des polices de texte (Plus Jakarta,
% Archivo Black) : rendus par leur équivalent mathématique, en texte comme en
% formule — sinon ils s'affichent en carré vide (ex. « Arithmétique dans ℤ »).
\usepackage{newunicodechar}
\newunicodechar{ℝ}{\ensuremath{\mathbb{R}}}
\newunicodechar{ℤ}{\ensuremath{\mathbb{Z}}}
\newunicodechar{ℂ}{\ensuremath{\mathbb{C}}}
\newunicodechar{ℕ}{\ensuremath{\mathbb{N}}}
\newunicodechar{ℚ}{\ensuremath{\mathbb{Q}}}
\newunicodechar{𝔻}{\ensuremath{\mathbb{D}}}
\newunicodechar{α}{\ensuremath{\alpha}}
\newunicodechar{β}{\ensuremath{\beta}}
\newunicodechar{γ}{\ensuremath{\gamma}}
\newunicodechar{δ}{\ensuremath{\delta}}
\newunicodechar{ε}{\ensuremath{\varepsilon}}
\newunicodechar{θ}{\ensuremath{\theta}}
\newunicodechar{λ}{\ensuremath{\lambda}}
\newunicodechar{μ}{\ensuremath{\mu}}
\newunicodechar{σ}{\ensuremath{\sigma}}
\newunicodechar{τ}{\ensuremath{\tau}}
\newunicodechar{φ}{\ensuremath{\varphi}}
\newunicodechar{ψ}{\ensuremath{\psi}}
\newunicodechar{ω}{\ensuremath{\omega}}
\newunicodechar{Σ}{\ensuremath{\Sigma}}
% majuscules produites par \MakeUppercase dans les titres (α-aminés -> Α)
\newunicodechar{Α}{\ensuremath{\alpha}}
\newunicodechar{Β}{\ensuremath{\beta}}
\newunicodechar{Γ}{\ensuremath{\Gamma}}
\newunicodechar{Δ}{\ensuremath{\Delta}}
\newunicodechar{Ε}{\ensuremath{\varepsilon}}
\newunicodechar{Θ}{\ensuremath{\Theta}}
\newunicodechar{Λ}{\ensuremath{\Lambda}}
\newunicodechar{Μ}{\ensuremath{\mu}}
\newunicodechar{Τ}{\ensuremath{\tau}}
\newunicodechar{Φ}{\ensuremath{\Phi}}
\newunicodechar{Ψ}{\ensuremath{\Psi}}
\newunicodechar{Ω}{\ensuremath{\Omega}}
\newunicodechar{↦}{\ensuremath{\mapsto}}
\newunicodechar{∈}{\ensuremath{\in}}
\newunicodechar{∉}{\ensuremath{\notin}}
\newunicodechar{∧}{\ensuremath{\wedge}}
\newunicodechar{⇔}{\ensuremath{\Leftrightarrow}}
\newunicodechar{⟺}{\ensuremath{\Longleftrightarrow}}
\newunicodechar{⇒}{\ensuremath{\Rightarrow}}
\newunicodechar{⟹}{\ensuremath{\Longrightarrow}}
\newunicodechar{∘}{\ensuremath{\circ}}
\newunicodechar{∪}{\ensuremath{\cup}}
\newunicodechar{∩}{\ensuremath{\cap}}
\newunicodechar{≡}{\ensuremath{\equiv}}
\newunicodechar{⊂}{\ensuremath{\subset}}
\newunicodechar{⊥}{\ensuremath{\perp}}
\newunicodechar{∃}{\ensuremath{\exists}}
\newunicodechar{∀}{\ensuremath{\forall}}
\newunicodechar{∛}{\ensuremath{\sqrt[3]{\,}}}
\newunicodechar{∜}{\ensuremath{\sqrt[4]{\,}}}
\newunicodechar{⃗}{\ensuremath{{}^{\to}}}
\newunicodechar{ⁿ}{\ifmmode^{n}\else\textsuperscript{\ensuremath{n}}\fi}
\newunicodechar{ˣ}{\ifmmode^{x}\else\textsuperscript{\ensuremath{x}}\fi}
\newunicodechar{ᵇ}{\ifmmode^{b}\else\textsuperscript{\ensuremath{b}}\fi}
\newunicodechar{ʳ}{\ifmmode^{r}\else\textsuperscript{\ensuremath{r}}\fi}
\newunicodechar{ᵘ}{\ifmmode^{u}\else\textsuperscript{\ensuremath{u}}\fi}
\newunicodechar{ʸ}{\ifmmode^{y}\else\textsuperscript{\ensuremath{y}}\fi}
\newunicodechar{ᵃ}{\ifmmode^{a}\else\textsuperscript{\ensuremath{a}}\fi}
\newunicodechar{ᵉ}{\ifmmode^{e}\else\textsuperscript{\ensuremath{e}}\fi}
\newunicodechar{ᵏ}{\ifmmode^{k}\else\textsuperscript{\ensuremath{k}}\fi}
\newunicodechar{ᵖ}{\ifmmode^{p}\else\textsuperscript{\ensuremath{p}}\fi}
\newunicodechar{ᴺ}{\ifmmode^{N}\else\textsuperscript{\ensuremath{N}}\fi}
\newunicodechar{ᶜ}{\ifmmode^{c}\else\textsuperscript{\ensuremath{c}}\fi}
\newunicodechar{ʰ}{\ifmmode^{h}\else\textsuperscript{\ensuremath{h}}\fi}
\newunicodechar{ᵠ}{\ifmmode^{q}\else\textsuperscript{\ensuremath{q}}\fi}
\newunicodechar{ₙ}{\ifmmode_{n}\else\textsubscript{\ensuremath{n}}\fi}
\newunicodechar{ₐ}{\ifmmode_{a}\else\textsubscript{\ensuremath{a}}\fi}
\newunicodechar{ᵢ}{\ifmmode_{i}\else\textsubscript{\ensuremath{i}}\fi}
\newunicodechar{ᵣ}{\ifmmode_{r}\else\textsubscript{\ensuremath{r}}\fi}
\newunicodechar{ₖ}{\ifmmode_{k}\else\textsubscript{\ensuremath{k}}\fi}
\newunicodechar{ᵦ}{\ifmmode_{\beta}\else\textsubscript{\ensuremath{\beta}}\fi}
\newunicodechar{ₓ}{\ifmmode_{x}\else\textsubscript{\ensuremath{x}}\fi}
\newfontfamily\popdisplay{ArchivoBlack-Regular.ttf}[Path=../../../fonts/]
\newfontfamily\poplogo{SpaceGrotesk-Bold.ttf}[Path=../../../fonts/]
\newfontfamily\popsemi{PlusJakartaSans-SemiBold.ttf}[Path=../../../fonts/]
\newfontfamily\popxbold{PlusJakartaSans-ExtraBold.ttf}[Path=../../../fonts/]
\newfontfamily\popxboldls{PlusJakartaSans-ExtraBold.ttf}[Path=../../../fonts/,LetterSpace=3.0]

\setlist[enumerate]{leftmargin=1.7em,itemsep=5pt,topsep=4pt}
\setlist[itemize]{leftmargin=1.7em,itemsep=4pt,topsep=4pt,label={\color{poporange}\textbullet}}
\setlength{\parindent}{0pt}
\setlength{\parskip}{5pt}
\renewcommand{\arraystretch}{1.25}

% pgfplots : style Pop par défaut
\pgfplotsset{
  every axis/.append style={axis line style={popink,line width=1pt},tick style={popink},
    grid style={popline,line width=0.5pt},label style={font=\small},tick label style={font=\footnotesize}},
  popcurve/.style={poporange,line width=1.8pt,no markers,smooth},
}
% Même style disponible dans un \draw TikZ simple (hors pgfplots)
\tikzset{popcurve/.style={poporange,line width=1.8pt}}
\tikzset{popgrid/.style={popline,very thin}}
\colorlet{popcurve}{poporange} % le nom du style est parfois utilisé comme couleur

% ============================================================
% Pill / squiggle
% ============================================================
\newcommand{\poppill}[3]{%
  \tikz[baseline=-0.55ex]\node[draw=popink,line width=1.2pt,rounded corners=2.6pt,%
    fill=#2,inner xsep=6.5pt,inner ysep=3pt]{%
    \popxboldls\fontsize{7}{7}\selectfont\color{#3}\MakeUppercase{#1}};%
}
\newcommand{\popsquiggle}[1]{%
  \tikz[baseline=-0.4ex]{
    \draw[#1,line width=2.6pt,line cap=round]
      plot [smooth] coordinates {(0,0.09) (0.34,0.01) (0.68,0.21) (1.02,0.01) (1.36,0.10)};
  }%
}
% Mot-clé surligné (marqueur orange sous le texte)
\newcommand{\cle}[1]{{\popxbold\color{popdark}#1}}

% ============================================================
% En-tête / pied
% ============================================================
\pagestyle{fancy}
\fancyhf{}
\renewcommand{\headrulewidth}{0pt}
\renewcommand{\footrulewidth}{0pt}
\lhead{\raisebox{-0.15ex}{\poplogo\fontsize{13}{13}\selectfont\color{popink}OnBuch\color{poporange}+}}
\rhead{\poppill{\DOCMATIERE\ \textbullet\ \DOCNIVEAU\ \textbullet\ Leçon \DOCLECON}{white}{popink}}
\lfoot{\popxbold\fontsize{7.6}{7.6}\selectfont\color{popmuted}\MakeUppercase{Cours signé OnBuch+}\ \textbullet\ {\popsemi\DOCTITRE}}
\rfoot{\tikz[baseline=-0.6ex]\node[rounded corners=8.5pt,fill=popink,minimum height=17pt,minimum width=17pt,inner xsep=5pt,inner sep=0]{\popxbold\fontsize{8}{8}\selectfont\color{popcream}\thepage\,/\,\pageref*{LastPage}};}

% ============================================================
% Titres
% ============================================================
\newcommand{\chaptitre}[2]{%
  \begin{tcolorbox}[enhanced,colback=poporange,colframe=popink,boxrule=2.6pt,arc=11pt,
      drop shadow=popink,left=15pt,right=15pt,top=12pt,bottom=11pt,before skip=3pt,after skip=18pt]
    \poppill{#2}{popink}{popcream}\par\vspace{9pt}
    {\raggedright\hyphenpenalty=10000\exhyphenpenalty=10000\popdisplay\fontsize{22}{24}\selectfont\color{popcream}\MakeUppercase{#1}\par}
  \end{tcolorbox}
}
% Section numérotée : pastille orange avec le numéro + titre Archivo
\newcounter{popsec}
\newcommand{\coursec}[1]{%
  \stepcounter{popsec}\setcounter{popsub}{0}%
  \par\vspace{16pt}\noindent
  \begin{minipage}{\linewidth}
  \tikz[baseline=-0.7ex]\node[draw=popink,line width=1.6pt,fill=poporange,rounded corners=4pt,minimum size=22pt,inner sep=0]{\popdisplay\fontsize{12}{12}\selectfont\color{popcream}\thepopsec};%
  \hspace{8pt}{\popdisplay\fontsize{15}{17}\selectfont\color{popink}\MakeUppercase{#1}}\par
  \vspace{4pt}\noindent{\color{poporange}\rule{36pt}{3.6pt}}
  \end{minipage}\par\nopagebreak\vspace{8pt}
}
\newcounter{popsub}
\newcommand{\courssub}[1]{%
  \stepcounter{popsub}%
  \par\vspace{9pt}\noindent{\popxbold\fontsize{11.5}{13}\selectfont\color{popink}{\color{poporange}\thepopsec.\thepopsub}\hspace{6pt}#1}\par\nopagebreak\vspace{4pt}
}

% ============================================================
% Boîtes pédagogiques — même recette : fond blanc, bordure épaisse colorée,
% ombre dure encre, badge plein. Titre optionnel [#1].
% ============================================================
\tcbset{popbase/.style={breakable,enhanced,colback=white,boxrule=2.3pt,arc=9pt,
  shadow={3pt}{-3pt}{0pt}{popink},left=14pt,right=14pt,top=11pt,bottom=11pt,before skip=11pt,after skip=15pt,
  fontupper=\popsemi\fontsize{10.6}{14.5}\selectfont\color{popink},
  fontlower=\popsemi\fontsize{10.6}{14.5}\selectfont\color{popink}}}
% Coche dessinée (le glyphe ✓ n'existe pas dans les polices du document)
\newcommand{\popcheck}[1]{\tikz[baseline=-0.5ex]\draw[#1,line width=1.6pt,line cap=round,line join=round](-0.13,0.0)--(-0.04,-0.09)--(0.13,0.1);}
\providecommand{\coche}{\popcheck{popgreen}}
\providecommand{\resultat}[1]{\par\smallskip\noindent\fcolorbox{popgreen}{popgreenL}{\parbox{\dimexpr\linewidth-2\fboxsep-2\fboxrule\relax}{\textbf{Résultat.} #1}}\par\smallskip}
\newcommand{\popboxhead}[3]{\poppill{#1}{#2}{white}\ifx\relax#3\relax\else\hspace{6pt}{\popxbold\fontsize{10.6}{12}\selectfont\color{popink}#3}\fi\par\vspace{7pt}}

\newtcolorbox{aretenir}[1][]{popbase,colframe=popgreen,before upper={\popboxhead{À retenir}{popgreen}{#1}}}
% Texte en anglais (document de lecture, dialogue, extrait) : cadre
% d'étude. Un titre optionnel (source, auteur) s'affiche dans l'en-tête.
\newtcolorbox{textebox}[1][]{popbase,colframe=popblue,before upper={\popboxhead{Text}{popblue}{#1}\begin{otherlanguage}{english}},after upper={\end{otherlanguage}}}
% Marques ✓ / ✗ (forme correcte / fautive) souvent utilisées par les modèles
\providecommand{\cross}{\textcolor{poppink}{\ensuremath{\times}}}
\providecommand{\xmark}{\cross}\providecommand{\crossmark}{\cross}\providecommand{\cmark}{\ensuremath{\checkmark}}
% Ligne de réponse / justification à compléter (inventée par les modèles)
\providecommand{\justification}[1]{\par\noindent\textit{Justification~:} \dotfill\par}
% \textipa (package tipa non chargé) : la police ne couvre pas l'API -> simple passage
\providecommand{\textipa}[1]{#1}
% Soulignement (ulem non chargé) : \uline, \ul
\providecommand{\uline}[1]{\underline{#1}}\providecommand{\ul}[1]{\underline{#1}}
% \glossaire{\item ...} inventée par les modèles : liste de glossaire
\providecommand{\glossaire}[1]{\begin{itemize}#1\end{itemize}}
% \alert (beamer) inventée par les modèles : texte mis en évidence
\providecommand{\alert}[1]{\textbf{\textcolor{poppink}{#1}}}
% variantes ASCII d'ordinaux écrites en macro
\providecommand{\iere}{\textsuperscript{re}}\providecommand{\ieme}{\textsuperscript{e}}\providecommand{\ere}{\textsuperscript{re}}\providecommand{\eme}{\textsuperscript{e}}\providecommand{\er}{\textsuperscript{er}}\providecommand{\ie}{\textsuperscript{e}}
% Ordinaux français écrits en macro par les modèles : 1\ière, 3\ième
\newcommand{\ière}{\textsuperscript{re}}\newcommand{\ième}{\textsuperscript{e}}
% \exemple (inventée par les modèles) : étiquette « Exemple : »
\providecommand{\exemple}{\textbf{Exemple~:}\ }
% \square utilisable aussi hors mode math (cases à cocher)
\AtBeginDocument{\let\squareMath\square\renewcommand{\square}{\ensuremath{\squareMath}}}
% Mot ou phrase en anglais à l'intérieur d'un paragraphe français
\newcommand{\eng}[1]{\ifmmode\text{\textit{#1}}\else\begingroup\selectlanguage{english}\itshape #1\endgroup\fi}
\let\esp\eng % alias toléré
% flèches d'intonation inventées par les modèles
\providecommand{\textsearrow}{}\renewcommand{\textsearrow}{\ensuremath{\searrow}}\providecommand{\textnearrow}{}\renewcommand{\textnearrow}{\ensuremath{\nearrow}}
% \fr{..} : mot français (inventé par les modèles)
\providecommand{\fr}[1]{\textit{#1}}
\newtcolorbox{exemplebox}[1][]{popbase,colframe=poporange,before upper={\popboxhead{Exemple}{poporange}{#1}}}
\newtcolorbox{definition}[1][]{popbase,colframe=popblue,before upper={\popboxhead{Définition}{popblue}{#1}}}
\newtcolorbox{propriete}[1][]{popbase,colframe=poppurple,before upper={\popboxhead{Propriété}{poppurple}{#1}}}
\newtcolorbox{methode}[1][]{popbase,colframe=popgold,before upper={\popboxhead{Méthode}{popgold}{#1}}}
\newtcolorbox{attention}[1][]{popbase,colframe=poppink,before upper={\popboxhead{Attention}{poppink}{#1}}}
\newtcolorbox{experience}[1][]{popbase,colframe=popblue,colback=popblueL,before upper={\popboxhead{Expérience}{popblue}{#1}}}
% « Le savais-tu ? » : carte violette pleine (contexte, culture, Cameroun)
\newtcolorbox{savaistu}[1][]{popbase,colback=poppurple,colframe=popink,
  fontupper=\popsemi\fontsize{10.4}{14.2}\selectfont\color{white},
  before upper={\poppill{Le savais-tu\,?}{white}{poppurple}\ifx\relax#1\relax\else\hspace{6pt}{\popxbold\color{white}#1}\fi\par\vspace{7pt}}}
% Exercice résolu : énoncé en haut, \tcblower puis la solution sur fond crème
\newcounter{popexo}\newcounter{popexr}
\newtcolorbox{exoresolu}[1][]{popbase,colframe=poporange,colbacklower=popcream,
  segmentation style={popink,line width=1.2pt,dashed},
  before upper={\stepcounter{popexr}\popboxhead{Exercice résolu \thepopexr}{poporange}{#1}},
  before lower={\poppill{Solution}{popink}{popcream}\par\vspace{6pt}}}
% Exercice (non résolu en place) — la correction est dans la partie Corrigés
\newtcolorbox{exercice}[1][]{popbase,colframe=popink,
  before upper={\stepcounter{popexo}\popboxhead{Exercice \thepopexo}{popink}{#1}}}
% Corrigé
\newtcolorbox{corrige}[1][]{popbase,colframe=popgreen,colback=popgreenL,
  before upper={\popboxhead{Corrigé}{popgreen}{#1}}}
% Objectifs / prérequis / bilan
\newtcolorbox{objectifs}{popbase,colframe=popink,colback=poporangeL,
  before upper={\popboxhead{Objectifs de la leçon}{popink}{}}}
\newtcolorbox{prerequis}{popbase,colframe=popink,colback=popblueL,
  before upper={\popboxhead{Prérequis}{popblue}{}}}
\newtcolorbox{bilan}{popbase,colframe=popink,colback=white,boxrule=2.8pt,
  before upper={\popboxhead{Fiche bilan}{poporange}{L'essentiel à connaître pour l'examen}}}
% Figure encadrée avec légende
\newtcolorbox{popfigure}[1][]{enhanced,breakable=false,colback=white,colframe=popline,boxrule=1.4pt,arc=8pt,
  left=10pt,right=10pt,top=10pt,bottom=8pt,before skip=10pt,after skip=12pt,halign=center,
  lower separated=false,halign lower=center,
  fontlower=\popsemi\fontsize{9}{11.5}\selectfont\color{popmuted},
  #1}
\newcommand{\legende}[1]{\par\vspace{6pt}{\popsemi\fontsize{9}{11.5}\selectfont\color{popmuted}#1}}

% ============================================================
% Signature OnBuch+
% ============================================================
\newcommand{\poplogoinline}[1]{{\poplogo\fontsize{#1}{#1}\selectfont\color{popink}OnBuch\color{poporange}+}}
\newcommand{\signatureonbuch}{%
  \begin{tcolorbox}[enhanced,colback=popink,colframe=popink,boxrule=2.2pt,arc=10pt,
      drop shadow=poporange,left=14pt,right=14pt,top=10pt,bottom=10pt,before skip=0pt,after skip=0pt]
    \tikz[baseline=-0.7ex]\node[circle,fill=popgreen,draw=popcream,line width=1.3pt,minimum size=22pt,inner sep=0]{\popcheck{white}};%
    \hspace{9pt}\begin{minipage}[c]{0.62\linewidth}
      {\popxbold\fontsize{12}{13}\selectfont\color{popcream}Signé\ \textbullet\ {\poplogo OnBuch\color{poporange}+}}\par\vspace{2pt}
      {\raggedright\popsemi\fontsize{8}{10.5}\selectfont\color{popline}\MakeUppercase{Équipe pédagogique OnBuch+}\\\MakeUppercase{Conforme au programme officiel MINESEC}\par}
    \end{minipage}\hfill
    {\poplogo\fontsize{22}{22}\selectfont\color{popcream}OnBuch\color{poporange}+}
  \end{tcolorbox}%
}

% ============================================================
% Page de garde
% #1 titre · #2 matière · #3 niveau · #4 module · #5 n° leçon · #6 durée ·
% #7 description / objectifs (texte ou itemize)
% ============================================================
\newcommand{\pagegarde}[7]{%
  \thispagestyle{empty}
  \newgeometry{a4paper,margin=1.7cm,top=1.6cm,bottom=1.4cm}%
  \begin{tikzpicture}[remember picture,overlay]
    \fill[poporange!18] ([xshift=-3.2cm,yshift=-3.6cm]current page.north east) circle (4.6cm);
    \fill[poppurple!14] ([xshift=2.4cm,yshift=7.6cm]current page.south west) circle (3.8cm);
  \end{tikzpicture}%
  {\poplogo\fontsize{30}{30}\selectfont\color{popink}OnBuch\color{poporange}+}\hfill
  \raisebox{0.6ex}{\poppill{Cours signé \textbullet\ Premium}{popink}{popcream}}\par
  \vspace{1pt}\popsquiggle{poppurple}\par
  \vspace{8pt}
  \begin{tcolorbox}[enhanced,colback=poporange,colframe=popink,boxrule=2.8pt,arc=13pt,
      drop shadow=popink,left=18pt,right=18pt,top=12pt,bottom=13pt,after skip=10pt]
    \poppill{#2 \textbullet\ #3}{popink}{popcream}\hspace{5pt}\poppill{#4}{white}{popink}\par\vspace{8pt}
    {\popdisplay\fontsize{14}{14}\selectfont\color{popink}LEÇON #5}\par\vspace{4pt}
    {\raggedright\hyphenpenalty=10000\exhyphenpenalty=10000\popdisplay\fontsize{25}{27}\selectfont\color{popcream}\MakeUppercase{#1}\par}
    \vspace{8pt}
    {\popxbold\fontsize{9.5}{9.5}\selectfont\color{popink}\MakeUppercase{Durée conseillée : #6}}
  \end{tcolorbox}
  \begin{tcolorbox}[enhanced,colback=white,colframe=popink,boxrule=2.2pt,arc=10pt,
      drop shadow=popink,left=16pt,right=16pt,top=10pt,bottom=10pt,after skip=8pt]
    \poppill{Dans cette leçon}{poppurple}{white}\par\vspace{5pt}
    \popsemi\fontsize{9.3}{12.6}\selectfont\color{popink}#7
  \end{tcolorbox}
  \noindent\begin{minipage}[t]{0.487\linewidth}
    \begin{tcolorbox}[enhanced,colback=popgreen,colframe=popink,boxrule=2.2pt,arc=10pt,
        drop shadow=popink,left=11pt,right=11pt,top=10pt,bottom=10pt,height=3.1cm,valign=center]
      \tcbox[colback=white,colframe=popink,boxrule=1.3pt,arc=5pt,boxsep=0pt,nobeforeafter,
        left=3pt,right=3pt,top=3pt,bottom=3pt,valign=center]{\qrcode[height=1.9cm]{https://play.google.com/store/apps/details?id=com.luvvix.onbuch&hl=fr}}%
      \hspace{8pt}\begin{minipage}[c]{0.5\linewidth}
        {\popdisplay\fontsize{11}{12.5}\selectfont\color{white}\MakeUppercase{Télécharge OnBuch}}\par\vspace{4pt}
        {\raggedright\popsemi\fontsize{8.4}{11}\selectfont\color{white}Gratuit \textbullet\ Google Play\\Tuteur IA, annales, concours, résultats\par}
      \end{minipage}
    \end{tcolorbox}
  \end{minipage}\hfill
  \begin{minipage}[t]{0.487\linewidth}
    \begin{tcolorbox}[enhanced,colback=poppurple,colframe=popink,boxrule=2.2pt,arc=10pt,
        drop shadow=popink,left=11pt,right=11pt,top=10pt,bottom=10pt,height=3.1cm,valign=center]
      \tikz[baseline=-0.7ex]\node[circle,fill=white,draw=popink,line width=1.3pt,minimum size=1.6cm,inner sep=0]{\popdisplay\fontsize{24}{24}\selectfont\color{poppurple}+};%
      \hspace{8pt}\begin{minipage}[c]{0.6\linewidth}
        {\popdisplay\fontsize{11}{12.5}\selectfont\color{white}\MakeUppercase{Passe à OnBuch+}}\par\vspace{4pt}
        {\raggedright\popsemi\fontsize{8.4}{11}\selectfont\color{white}Tous les cours signés, Tuteur IA illimité, fiches PDF — onglet OnBuch+ de l'app\par}
      \end{minipage}
    \end{tcolorbox}
  \end{minipage}
  \vspace{8pt}
  \popcontenu
  \vfill
  \signatureonbuch
  \restoregeometry
  \newpage
}

% Rangée « Ce cours contient » (page de garde)
\newcommand{\poptile}[3]{% #1 couleur #2 gros texte #3 légende
  \begin{tcolorbox}[enhanced,colback=#1,colframe=popink,boxrule=2pt,arc=9pt,shadow={2.5pt}{-2.5pt}{0pt}{popink},
      left=6pt,right=6pt,top=9pt,bottom=9pt,halign=center,valign=center,height=1.75cm,width=\linewidth,nobeforeafter]
    {\popdisplay\fontsize{12.5}{13}\selectfont\color{white}\MakeUppercase{#2}}\par\vspace{3pt}
    {\centering\popsemi\fontsize{7.8}{9.5}\selectfont\color{white}#3\par}
  \end{tcolorbox}}
\newcommand{\popcontenu}{%
  \par\noindent{\popxboldls\fontsize{7.5}{7.5}\selectfont\color{popmuted}CE COURS SIGNÉ CONTIENT}\par\vspace{7pt}
  \noindent\begin{minipage}[t]{0.235\linewidth}\poptile{popblue}{Cours}{complet, illustré, pas à pas}\end{minipage}\hfill
  \begin{minipage}[t]{0.235\linewidth}\poptile{poporange}{Méthodes}{et exercices résolus}\end{minipage}\hfill
  \begin{minipage}[t]{0.235\linewidth}\poptile{poppink}{Exercices}{type examen + corrigés}\end{minipage}\hfill
  \begin{minipage}[t]{0.235\linewidth}\poptile{popgreen}{Fiche bilan}{l'essentiel à retenir}\end{minipage}\par}

% Sommaire
\newtcolorbox{sommaire}{enhanced,colback=white,colframe=popink,boxrule=2.2pt,arc=10pt,
  drop shadow=popink,left=18pt,right=18pt,top=13pt,bottom=13pt,before skip=6pt,after skip=20pt,
  before upper={\poppill{Sommaire}{poppurple}{white}\par\vspace{9pt}\popsemi\fontsize{10.6}{14.5}\selectfont\color{popink}}}

% Signature de fin de cours — poussée tout en bas de la dernière page
\newcommand{\findecours}{%
  \par\vfill\nopagebreak
  \begin{center}\popsquiggle{poporange}\end{center}\vspace{4pt}
  \signatureonbuch
}
\AtBeginDocument{\def\llbracket{\mathopen{[\mkern-3.5mu[}}\def\rrbracket{\mathclose{]\mkern-3.5mu]}}} % intervalles d'entiers (glyphe ⟦ absent de la police)
% Glyphes absents de Fira Math : redéfinis à partir de glyphes disponibles
\AtBeginDocument{%
  \def\setminus{\mathbin{\backslash}}\let\smallsetminus\setminus
  \def\vdots{\mathord{\vbox{\baselineskip3.2pt\lineskiplimit0pt\kern4pt\hbox{.}\hbox{.}\hbox{.}}}}%
  \def\ddots{\mathinner{\mkern1mu\raise6.5pt\hbox{.}\mkern2mu\raise3.6pt\hbox{.}\mkern2mu\raise0.7pt\hbox{.}\mkern1mu}}%
  \def\triangle{\mathord{\scalebox{0.9}{$\symup{\Delta}$}}}%
  \def\nabla{\mathord{\raisebox{\depth}{\scalebox{1}[-1]{$\symup{\Delta}$}}}}%
  \def\checkmark{\popcheck{popdark}}%
  \def\star{\mathbin{\ast}}\def\bigstar{\mathord{\ast}}%
  \def\diamond{\mathbin{\scalebox{0.6}{$\Diamond$}}}%
  \def\bigcup{\mathop{\mathchoice{\raisebox{-0.3em}{\scalebox{1.7}{$\cup$}}}{\scalebox{1.25}{$\cup$}}{\cup}{\cup}}\displaylimits}%
  \def\bigcap{\mathop{\mathchoice{\raisebox{-0.3em}{\scalebox{1.7}{$\cap$}}}{\scalebox{1.25}{$\cap$}}{\cap}{\cap}}\displaylimits}%
  \def\lesseqqgtr{\mathrel{\lessgtr}}\def\bigoplus{\mathop{\oplus}\displaylimits}%
}
\providecommand{\ssi}{si et seulement si} % abréviation fréquente
\AtBeginDocument{\providecommand{\wideparen}{\overparen}} % arc de cercle

\begin{document}

\pagegarde{Lesson 21 — Grammar: Adjectives as nouns; "So do I", "Neither did he"}{Anglais}{1ère A}{Chapter 2 : Economic Life and Occupations}{ENA21}{2 h}{Cette leçon de grammaire étudie deux structures essentielles de l'anglais : l'emploi des adjectifs comme noms (the rich, the unemployed) et les réponses courtes d'accord ou de désaccord (so do I, neither did he). L'élève apprend à les construire, à éviter les pièges du français et à les utiliser dans des échanges authentiques.
\vspace{4pt}
\begin{multicols}{2}\small
\begin{itemize}
\item À la fin de cette leçon, je sais identifier un adjectif utilisé comme nom dans une phrase.
\item Je sais former le nom collectif à partir d'un adjectif avec « the + adjective » (the poor, the elderly).
\item Je sais distinguer les adjectifs qui prennent « the » (sens collectif, pluriel) de ceux qui deviennent de vrais noms (a Frenchman, the French).
\item Je sais employer « so + auxiliaire + sujet » pour exprimer l'accord avec une phrase affirmative.
\item Je sais employer « neither/nor + auxiliaire + sujet » pour exprimer l'accord avec une phrase négative.
\item Je sais choisir le bon auxiliaire (do/does/did, be, have, can, will...) selon le verbe de la phrase de départ.
\end{itemize}\end{multicols}}

\begin{sommaire}
\begin{multicols}{2}
\begin{enumerate}
\item Observation : adjectifs employés comme noms
\item Règle et emplois de « the + adjective »
\item Comparaison français/anglais et pièges
\item Observation : « So do I » / « Neither did he »
\item Règle, tableaux de formes et choix de l'auxiliaire
\item Comparaison français/anglais et entraînement aux réponses courtes
\item Activité d'intégration
\item Méthodes et exercices résolus
\item Exercices
\item Corrigés des exercices
\item Fiche bilan
\end{enumerate}
\end{multicols}
\end{sommaire}

\begin{objectifs}
\begin{itemize}
\item À la fin de cette leçon, je sais identifier un adjectif utilisé comme nom dans une phrase.
\item Je sais former le nom collectif à partir d'un adjectif avec « the + adjective » (the poor, the elderly).
\item Je sais distinguer les adjectifs qui prennent « the » (sens collectif, pluriel) de ceux qui deviennent de vrais noms (a Frenchman, the French).
\item Je sais employer « so + auxiliaire + sujet » pour exprimer l'accord avec une phrase affirmative.
\item Je sais employer « neither/nor + auxiliaire + sujet » pour exprimer l'accord avec une phrase négative.
\item Je sais choisir le bon auxiliaire (do/does/did, be, have, can, will...) selon le verbe de la phrase de départ.
\item Je sais éviter les erreurs fréquentes des francophones (moi aussi, me too, so I do...).
\item Je sais utiliser ces structures dans des dialogues et de courts textes sur la vie économique et les métiers.
\end{itemize}
\end{objectifs}

\begin{prerequis}
\begin{itemize}
\item Les auxiliaires anglais : do/does/did, be, have, can, will, must (Seconde).
\item La formation des phrases affirmatives et négatives aux temps simples (Seconde).
\item L'article défini « the » et ses emplois (Seconde).
\item Les adjectifs qualificatifs et leur place en anglais (Seconde).
\item Le vocabulaire de base des métiers et de la vie économique (Chapter 2, leçons précédentes).
\end{itemize}
\end{prerequis}

\newpage

\coursec{Observation : adjectifs employés comme noms}

\courssub{Situation d'accroche et problématique}

Pendant la campagne électorale à Bamenda, un journaliste de la télévision interviewe des habitants dans la rue. Au micro, un vieil homme déclare : \eng{The young want jobs, the old want peace, and the rich don't always understand the poor.} « Les jeunes veulent du travail, les vieux veulent la paix, et les riches ne comprennent pas toujours les pauvres. » Dans la salle de classe, un élève lève la main, intrigué : \emph{Monsieur, comment un adjectif comme} \eng{young} \emph{peut-il devenir un nom ?} Excellente question ! En français, nous disons « les jeunes », « les riches », « les pauvres » : l'adjectif, précédé de l'article, désigne toute une catégorie de personnes. L'anglais fait exactement la même chose — mais avec ses propres règles, et quelques pièges.

Plus tard, au marché de Bamenda, deux amies discutent : \eng{I love fresh mangoes. — So do I!} « J'adore les mangues fraîches. — Moi aussi ! » puis \eng{He didn't buy fish yesterday. — Neither did she.} « Il n'a pas acheté de poisson hier. — Elle non plus. » Ces petites réponses toutes faites, si fréquentes dans les conversations anglaises, obéissent elles aussi à des règles précises.

\textbf{Problématique de la leçon :} comment l'anglais transforme-t-il un adjectif en nom de groupe, et comment construit-on les réponses courtes d'accord ou de désaccord \eng{so do I} et \eng{neither did he} ? Dans cette première partie, nous observons d'abord les adjectifs employés comme noms.

\courssub{Lecture : un texte sur la vie économique au Cameroun}

Lisons attentivement ce court texte sur la vie à Douala. Pendant la lecture, souligne mentalement tous les groupes de mots formés de \eng{the} suivi d'un adjectif.

\begin{textebox}[Economic life in Douala — reading text]
In Douala, the economic capital of Cameroon, life is not the same for everybody. The rich live in beautiful houses in Bonanjo, while the poor struggle every day in the crowded quarters of New Bell and Ndogpassi. The young dream of jobs abroad, in Europe or in America, but the old prefer to stay near their families and their villages.

Every morning, the unemployed wait near the port, hoping for work. Some carry heavy bags for a few francs; others sell phone credit or groundnuts in the streets. The disabled often sit at the doors of shops and offices, asking for help from people who pass by.

The government says it wants to help the poor and create jobs for the young. New training centres are opening for the unemployed, and special programmes exist for the disabled. But journalists ask difficult questions: are the rich paying enough taxes? Do the old receive their pensions on time? In Douala, as in every big city, the distance between the rich and the poor is still very wide.
\end{textebox}

\textbf{Glossaire :}

\begin{center}
\begin{tabularx}{\linewidth}{|l|l|X|}
\hline
\rowcolor{popblueL} \textbf{Mot anglais} & \textbf{Nature} & \textbf{Traduction} \\
\hline
crowded & adjectif & bondé, surpeuplé \\
\hline
abroad & adverbe & à l'étranger \\
\hline
to struggle & verbe & lutter, peiner \\
\hline
to hope for & verbe & espérer (quelque chose) \\
\hline
groundnuts & nom pluriel & arachides \\
\hline
pension & nom & retraite, pension \\
\hline
tax & nom & impôt, taxe \\
\hline
wide & adjectif & large, grand \\
\hline
\end{tabularx}
\end{center}

\textbf{Questions de compréhension (reading) :}

\begin{enumerate}
\item Where do the rich live in Douala?
\item What do the young dream of?
\item Where do the unemployed wait every morning, and why?
\item What two questions do journalists ask?
\item According to the text, what is still very wide in Douala?
\end{enumerate}

\courssub{Repérage guidé : quel mot précède l'adjectif ?}

Relevons maintenant, dans le texte, tous les groupes de mots qui désignent des \emph{catégories de personnes} :

\begin{center}
\begin{tabularx}{\linewidth}{|X|X|X|}
\hline
\rowcolor{popblueL} \textbf{Groupe relevé dans le texte} & \textbf{Traduction} & \textbf{Catégorie désignée} \\
\hline
the rich & les riches & les personnes riches \\
\hline
the poor & les pauvres & les personnes pauvres \\
\hline
the young & les jeunes & les jeunes gens \\
\hline
the old & les vieux & les personnes âgées \\
\hline
the unemployed & les chômeurs & les personnes sans emploi \\
\hline
the disabled & les handicapés & les personnes handicapées \\
\hline
\end{tabularx}
\end{center}

\begin{aretenir}[Premier constat]
Dans chaque groupe, le mot qui précède l'adjectif est \textbf{toujours} l'article défini \cle{the}. Jamais \eng{a}, jamais \eng{some}, jamais l'article zéro. La structure est donc : \textbf{THE + adjectif}.
\end{aretenir}

\textbf{Question déclencheuse :} ces groupes de mots désignent-ils \emph{une} personne ou une \emph{catégorie} de personnes ? Regarde à nouveau le texte : \eng{the unemployed wait near the port} — est-ce un seul chômeur qui attend ? Non : c'est \emph{tous} les chômeurs, la catégorie entière. De même, \eng{the young dream of jobs abroad} parle de la jeunesse en général, pas d'un jeune en particulier.

\begin{definition}[Adjectif employé comme nom]
En anglais, \textbf{the + adjectif} désigne une \cle{catégorie de personnes} tout entière, considérée comme un groupe : \eng{the rich} = « les riches » (toutes les personnes riches), \eng{the blind} = « les aveugles », \eng{the homeless} = « les sans-abri ». L'adjectif joue alors le rôle d'un \emph{nom pluriel}.
\end{definition}

\courssub{Observation de l'accord du verbe : toujours le pluriel}

Revenons à une phrase du texte : \eng{The poor struggle every day in the crowded quarters.} Le verbe \eng{struggle} est à la forme du \textbf{pluriel} (sans -s à la troisième personne). De même, dans notre phrase titre : \eng{The poor ARE suffering.} « Les pauvres souffrent. » — le verbe \eng{to be} est à la forme \eng{are}, jamais \eng{is}.

\begin{propriete}[Accord du verbe avec « the + adjectif »]
Le groupe \textbf{the + adjectif} est toujours suivi d'un verbe au \cle{pluriel}, car il désigne un ensemble de personnes :
\begin{itemize}
\item \eng{The rich get richer.} « Les riches deviennent plus riches. » (et non \emph{gets})
\item \eng{The unemployed need help.} « Les chômeurs ont besoin d'aide. » (et non \emph{needs})
\item \eng{The old are often lonely.} « Les personnes âgées sont souvent seules. » (et non \emph{is})
\end{itemize}
\end{propriete}

\begin{exemplebox}[Exemples traduits]
\begin{itemize}
\item \eng{The rich get richer.} « Les riches deviennent plus riches. »
\item \eng{The unemployed need help.} « Les chômeurs ont besoin d'aide. »
\item \eng{The young have energy, but the old have experience.} « Les jeunes ont de l'énergie, mais les vieux ont de l'expérience. »
\item \eng{The disabled deserve respect.} « Les handicapés méritent le respect. »
\item \eng{In Buea, the homeless sleep near the market.} « À Buea, les sans-abri dorment près du marché. »
\item \eng{The deaf use sign language to communicate.} « Les sourds utilisent la langue des signes pour communiquer. »
\end{itemize}
\end{exemplebox}

\courssub{Comparaison : « a poor man » ou « the poor » ?}

C'est le point capital de cette observation. Comparons deux phrases :

\begin{itemize}
\item \eng{A poor man was sitting near the church.} « Un pauvre homme était assis près de l'église. » → \textbf{une seule personne}, verbe au singulier (\eng{was}).
\item \eng{The poor were sitting near the church.} « Les pauvres étaient assis près de l'église. » → \textbf{toute la catégorie}, verbe au pluriel (\eng{were}).
\end{itemize}

\begin{center}
\begin{tabularx}{\linewidth}{|X|X|X|}
\hline
\rowcolor{popblueL} \textbf{Structure} & \textbf{Sens} & \textbf{Accord du verbe} \\
\hline
a poor man / a young girl & une seule personne & singulier : \eng{A poor man \textbf{is} waiting.} \\
\hline
the poor / the young & toute la catégorie & pluriel : \eng{The poor \textbf{are} waiting.} \\
\hline
\end{tabularx}
\end{center}

\begin{attention}[Piège : une personne, pas « the + adjectif »]
Pour désigner \textbf{une seule personne}, l'anglais ne peut pas utiliser \eng{the + adjectif} seul : il faut ajouter un nom (\eng{man}, \eng{woman}, \eng{person}, \eng{people}). On dit \eng{a rich man} « un homme riche », \eng{a disabled person} « une personne handicapée » — jamais \emph{a rich} ou \emph{a disabled} tout seul. La structure \eng{the + adjectif} sert uniquement pour le \textbf{groupe entier}.
\end{attention}

\courssub{Carte mentale : la famille de « THE + ADJECTIVE »}

\begin{popfigure}
\begin{tikzpicture}[mindmap, grow cyclic,
  every node/.style={concept, align=center},
  root concept/.append style={concept color=popblue, font=\bfseries, minimum size=2.6cm},
  level 1 concept/.append style={sibling angle=60, level distance=3.6cm, minimum size=1.7cm, font=\small\bfseries},
  level 2 concept/.append style={level distance=2.6cm, minimum size=2.1cm, font=\scriptsize, concept color=popgreenL}]
\node[root concept]{THE + ADJECTIVE}
  child[concept color=poporangeL]{ node {the rich} child { node[align=center]{a group of people\\plural verb} } }
  child[concept color=poporangeL]{ node {the poor} child { node[align=center]{a group of people\\plural verb} } }
  child[concept color=poporangeL]{ node {the young} child { node[align=center]{a group of people\\plural verb} } }
  child[concept color=poporangeL]{ node {the old} child { node[align=center]{a group of people\\plural verb} } }
  child[concept color=poporangeL]{ node {the unemployed} child { node[align=center]{a group of people\\plural verb} } }
  child[concept color=poporangeL]{ node {the disabled} child { node[align=center]{a group of people\\plural verb} } };
\end{tikzpicture}
\legende{Carte mentale : chaque groupe « the + adjectif » désigne une catégorie de personnes et exige un verbe au pluriel.}
\end{popfigure}

\courssub{Comparaison avec le français}

Bonne nouvelle : le français possède exactement le même procédé ! « Les riches », « les pauvres », « les jeunes » fonctionnent comme \eng{the rich}, \eng{the poor}, \eng{the young}. Mais il y a deux différences importantes à retenir.

\begin{center}
\begin{tabularx}{\linewidth}{|X|X|X|}
\hline
\rowcolor{popblueL} \textbf{Français} & \textbf{English} & \textbf{Remarque} \\
\hline
les riches (avec un -s de pluriel) & the rich (jamais de -s) & l'adjectif anglais est invariable \\
\hline
les jeunes & the young & pas de -s en anglais \\
\hline
les chômeurs (un vrai nom) & the unemployed (adjectif + the) & le français nominalise, l'anglais garde l'adjectif \\
\hline
les handicapés & the disabled & même construction, même sens de groupe \\
\hline
un aveugle (une personne) & a blind man / a blind person & pour une personne, l'anglais ajoute un nom \\
\hline
\end{tabularx}
\end{center}

\begin{attention}[Erreur fréquente des francophones : le -s interdit]
En français, on écrit « les pauvre\emph{s} », « les riche\emph{s} » avec un -s de pluriel. En anglais, c'est impossible : l'adjectif anglais est \textbf{invariable}. Ne dis jamais \emph{the poors} ou \emph{the riches} ! La forme correcte est \eng{the poor}, \eng{the rich}, sans -s, même si le sens est pluriel et même si le verbe est au pluriel : \eng{The poor \textbf{are} suffering.} « Les pauvres souffrent. »
\end{attention}

\courssub{Première généralisation orale}

\begin{experience}[À vous de formuler la règle !]
Avant de lire la règle officielle à la section suivante, essayez, en groupes de trois, de compléter oralement cette généralisation à partir de vos observations :
\begin{enumerate}
\item Pour parler d'une catégorie de personnes en anglais, on utilise l'article \ldots\ suivi d'un \ldots
\item Ce groupe de mots désigne \ldots\ personne / une \ldots\ de personnes. (barre la mauvaise réponse)
\item Le verbe qui suit se met au \ldots\ (singulier / pluriel).
\item L'adjectif anglais prend-il un -s ? \ldots
\end{enumerate}
Chaque groupe présente sa règle à la classe, puis le professeur la confronte aux exemples du texte : \eng{The rich live in Bonanjo}, \eng{The unemployed wait near the port}, \eng{The old prefer to stay}. La règle écrite sera fixée à la section 2.
\end{experience}

\courssub{Exercice résolu d'observation}

\begin{exoresolu}[Repérage et accord]
\textbf{Énoncé.} Dans les phrases suivantes, (a) souligne le groupe « the + adjectif », (b) traduis la phrase, (c) dis si le verbe est au singulier ou au pluriel et justifie.

1. \eng{The young need jobs and hope.}
2. \eng{A young man is waiting for the bus.}
3. \eng{The disabled receive help from the government.}
\tcblower
\textbf{Solution détaillée.}

1. Groupe : \eng{the young}. Traduction : « Les jeunes ont besoin de travail et d'espoir. » Le verbe \eng{need} est au \textbf{pluriel} (pas de -s), car \eng{the young} désigne toute la catégorie des jeunes.

2. Ici, il n'y a \textbf{pas} de groupe « the + adjectif » : \eng{a young man} = « un jeune homme », une seule personne. Le verbe \eng{is} est donc au \textbf{singulier}. C'est la différence capitale entre \eng{a young man} (une personne) et \eng{the young} (tous les jeunes).

3. Groupe : \eng{the disabled}. Traduction : « Les handicapés reçoivent de l'aide du gouvernement. » Le verbe \eng{receive} est au \textbf{pluriel} (pas de -s), car il s'agit d'une catégorie entière de personnes. Et remarque : \eng{disabled} ne prend aucun -s, même si le sens est pluriel.
\end{exoresolu}

\begin{exercice}[À ton tour]
Traduis en anglais les phrases suivantes, en faisant attention à l'article, à l'absence de -s et à l'accord du verbe :
\begin{enumerate}
\item Les riches ne comprennent pas toujours les pauvres.
\item Les chômeurs attendent près du port de Douala.
\item Un vieil homme est assis devant sa maison. (une seule personne !)
\item Les jeunes rêvent de voyager à l'étranger.
\end{enumerate}
\end{exercice}

\begin{corrige}[Exercice « À ton tour »]
\begin{enumerate}
\item \eng{The rich do not always understand the poor.} (deux catégories, deux verbes au pluriel, aucun -s sur les adjectifs)
\item \eng{The unemployed wait near the port of Douala.} (\eng{wait} sans -s : pluriel)
\item \eng{An old man is sitting in front of his house.} (une personne : \eng{an old man} + verbe au singulier \eng{is})
\item \eng{The young dream of travelling abroad.} (attention : \eng{abroad} « à l'étranger », à ne pas confondre avec \eng{aboard} « à bord »)
\end{enumerate}
\end{corrige}

\begin{savaistu}[Et les choses abstraites ?]
La structure \eng{the + adjectif} peut aussi désigner des notions abstraites, mais alors le verbe est au \textbf{singulier} : \eng{the beautiful} « le beau », \eng{the unknown} « l'inconnu », \eng{the impossible} « l'impossible ». Exemple : \eng{The unexpected always happens.} « L'inattendu arrive toujours. » Retiens le contraste : \eng{the + adjectif} pour des \emph{personnes} → pluriel ; pour une \emph{notion abstraite} → singulier. Au lycée, l'emploi le plus fréquent reste celui des catégories de personnes, comme dans notre texte sur Douala.
\end{savaistu}

\begin{aretenir}[Bilan de l'observation]
\begin{itemize}
\item Structure : \textbf{the + adjectif} = une catégorie entière de personnes.
\item Le verbe qui suit est toujours au \textbf{pluriel} : \eng{The poor are suffering.}
\item L'adjectif ne prend \textbf{jamais de -s} : \eng{the rich}, jamais \emph{the riches}.
\item Pour une seule personne : \eng{a + adjectif + nom} : \eng{a poor man}, \eng{a disabled person}.
\end{itemize}
\end{aretenir}

\coursec{Règle et emplois de « the + adjective »}

Dans la section précédente, nous avons \emph{observé} des phrases comme \eng{The rich get richer and the poor get poorer} ou \eng{The government should help the unemployed}. Nous avons remarqué qu'un adjectif anglais, précédé de \eng{the}, peut désigner tout un groupe de personnes, sans qu'aucun nom ne le suive. Il est temps maintenant de transformer cette observation en \cle{règle} précise : quelle est la forme exacte ? Quel sens prend-elle ? Quel nombre (singulier ou pluriel) faut-il utiliser pour le verbe ? Quelles sont les exceptions ? C'est tout l'objet de cette section.

\courssub{La règle générale : forme, sens, nombre du verbe}

\begin{propriete}[« THE + ADJECTIVE » = un groupe de personnes]
La structure \cle{the + adjectif} (sans nom après l'adjectif) désigne \textbf{toute une catégorie de personnes} partageant la même caractéristique. Elle a un \textbf{sens collectif} et le verbe qui suit se met toujours au \textbf{PLURIEL}.
\begin{center}
\eng{The elderly ARE respected in our villages.}\\
« Les personnes âgées SONT respectées dans nos villages. »
\end{center}
Autrement dit : \eng{the poor} = \eng{poor people} (les pauvres, les personnes pauvres) ; \eng{the young} = \eng{young people} (les jeunes).
\end{propriete}

Trois points sont à retenir absolument :

\begin{itemize}
\item \textbf{Forme} : \eng{the} + adjectif, \emph{et rien d'autre}. On n'ajoute jamais de \eng{-s} à l'adjectif : on dit \eng{the rich}, jamais \eng{the riches}.
\item \textbf{Sens} : la structure ne désigne pas une personne précise, mais \emph{tous} les membres de la catégorie, partout dans le monde ou dans le pays dont on parle.
\item \textbf{Accord} : le verbe est au pluriel (\eng{are}, \eng{have}, \eng{need}, \eng{live}...), car on parle de plusieurs personnes.
\end{itemize}

\begin{exemplebox}[Premiers exemples traduits]
\begin{itemize}
\item \eng{The blind need special schools.} = « Les aveugles ont besoin d'écoles spéciales. »
\item \eng{The young often prefer living in Douala or Yaoundé.} = « Les jeunes préfèrent souvent vivre à Douala ou à Yaoundé. »
\item \eng{The unemployed are asking the government for more jobs.} = « Les chômeurs demandent plus d'emplois au gouvernement. »
\item \eng{The sick must be taken to hospital quickly.} = « Les malades doivent être transportés rapidement à l'hôpital. »
\end{itemize}
\end{exemplebox}

\begin{attention}[Pas de « s », pas de nom !]
Deux erreurs sont fréquentes chez les francophones :
\begin{enumerate}
\item Ajouter un \eng{-s} à l'adjectif, sur le modèle français « les pauvre\emph{s} » : on ne dit jamais \eng{the poors} ni \eng{the youngs}. En anglais, l'adjectif est \textbf{invariable}.
\item Ajouter un nom après l'adjectif quand on veut parler du groupe entier : \eng{the poor people} est grammaticalement possible mais lourd ; la forme idiomatique et attendue à l'examen est \eng{the poor}.
\end{enumerate}
\end{attention}

\courssub{Emploi 1 : les catégories sociales}

C'est l'emploi le plus fréquent, surtout dans les textes sur la vie économique (notre chapitre \emph{Economic Life and Occupations}). La structure sert à désigner des groupes définis par leur situation sociale ou économique.

\begin{center}
\begin{tabularx}{\linewidth}{|X|X|X|}
\hline
\rowcolor{popblueL} English & Français & Remarque \\
\hline
the rich & les riches & verbe au pluriel \\
\hline
the poor & les pauvres & jamais \eng{the poors} \\
\hline
the homeless & les sans-abri & \eng{home} + suffixe \eng{-less} (sans) \\
\hline
the unemployed & les chômeurs & participe passé employé comme adjectif \\
\hline
the employed & les personnes qui ont un emploi & contraire du précédent \\
\hline
the underprivileged & les défavorisés & très fréquent dans la presse \\
\hline
the privileged & les privilégiés & \\
\hline
the needy & les nécessiteux & de \eng{need} (besoin) \\
\hline
the powerful & les puissants & \\
\hline
the weak & les faibles & \\
\hline
\end{tabularx}
\end{center}

\begin{exemplebox}[Dans des phrases]
\begin{itemize}
\item \eng{In every country, the rich pay more taxes than the poor.} = « Dans tous les pays, les riches paient plus d'impôts que les pauvres. »
\item \eng{During the rainy season, the homeless suffer a lot in big cities.} = « Pendant la saison des pluies, les sans-abri souffrent beaucoup dans les grandes villes. »
\item \eng{The unemployed of Bamenda organised a peaceful march.} = « Les chômeurs de Bamenda ont organisé une marche pacifique. »
\end{itemize}
\end{exemplebox}

\courssub{Emploi 2 : les catégories d'âge}

\begin{center}
\begin{tabularx}{\linewidth}{|X|X|X|}
\hline
\rowcolor{popblueL} English & Français & Remarque \\
\hline
the young & les jeunes & \\
\hline
the old & les vieux, les personnes âgées & parfois jugé un peu direct \\
\hline
the elderly & les personnes âgées & plus poli que \eng{the old} \\
\hline
the middle-aged & les personnes d'âge mûr & \eng{middle-aged} se prononce « mi-deul-èidjd » \\
\hline
\end{tabularx}
\end{center}

\eng{The young have energy; the elderly have experience.} = « Les jeunes ont de l'énergie ; les personnes âgées ont de l'expérience. »

\begin{savaistu}[Politesse linguistique]
Dans les pays anglophones, on préfère de plus en plus \eng{the elderly} ou \eng{older people} à \eng{the old}, jugé un peu brutal. De même, on évite souvent \eng{the disabled} au profit de \eng{people with disabilities} dans les documents officiels. La langue évolue avec le respect dû aux personnes !
\end{savaistu}

\courssub{Emploi 3 : caractéristiques physiques et situations}

La structure désigne aussi les personnes partageant une caractéristique physique ou une situation particulière :

\begin{itemize}
\item \eng{the blind} = les aveugles ; \eng{the deaf} = les sourds ; \eng{the disabled} = les handicapés.
\item \eng{the sick} = les malades ; \eng{the wounded} = les blessés (à la guerre) ; \eng{the injured} = les blessés (accident).
\item \eng{the dead} = les morts ; \eng{the living} = les vivants.
\end{itemize}

\eng{The dead cannot speak, but the living must remember them.} = « Les morts ne peuvent pas parler, mais les vivants doivent se souvenir d'eux. »

\eng{In Buea, a special school takes care of the deaf.} = « À Buea, une école spéciale prend soin des sourds. »

\courssub{Emploi 4 : les nationalités en -sh, -ch, -ese}

Pour les nationalités, la structure \eng{the + adjectif} désigne \textbf{tout le peuple} d'un pays — mais seulement pour les adjectifs de nationalité se terminant par \textbf{-sh}, \textbf{-ch} ou \textbf{-ese} (et pour \eng{the Swiss}) :

\begin{center}
\begin{tabularx}{\linewidth}{|X|X|X|}
\hline
\rowcolor{popblueL} English & Français & Terminaison \\
\hline
the English & les Anglais & -sh \\
\hline
the Irish & les Irlandais & -sh \\
\hline
the Welsh & les Gallois & -sh \\
\hline
the Spanish & les Espagnols & -sh \\
\hline
the French & les Français & -ch \\
\hline
the Dutch & les Néerlandais & -tch (cas voisin) \\
\hline
the Swiss & les Suisses & -ss (invariable) \\
\hline
the Chinese & les Chinois & -ese \\
\hline
the Japanese & les Japonais & -ese \\
\hline
the Portuguese & les Portugais & -ese \\
\hline
\end{tabularx}
\end{center}

\eng{The English love tea.} = « Les Anglais aiment le thé. »

\eng{The French and the English have a long common history.} = « Les Français et les Anglais ont une longue histoire commune. »

\begin{attention}[The Americans ? Non !]
Cette structure ne fonctionne \textbf{pas} avec les adjectifs de nationalité en \eng{-an}, \eng{-ian} ou \eng{-i}. Ceux-ci sont de \emph{vrais noms} qui prennent un \eng{-s} au pluriel :
\begin{itemize}
\item on dit \eng{the Americans} (avec un \eng{-s}), jamais \eng{the American} pour tout le peuple ;
\item de même : \eng{the Cameroonians}, \eng{the Nigerians}, \eng{the Germans}, \eng{the Pakistanis}.
\end{itemize}
Comparez : \eng{The Chinese are hard-working.} (adjectif en \eng{-ese}, invariable) mais \eng{The Cameroonians are hard-working.} (nom pluriel régulier).
\end{attention}

\courssub{Exception : pour parler d'UNE seule personne}

La structure \eng{the + adjectif} est \textbf{toujours collective}. Elle ne peut jamais désigner une seule personne. Pour parler d'un individu, il faut ajouter un \textbf{nom} après l'adjectif :

\begin{center}
\begin{tabularx}{\linewidth}{|X|X|}
\hline
\rowcolor{popblueL} Le groupe (pluriel) & Une personne (singulier) \\
\hline
the poor & a poor person / a poor man / a poor woman \\
\hline
the blind & a blind man / a blind girl \\
\hline
the young & a young man / a young lady \\
\hline
the unemployed & an unemployed worker \\
\hline
the French & a Frenchman / a Frenchwoman / a French citizen \\
\hline
the English & an Englishman / an Englishwoman \\
\hline
the Japanese & a Japanese student (adjectif + nom) \\
\hline
\end{tabularx}
\end{center}

\begin{attention}[Jamais « a poor » !]
Ne dites jamais \eng{a poor} ou \eng{a blind} pour parler d'une personne : l'adjectif anglais ne peut pas être précédé de \eng{a/an} sans nom. Dites \eng{a poor person}, \eng{a blind man}. Erreur typique du francophone qui traduit mot à mot « un pauvre ». En revanche, pour les nationalités en \eng{-an}, le nom existe : \eng{an American}, \eng{a Cameroonian} sont corrects.
\end{attention}

\courssub{Emploi abstrait : le verbe au SINGULIER}

Il existe un second emploi, plus rare mais très élégant, où \eng{the + adjectif} ne désigne pas des personnes mais une \textbf{notion abstraite}. Dans ce cas, le verbe se met au \textbf{singulier} :

\begin{itemize}
\item \eng{the good} = ce qui est bien, le bien ; \eng{the beautiful} = le beau ; \eng{the true} = le vrai.
\item \eng{the unknown} = l'inconnu ; \eng{the impossible} = l'impossible ; \eng{the supernatural} = le surnaturel.
\end{itemize}

\eng{The unknown worries him.} = « L'inconnu l'inquiète. » (verbe \eng{worries} avec un \eng{-s} : singulier !)

\eng{The beautiful does not always cost money.} = « Le beau ne coûte pas toujours de l'argent. »

\eng{We must choose between the good and the bad.} = « Nous devons choisir entre le bien et le mal. »

\begin{popfigure}
\begin{center}
\begin{tikzpicture}
\node[draw=popblue, fill=popblueL, rounded corners, thick, align=center, text width=6.2cm, minimum height=3.6cm] (A) at (0,0)
{\textbf{THE + ADJECTIVE}\\[0.15cm]
\textbf{Sens collectif}\\[0.1cm]
= un groupe de personnes\\
\eng{the poor} = \eng{poor people}\\
\eng{the young} = \eng{young people}\\[0.15cm]
Verbe au \textbf{PLURIEL}\\
\eng{The poor NEED help.}};
\node[draw=poporange, fill=poporangeL, rounded corners, thick, align=center, text width=6.2cm, minimum height=3.6cm] (B) at (9,0)
{\textbf{THE + ADJECTIVE}\\[0.15cm]
\textbf{Sens abstrait}\\[0.1cm]
= une notion, une qualité\\
\eng{the beautiful} = what is beautiful\\
\eng{the unknown} = what is unknown\\[0.15cm]
Verbe au \textbf{SINGULIER}\\
\eng{The unknown FRIGHTENS us.}};
\draw[->, very thick, popdark] (A) -- (B) node[midway, above, align=center, text width=2.2cm]{\footnotesize même forme, deux sens};
\end{tikzpicture}
\end{center}
\legende{Les deux sens de « the + adjective » : collectif (pluriel) ou abstrait (singulier).}
\end{popfigure}

\courssub{Cas particulier utile à l'examen : les adjectifs devenus de vrais noms}

Certains adjectifs sont devenus de véritables \textbf{noms} : ils peuvent alors prendre \eng{a/an} au singulier et un \eng{-s} au pluriel. Attention à ne pas les confondre avec la structure collective :

\begin{center}
\begin{tabularx}{\linewidth}{|X|X|X|}
\hline
\rowcolor{popblueL} Singulier & Pluriel & Français \\
\hline
a criminal & criminals & un criminel / des criminels \\
\hline
a professional & professionals & un professionnel / des professionnels \\
\hline
a liberal & liberals & un libéral / des libéraux \\
\hline
a conservative & conservatives & un conservateur / des conservateurs \\
\hline
an adult & adults & un adulte / des adultes \\
\hline
an intellectual & intellectuals & un intellectuel / des intellectuels \\
\hline
\end{tabularx}
\end{center}

\eng{Young professionals are starting businesses in Yaoundé.} = « De jeunes professionnels créent des entreprises à Yaoundé. »

Comparez bien : \eng{the rich} (groupe entier, pas de \eng{-s}) mais \eng{criminals} (vrai nom, pluriel avec \eng{-s}).

\courssub{Exercice résolu et piège final}

\begin{exoresolu}[Accordez le verbe et corrigez si nécessaire]
Complétez avec la bonne forme du verbe entre parenthèses, ou corrigez la phrase :\\
1. \eng{The unemployed (to be) worried about the future.}\\
2. \eng{The unknown (to frighten) many people.}\\
3. \eng{The poors need our help.}\\
4. \eng{I saw a blind yesterday at the market.}
\tcblower
\textbf{Solution détaillée :}\\
1. \eng{The unemployed ARE worried about the future.} — sens collectif (les chômeurs), donc verbe au pluriel.\\
2. \eng{The unknown FRIGHTENS many people.} — sens abstrait (l'inconnu), donc verbe au singulier avec \eng{-s}.\\
3. Faux : on ne met jamais de \eng{-s} à l'adjectif. Correction : \eng{The POOR need our help.}\\
4. Faux : pour une seule personne, il faut ajouter un nom. Correction : \eng{I saw a BLIND MAN (or: a blind person) at the market yesterday.}
\end{exoresolu}

\begin{aretenir}[L'essentiel de la section]
\begin{itemize}
\item \eng{The + adjectif} = tout un groupe de personnes ; verbe au \textbf{pluriel} : \eng{The elderly are respected.}
\item Quatre grands emplois : catégories sociales (\eng{the rich}, \eng{the homeless}), âge (\eng{the young}, \eng{the elderly}), caractéristiques physiques (\eng{the blind}, \eng{the sick}), nationalités en \eng{-sh / -ch / -ese} (\eng{the English}, \eng{the Chinese}).
\item Jamais de \eng{-s} sur l'adjectif ; jamais \eng{a + adjectif} seul pour une personne : dites \eng{a poor person}.
\item Sens abstrait possible, avec verbe au \textbf{singulier} : \eng{The unknown frightens us.}
\item Certains adjectifs sont devenus de vrais noms : \eng{a criminal}, \eng{criminals}.
\end{itemize}
\end{aretenir}

\coursec{Comparaison français/anglais et pièges}

Nous venons de voir la règle générale : en anglais, \eng{the + adjective} désigne un groupe de personnes (\eng{the poor}, \eng{the young}, \eng{the unemployed}). Mais pour un élève francophone, le vrai danger n'est pas de comprendre la règle : c'est de la \textbf{trahir au moment de traduire}, parce que le français et l'anglais ne construisent pas du tout cette idée de la même façon. Cette section est donc consacrée aux \cle{différences} entre les deux langues et aux \cle{pièges} qu'elles tendent. Observons d'abord, traduisons ensuite.

\courssub{Premier choc : le français ajoute un -s, l'anglais jamais}

\begin{exemplebox}[Observation : deux langues, deux logiques]
Comparons ces paires de phrases :

\eng{The rich get richer and the poor get poorer.} « Les riches s'enrichissent et les pauvres s'appauvrissent. »

\eng{The young today love technology.} « Les jeunes d'aujourd'hui aiment la technologie. »

\eng{The unemployed need government support.} « Les chômeurs ont besoin de l'aide de l'État. »

\eng{The blind learn to read with their fingers.} « Les aveugles apprennent à lire avec leurs doigts. »

Remarquez : en français, l'adjectif prend la marque du pluriel (\emph{pauvres}, \emph{jeunes}, \emph{aveugles}). En anglais, l'adjectif reste \textbf{toujours invariable} : \eng{the poor}, jamais \eng{the poors}.
\end{exemplebox}

\begin{propriete}[Règle : invariabilité de l'adjectif-nom anglais]
La structure anglaise est : \textbf{the + adjectif (sans -s)}. L'adjectif anglais est par nature invariable ; même employé comme nom de groupe, il ne prend jamais la marque du pluriel. Le verbe qui suit, lui, est au \textbf{pluriel} : \eng{The poor \textbf{are} many.} « Les pauvres sont nombreux. »

\begin{center}
\begin{tabularx}{\linewidth}{|X|X|X|}
\hline
\rowcolor{popblueL} Français & English & Remarque \\
\hline
les pauvres & \eng{the poor} & jamais \eng{the poors} \\
\hline
les riches & \eng{the rich} & jamais \eng{the riches} \\
\hline
les jeunes & \eng{the young} & jamais \eng{the youngs} \\
\hline
les vieux / les personnes âgées & \eng{the old} / \eng{the elderly} & \eng{the elderly} est plus poli \\
\hline
les chômeurs & \eng{the unemployed} & participe passé employé comme adjectif \\
\hline
les aveugles & \eng{the blind} & \\
\hline
les sourds & \eng{the deaf} & \\
\hline
les malades & \eng{the sick} & \\
\hline
les blessés & \eng{the injured} & \\
\hline
les morts & \eng{the dead} & \\
\hline
\end{tabularx}
\end{center}
\end{propriete}

\begin{attention}[Erreur n°1 des francophones : ajouter un -s]
C'est la faute la plus fréquente dans les copies camerounaises : le français dit « les pauvres », donc l'élève écrit \eng{the poors}. \textbf{Faux !}

\begin{itemize}
\item \eng{The poors suffer a lot.} ✗ → \eng{The poor suffer a lot.} ✓ « Les pauvres souffrent beaucoup. »
\item \eng{The youngs like music.} ✗ → \eng{The young like music.} ✓ « Les jeunes aiment la musique. »
\item \eng{The richs live in Bastos.} ✗ → \eng{The rich live in Bastos.} ✓ « Les riches habitent à Bastos. »
\end{itemize}

Astuce : en anglais, \textbf{aucun adjectif ne prend jamais de -s}, quelle que soit sa fonction. Retenez : \emph{adjectives never take -s in English}.
\end{attention}

\courssub{Deuxième choc : « un pauvre » ne se traduit pas par « a poor »}

\begin{exemplebox}[Observation]
\eng{A poor man came to my office this morning.} « Un pauvre (homme) est venu à mon bureau ce matin. »

\eng{She married a rich man from Douala.} « Elle a épousé un riche (homme) de Douala. »

\eng{The hospital received an injured person.} « L'hôpital a reçu un blessé. »
\end{exemplebox}

\begin{propriete}[Règle : au singulier, l'anglais exige un nom]
Le français peut nominaliser un adjectif au singulier : \emph{un pauvre, une aveugle, un blessé}. L'anglais le refuse : \eng{the + adjectif} désigne toujours un \textbf{groupe entier} (pluriel), jamais un individu. Pour parler d'\textbf{une seule personne}, il faut ajouter un nom : \eng{a poor man}, \eng{a poor woman}, \eng{a poor person}, \eng{an injured man}.

\begin{center}
\begin{tabularx}{\linewidth}{|X|X|X|}
\hline
\rowcolor{popblueL} Français & English & Remarque \\
\hline
un pauvre & \eng{a poor man} / \eng{a poor person} & jamais \eng{a poor} \\
\hline
une aveugle & \eng{a blind woman} & jamais \eng{a blind} \\
\hline
un chômeur & \eng{an unemployed man} / \eng{a jobless person} & \\
\hline
un blessé & \eng{an injured man} & \\
\hline
un malade & \eng{a sick person} / \eng{a patient} & \\
\hline
\end{tabularx}
\end{center}
\end{propriete}

\begin{attention}[Piège : « a poor » tout seul n'existe pas]
\eng{I saw a poor in the street.} ✗ → \eng{I saw a poor man in the street.} ✓ « J'ai vu un pauvre dans la rue. »

L'adjectif anglais ne peut pas porter l'article indéfini \eng{a/an} tout seul : il a besoin d'un nom derrière lui. Seule exception : quelques adjectifs devenus de vrais noms par l'usage (\eng{a daily} = un quotidien, \eng{a classic} = un classique), mais ce sont des cas figés à apprendre un par un.
\end{attention}

\courssub{Troisième choc : les nationalités}

\begin{exemplebox}[Observation]
\eng{The French love good food.} « Les Français aiment la bonne cuisine. »

\eng{A Frenchman opened a bakery in Yaoundé.} « Un Français a ouvert une boulangerie à Yaoundé. »

\eng{The English drink a lot of tea.} « Les Anglais boivent beaucoup de thé. »

\eng{Two Englishmen visited our school.} « Deux Anglais ont visité notre école. »
\end{exemplebox}

\begin{propriete}[Règle : nationalités en -sh, -ch, -se]
Pour les nationalités en \textbf{-sh, -ch, -se} (\eng{French, English, Welsh, Dutch, Spanish, Japanese, Chinese}), le groupe se dit \eng{the + adjectif} : \eng{the French} « les Français », \eng{the English} « les Anglais ». Mais au singulier, on ne dit jamais \eng{a French} : il faut \eng{a Frenchman / a Frenchwoman}, \eng{an Englishman / an Englishwoman}, ou plus neutre : \eng{a French person}, \eng{a Spanish student}.

Pour les autres nationalités (\eng{German, American, Cameroonian, Nigerian}), l'adjectif est déjà un vrai nom et prend un -s normal : \eng{a German}, \eng{two Germans}, \eng{the Germans}.

\begin{center}
\begin{tabularx}{\linewidth}{|X|X|X|}
\hline
\rowcolor{popblueL} Français & English (groupe) & English (individu) \\
\hline
les Français & \eng{the French} & \eng{a Frenchman} / \eng{a Frenchwoman} \\
\hline
les Anglais & \eng{the English} & \eng{an Englishman} / \eng{an Englishwoman} \\
\hline
les Espagnols & \eng{the Spanish} & \eng{a Spaniard} / \eng{a Spanish person} \\
\hline
les Japonais & \eng{the Japanese} & \eng{a Japanese} (invariable) \\
\hline
les Chinois & \eng{the Chinese} & \eng{a Chinese} (invariable) \\
\hline
les Néerlandais & \eng{the Dutch} & \eng{a Dutchman} / \eng{a Dutch person} \\
\hline
les Gallois & \eng{the Welsh} & \eng{a Welshman} / \eng{a Welsh person} \\
\hline
les Allemands & \eng{the Germans} & \eng{a German} \\
\hline
les Camerounais & \eng{the Cameroonians} & \eng{a Cameroonian} \\
\hline
\end{tabularx}
\end{center}
\end{propriete}

\begin{attention}[Piège : « a French » n'existe pas]
\eng{My neighbour is a French.} ✗ → \eng{My neighbour is a Frenchman.} ✓ ou \eng{My neighbour is French.} ✓ « Mon voisin est français / un Français. »

Deux solutions correctes : soit l'adjectif seul après \eng{to be} (\eng{He is French.}), soit le nom complet (\eng{He is a Frenchman.}). Mais jamais \eng{a French} !
\end{attention}

\courssub{Quatrième piège : l'article « the » change le sens}

\begin{exemplebox}[Observation : comparez]
\eng{Poor people suffer everywhere in the world.} « Les pauvres (en général, partout) souffrent dans le monde entier. »

\eng{The poor of this town suffer from the lack of water.} « Les pauvres de cette ville (un groupe précis) souffrent du manque d'eau. »
\end{exemplebox}

\begin{propriete}[Règle : avec ou sans « the »]
\begin{itemize}
\item \textbf{the + adjectif} = un groupe \textbf{défini}, que l'on peut identifier : \eng{the poor of this town}, \eng{the young of this school}.
\item \textbf{adjectif + people} (sans \eng{the}) = une généralité absolue : \eng{Poor people need help.} « Les pauvres (en général) ont besoin d'aide. »
\end{itemize}
Les deux sont corrects, mais le sens n'est pas identique : sans \eng{the}, l'adjectif redevient un simple adjectif épithète placé devant \eng{people}. En revanche, on ne peut jamais écrire \eng{Poors suffer} ✗ ni \eng{The poors suffer} ✗.
\end{propriete}

\begin{methode}[Le test du mot « people »]
Pour vérifier votre phrase, remplacez mentalement la structure par \eng{people} :

\begin{enumerate}
\item Vous voulez écrire \eng{the young today love technology}.
\item Testez : \eng{young people today love technology} — la phrase marche avec \eng{... people} ?
\item Oui → la structure \eng{the + adjectif} est correcte : \eng{The young today love technology.} ✓
\item Si vous hésitez entre \eng{the youngs} et \eng{the young}, rappelez-vous : l'adjectif anglais est invariable → \eng{the young}. ✓
\end{enumerate}
\end{methode}

\begin{popfigure}
\begin{center}
\begin{tikzpicture}[
  box/.style={draw=popblue, fill=popblueL, rounded corners=3pt, align=center, minimum width=4.6cm, minimum height=0.9cm, font=\small},
  ebox/.style={draw=poporange, fill=poporangeL, rounded corners=3pt, align=center, minimum width=4.6cm, minimum height=0.9cm, font=\small},
  fl/.style={-{Stealth[length=3mm]}, thick, poppurple}
]
\node[box] (f1) at (0,0) {\textbf{les riches}};
\node[ebox] (e1) at (7,0) {\textbf{\eng{the rich}} (jamais \emph{\eng{the richs}})};
\draw[fl] (f1) -- (e1);

\node[box] (f2) at (0,-1.4) {\textbf{un riche}};
\node[ebox] (e2) at (7,-1.4) {\textbf{\eng{a rich man}} (jamais \emph{\eng{a rich}})};
\draw[fl] (f2) -- (e2);

\node[box] (f3) at (0,-2.8) {\textbf{les Anglais}};
\node[ebox] (e3) at (7,-2.8) {\textbf{\eng{the English}}};
\draw[fl] (f3) -- (e3);

\node[box] (f4) at (0,-4.2) {\textbf{un Anglais}};
\node[ebox] (e4) at (7,-4.2) {\textbf{\eng{an Englishman}} (jamais \emph{\eng{an English}})};
\draw[fl] (f4) -- (e4);

\node[font=\bfseries, popblue] at (0,1) {FRANÇAIS};
\node[font=\bfseries, poporange] at (7,1) {ENGLISH};
\end{tikzpicture}
\end{center}
\legende{Tableau comparatif : le français nominalise l'adjectif librement ; l'anglais impose \eng{the + adjectif} pour le groupe et exige un nom pour l'individu.}
\end{popfigure}

\begin{exoresolu}[Exercice flash : repérez et corrigez les erreurs]
Énoncé — Voici cinq phrases écrites par des élèves francophones de Première. Chacune contient une erreur sur les adjectifs-noms. Trouvez l'erreur et corrigez-la.

\begin{enumerate}
\item \eng{The poors of our village need clean water.}
\item \eng{I met a blind at the market yesterday.}
\item \eng{The youngs today spend too much time on their phones.}
\item \eng{My uncle's new boss is a French.}
\item \eng{The riches should pay more taxes.}
\end{enumerate}

\tcblower
Solution détaillée :

\begin{enumerate}
\item \eng{The poors} ✗ → \eng{The poor of our village need clean water.} ✓ L'adjectif anglais est invariable : jamais de -s. « Les pauvres de notre village ont besoin d'eau potable. »
\item \eng{a blind} ✗ → \eng{I met a blind man (person) at the market yesterday.} ✓ Au singulier, l'anglais exige un nom après l'adjectif. « J'ai rencontré un aveugle au marché hier. »
\item \eng{The youngs} ✗ → \eng{The young today spend too much time on their phones.} ✓ « Les jeunes d'aujourd'hui passent trop de temps sur leur téléphone. »
\item \eng{a French} ✗ → \eng{My uncle's new boss is a Frenchman.} ✓ (ou \eng{... is French.}) « Le nouveau patron de mon oncle est un Français. »
\item \eng{The riches} ✗ → \eng{The rich should pay more taxes.} ✓ « Les riches devraient payer plus d'impôts. »
\end{enumerate}
\end{exoresolu}

\begin{aretenir}[Les quatre réflexes à acquérir]
\begin{enumerate}
\item Groupe de personnes : \eng{the + adjectif}, \textbf{sans -s} : \eng{the poor}, \eng{the young}, \eng{the unemployed}.
\item Une seule personne : adjectif + nom : \eng{a poor man}, \eng{a blind woman}, jamais \eng{a poor}.
\item Nationalités en -sh/-ch/-se : groupe = \eng{the French} ; individu = \eng{a Frenchman} ou \eng{He is French}.
\item Sans \eng{the}, utilisez \eng{adjectif + people} pour une généralité : \eng{Poor people suffer.} ≠ \eng{The poor of this town suffer.}
\end{enumerate}
\end{aretenir}

\begin{exercice}[À vous de jouer]
Traduisez en anglais :
\begin{enumerate}
\item Les jeunes d'aujourd'hui aiment la technologie.
\item Un blessé a été transporté à l'hôpital de Buea.
\item Les Anglais adorent le football.
\item J'ai parlé à un Anglais au marché de Bamenda.
\item Les chômeurs demandent du travail.
\end{enumerate}
\end{exercice}

\begin{corrige}[Exercice « À vous de jouer »]
\begin{enumerate}
\item \eng{The young today love technology.}
\item \eng{An injured man was taken to the Buea hospital.}
\item \eng{The English love football.}
\item \eng{I spoke to an Englishman at the Bamenda market.}
\item \eng{The unemployed are asking for jobs.}
\end{enumerate}
\end{corrige}

\coursec{Observation : « So do I » / « Neither did he »}

Après avoir comparé le français et l'anglais pour les adjectifs employés comme noms, et repéré les pièges qui guettent les francophones (\eng{the poors} au lieu de \eng{the poor}, par exemple), nous allons maintenant aborder une deuxième structure très fréquente dans la conversation quotidienne : les \cle{réponses courtes} qui permettent de dire « moi aussi » ou « moi non plus » en anglais. Là encore, l'anglais procède d'une manière très différente du français, et c'est précisément en observant d'abord des exemples authentiques que nous découvrirons la règle.

\courssub{Un dialogue au marché de Yaoundé}

\begin{textebox}[At Mokolo Market, Yaoundé]
Two traders, Mama Ngo and Mr Ewusi, meet at Mokolo Market in Yaoundé. Listen to their conversation.

A: I sell cocoa from the Centre region. What about you?

B: So do I! I have a small farm near Obala.

A: Really? I don't sell coffee this year. The price is too low.

B: Neither do I. My brother grows coffee, but I prefer cocoa.

A: I work very hard, you know. I wake up at five every morning.

B: So does my wife! She opens her shop before six.

A: I didn't go to the farm on Saturday. It was raining.

B: Neither did I. I stayed at home with my children.

A: I can speak a little English with the buyers from Nigeria.

B: So can my daughter! She learnt it at school.

A: Well, I'm not tired today. Business is good!

B: Neither am I. Let's have a cup of tea together.
\end{textebox}

\textbf{Glossaire.}

\begin{center}
\begin{tabularx}{\linewidth}{|l|l|X|}
\hline
\rowcolor{popblueL} \textbf{Mot anglais} & \textbf{Nature} & \textbf{Traduction} \\
\hline
trader & nom & commerçant, commerçante \\
\hline
to sell & verbe & vendre \\
\hline
cocoa & nom & cacao \\
\hline
to grow & verbe & cultiver \\
\hline
to wake up & verbe & se réveiller \\
\hline
buyer & nom & acheteur \\
\hline
business & nom & les affaires, le commerce \\
\hline
\end{tabularx}
\end{center}

\courssub{Questions d'observation}

Avant d'énoncer la règle, observons attentivement les réponses de B dans le dialogue. Répondez mentalement aux questions suivantes :

\begin{enumerate}
\item Dans la réponse \eng{So do I!}, la phrase de départ (\eng{I sell cocoa}) est-elle affirmative ou négative ?
\item Dans la réponse \eng{Neither do I.}, la phrase de départ (\eng{I don't sell coffee}) est-elle affirmative ou négative ?
\item Dans \eng{So do I}, où se trouve l'auxiliaire \eng{do} : avant ou après le sujet \eng{I} ?
\item Dans \eng{So does my wife}, pourquoi utilise-t-on \eng{does} et non \eng{do} ?
\item Dans \eng{Neither did I}, d'où vient l'auxiliaire \eng{did} ? Regardez la phrase de départ : \eng{I didn't go to the farm}.
\item Dans \eng{So can my daughter} et \eng{Neither am I}, quels auxiliaires sont repris ?
\end{enumerate}

Vous avez certainement remarqué les points essentiels : \cle{so} suit une phrase \textbf{affirmative}, \cle{neither} (ou \cle{nor}) suit une phrase \textbf{négative}, et l'auxiliaire se place \textbf{avant} le sujet. C'est exactement ce que nous allons formaliser.

\courssub{Le repérage : so après l'affirmatif, neither après le négatif}

\begin{exemplebox}[Deux exemples traduits pour bien comprendre]
\eng{I like my job. — So do I.} = « J'aime mon métier. — Moi aussi. »

\eng{He didn't come. — Neither did she.} = « Il n'est pas venu. — Elle non plus. »
\end{exemplebox}

La logique est simple et ne souffre aucune exception :

\begin{aretenir}[Le choix entre so et neither]
\begin{itemize}
\item Phrase de départ \textbf{affirmative} $\rightarrow$ réponse avec \cle{so} : \eng{So do I.} (moi aussi).
\item Phrase de départ \textbf{négative} $\rightarrow$ réponse avec \cle{neither} ou \cle{nor} : \eng{Neither do I.} / \eng{Nor do I.} (moi non plus).
\end{itemize}
\eng{Neither} et \eng{nor} sont parfaitement interchangeables ici ; \eng{neither} est un peu plus fréquent dans la conversation. Attention : après \eng{neither/nor}, la phrase est déjà négative de sens, on ne remet donc \textbf{jamais} \eng{not} : on dit \eng{Neither did I}, jamais \eng{Neither didn't I}.
\end{aretenir}

\courssub{L'inversion : l'auxiliaire avant le sujet}

Deuxième observation capitale : dans ces réponses courtes, l'anglais place l'\textbf{auxiliaire avant le sujet}, comme dans une question. On dit \eng{So do I}, et non \eng{So I do}. Cette inversion est obligatoire dans la réponse courte d'accord.

\begin{center}
\begin{tabularx}{\linewidth}{|X|X|X|}
\hline
\rowcolor{popblueL} \textbf{Français} & \textbf{English} & \textbf{Remarque} \\
\hline
Moi aussi. & \eng{So do I.} & inversion : auxiliaire + sujet \\
\hline
Moi non plus. & \eng{Neither do I.} & inversion : auxiliaire + sujet \\
\hline
Lui aussi. & \eng{So does he.} & \eng{does} à la 3\textsuperscript{e} personne \\
\hline
Elle non plus. & \eng{Neither did she.} & \eng{did} pour le prétérit \\
\hline
\end{tabularx}
\end{center}

\courssub{Le choix de l'auxiliaire : on reprend celui de la phrase de départ}

Troisième observation : l'auxiliaire de la réponse courte \textbf{reprend celui de la phrase de départ}. Regardons les cas possibles.

\begin{propriete}[Quel auxiliaire choisir ?]
\begin{enumerate}
\item Si la phrase de départ contient un auxiliaire ou un verbe modal (\eng{is, are, was, has, can, will, must, would}\ldots), on le \textbf{reprend tel quel} :
\begin{itemize}
\item \eng{I can swim. — So can I.} (« Je sais nager. — Moi aussi. »)
\item \eng{She is tired. — So am I.} (« Elle est fatiguée. — Moi aussi. »)
\item \eng{They have finished. — So have we.} (« Ils ont fini. — Nous aussi. »)
\item \eng{He will come. — So will she.} (« Il viendra. — Elle aussi. »)
\item \eng{I wasn't there. — Neither was he.} (« Je n'y étais pas. — Lui non plus. »)
\end{itemize}
\item Si la phrase de départ contient un \textbf{verbe lexical sans auxiliaire apparent} (présent simple ou prétérit simple), on utilise \cle{do}, \cle{does} ou \cle{did} selon le temps et la personne :
\begin{itemize}
\item \eng{She works hard. — So does he.} (« Elle travaille dur. — Lui aussi. ») : présent simple, 3\textsuperscript{e} personne $\rightarrow$ \eng{does}.
\item \eng{They play football. — So do we.} (« Ils jouent au football. — Nous aussi. ») : présent simple $\rightarrow$ \eng{do}.
\item \eng{He bought a bendskin. — So did his friend.} (« Il a acheté une moto-taxi. — Son ami aussi. ») : prétérit $\rightarrow$ \eng{did}.
\end{itemize}
\end{enumerate}
\end{propriete}

\begin{center}
\begin{tabular}{|l|l|l|}
\hline
\rowcolor{popblueL} \textbf{Phrase de départ} & \textbf{Réponse « aussi »} & \textbf{Réponse « non plus »} \\
\hline
\eng{I work in a bank.} & \eng{So do I.} & --- \\
\hline
\eng{I don't work on Sundays.} & --- & \eng{Neither do I.} \\
\hline
\eng{She lives in Douala.} & \eng{So does her sister.} & --- \\
\hline
\eng{He didn't pass the exam.} & --- & \eng{Neither did his brother.} \\
\hline
\eng{I am a student.} & \eng{So am I.} & --- \\
\hline
\eng{I am not ready.} & --- & \eng{Neither am I.} \\
\hline
\eng{They have eaten.} & \eng{So have we.} & --- \\
\hline
\eng{She can drive.} & \eng{So can I.} & --- \\
\hline
\eng{We will travel.} & \eng{So will they.} & --- \\
\hline
\eng{He went to Buea.} & \eng{So did she.} & --- \\
\hline
\end{tabular}
\end{center}

\courssub{Schéma de décision}

\begin{popfigure}
\begin{tikzpicture}[
  node distance=1.5cm and 2.5cm,
  startstop/.style={rectangle, rounded corners, draw=popblue, fill=popblueL, text width=3.8cm, align=center, minimum height=1cm, font=\small},
  decision/.style={diamond, draw=poporange, fill=poporangeL, text width=2.8cm, align=center, inner sep=1pt, font=\small},
  result/.style={rectangle, rounded corners, draw=popgreen, fill=popgreenL, text width=4.5cm, align=center, minimum height=1cm, font=\small},
  resultb/.style={rectangle, rounded corners, draw=poppurple, fill=poppurpleL, text width=4.5cm, align=center, minimum height=1cm, font=\small},
  arrow/.style={-{Stealth[length=3mm]}, thick, draw=popdark}
]
\node[startstop, align=center] (start){Phrase de départ\\ \eng{I sell cocoa.}};
\node[decision, below=of start] (dec) {Affirmative ?};
\node[result, below left=of dec, align=center] (yes){\textbf{SO} + auxiliaire + sujet\\ \eng{So do I.}};
\node[resultb, below right=of dec, align=center] (no){\textbf{NEITHER / NOR} + auxiliaire + sujet\\ \eng{Neither do I.}};
\draw[arrow] (start) -- (dec);
\draw[arrow] (dec) -- node[above left, font=\small\itshape]{oui} (yes);
\draw[arrow] (dec) -- node[above right, font=\small\itshape]{non} (no);
\end{tikzpicture}
\legende{Organigramme de choix : \eng{so} après une phrase affirmative, \eng{neither/nor} après une phrase négative, toujours avec inversion de l'auxiliaire.}
\end{popfigure}

\courssub{Première formulation de la règle par les élèves}

\begin{methode}[À vous de formuler la règle]
Avant de lire la règle officielle, essayez de compléter cette formulation avec votre classe :
\begin{enumerate}
\item Pour dire « moi aussi » après une phrase affirmative, j'utilise \cle{so} + \ldots + \ldots
\item Pour dire « moi non plus » après une phrase négative, j'utilise \cle{neither} ou \cle{nor} + \ldots + \ldots
\item L'auxiliaire que je choisis est celui de \ldots ; s'il n'y en a pas, j'utilise \ldots
\item L'ordre des mots est : auxiliaire \ldots{} le sujet (inversion).
\end{enumerate}
Formulation attendue : \eng{So / Neither + auxiliaire (celui de la phrase de départ, ou do/does/did) + sujet}, avec l'auxiliaire placé avant le sujet.
\end{methode}

\begin{propriete}[Règle : les réponses courtes d'accord]
\textbf{Forme} : \cle{so} / \cle{neither} / \cle{nor} + auxiliaire + sujet.

\textbf{Emploi} : pour exprimer l'accord avec ce qui vient d'être dit, sans répéter toute la phrase.
\begin{itemize}
\item \eng{I like my job. — So do I.} (« J'aime mon métier. — Moi aussi. »)
\item \eng{He didn't come. — Neither did she.} (« Il n'est pas venu. — Elle non plus. »)
\end{itemize}

\textbf{Exceptions et points délicats} :
\begin{itemize}
\item Pas de double négation : \eng{Neither did I}, jamais \eng{Neither didn't I}.
\item Le verbe \eng{be} seul sert d'auxiliaire : \eng{I'm tired. — So am I.}
\item Au présent simple, n'oubliez pas le \eng{-s} de \eng{does} à la 3\textsuperscript{e} personne : \eng{So does she.}
\end{itemize}
\end{propriete}

\begin{exoresolu}[Complétez les réponses courtes]
Énoncé. Complétez chaque réponse avec \eng{so} ou \eng{neither} et le bon auxiliaire.

1. A: I speak French. B: \ldots\ldots\ldots\ldots{} I.

2. A: She didn't buy any plantains. B: \ldots\ldots\ldots\ldots{} her husband.

3. A: We are going to Bamenda tomorrow. B: \ldots\ldots\ldots\ldots{} we.

4. A: My father can't drive. B: \ldots\ldots\ldots\ldots{} mine.

5. A: The students passed the test. B: \ldots\ldots\ldots\ldots{} the teacher's son.
\tcblower
Solution détaillée.

1. \eng{So do I.} Phrase affirmative au présent simple, verbe lexical \eng{speak} $\rightarrow$ auxiliaire \eng{do}.

2. \eng{Neither did her husband.} Phrase négative au prétérit (\eng{didn't buy}) $\rightarrow$ \eng{neither} + \eng{did}.

3. \eng{So are we.} Phrase affirmative avec \eng{be going to} ; l'auxiliaire est \eng{are}, on le reprend.

4. \eng{Neither can mine.} Phrase négative avec le modal \eng{can't} $\rightarrow$ \eng{neither} + \eng{can} ; \eng{mine} = « le mien ».

5. \eng{So did the teacher's son.} Phrase affirmative au prétérit simple (\eng{passed}) $\rightarrow$ auxiliaire \eng{did}.
\end{exoresolu}

\begin{attention}[Ne confondez pas « So do I » et « So I do » !]
L'ordre \eng{So I do} (sans inversion) \textbf{existe} en anglais, mais il a un sens complètement différent : il exprime l'\textbf{insistance}, « c'est vrai que je le fais ! », « mais si, je le fais ! ».

Comparez :
\begin{itemize}
\item \eng{I work hard. — So do I.} = « Je travaille dur. — Moi aussi. » (accord, même situation)
\item \eng{You never help me! — So I do!} = « Tu ne m'aides jamais ! — Mais si, je t'aide ! » (insistance, même sujet)
\end{itemize}

Autre piège majeur pour les francophones : à l'écrit formel et dans les exercices de grammaire, on évite \eng{Me too} (familier) et on utilise la structure complète \eng{So do I}. Ne dites jamais \eng{So I am} pour « moi aussi » quand la phrase de départ a un verbe lexical : \eng{I like fish. — So do I} (et non \eng{So am I}). Enfin, rappelez-vous : après \eng{neither}, jamais de \eng{not} supplémentaire.
\end{attention}

\begin{experience}[Chaîne orale en classe]
Le professeur dit une phrase, affirmative ou négative, en lien avec la vie économique : \eng{I sell tomatoes at the market.}, \eng{I don't like waking up early.}, \eng{I can count money very fast.}, \eng{I didn't go to school on Saturday.} Chaque élève, à tour de rôle, répond par une réponse courte correcte (\eng{So do I!}, \eng{Neither can I.}\ldots), puis invente une nouvelle phrase pour l'élève suivant. Objectif : réagir en moins de trois secondes avec le bon auxiliaire et l'inversion correcte.
\end{experience}

\begin{savaistu}[Pourquoi cette inversion ?]
L'inversion de l'auxiliaire dans \eng{So do I} n'est pas une bizarrerie isolée : c'est la même logique que dans les questions (\eng{Do you speak English?}) et que dans les questions de confirmation (\eng{You live in Yaoundé, don't you?}). En anglais, dès qu'une phrase courte commence par certains mots comme \eng{so}, \eng{neither}, \eng{nor} ou \eng{never}, l'auxiliaire se place avant le sujet. C'est un réflexe fondamental de la syntaxe anglaise que vous retrouverez dans de nombreuses autres structures au fil de vos études.
\end{savaistu}

\coursec{Lesson 21 — Grammar: Adjectives as nouns; “So do I”, “Neither did he”}

\courssub{Adjectives used as nouns : “The + adjective”}

\begin{textebox}[Observation : Article de presse — The widening gap]
\textbf{The rich} get richer while \textbf{the poor} struggle to survive. In Yaoundé, \textbf{the unemployed} queue at the National Employment Fund every morning. \textbf{The young} dream of stable contracts, but \textbf{the elderly} often rely on informal trade. Meanwhile, \textbf{the disabled} face accessibility barriers in public buildings. The government promises to support \textbf{the needy} and protect \textbf{the vulnerable}.
\end{textebox}

\begin{definition}[Adjectifs substantivés : “The + adjective”]
En anglais, l'article défini \eng{the} suivi d'un \textbf{adjectif} (sans nom) fonctionne comme un \textbf{nom collectif} désignant un groupe de personnes partageant cette caractéristique.
\begin{itemize}
\item \eng{The rich} = les riches (le groupe des personnes riches).
\item \eng{The unemployed} = les chômeurs.
\item \eng{The young} = les jeunes.
\end{itemize}
\emph{Traduction française} : souvent \textit{les + adjectif} (substantivé) ou \textit{ceux qui sont + adjectif}.
\end{definition}

\begin{propriete}[Règles de base : accord et sens]
\begin{enumerate}
\item \textbf{Sens collectif, verbe au pluriel.} \eng{The poor \textbf{are} suffering.} (Les pauvres souffrent.) — On ne dit jamais \eng{The poors} ni \eng{The poor is}.
\item \textbf{Pas de “ones” / “people” après l'adjectif.} \eng{The injured \textbf{were taken} to hospital.} ✓ — \eng{The injured people...} possible mais lourd ; \eng{The injured ones...} ✗.
\item \textbf{Nationalités en \eng{-ch, -sh, -ese, -ss} :} \eng{the English, the French, the Spanish, the Chinese, the Swiss, the Dutch.} Verbe au pluriel. \eng{The French \textbf{are} proud of their cuisine.}
\item \textbf{Autres nationalités} : on utilise généralement le nom au pluriel (\eng{the Americans, the Cameroonians, the Germans}) ou \eng{the + adjective + people} (\eng{the Belgian people}).
\item \textbf{Valeur abstraite (singulier) :} \eng{the unknown, the unexpected, the best, the worst, the supernatural.} Verbe au singulier. \eng{The unknown \textbf{frightens} me.}
\end{enumerate}
\end{propriete}

\begin{center}
\begin{tabularx}{\linewidth}{|X|X|X|}
\hline
\rowcolor{popblueL} Français & English & Remarque \\
\hline
Les riches / les pauvres & The rich / the poor & Verbe au pluriel. Jamais \eng{the riches/poors}. \\
\hline
Les chômeurs & The unemployed & Préféré à \eng{the jobless}. \\
\hline
Les jeunes / les personnes âgées & The young / the elderly & \eng{The old} possible mais \eng{the elderly} plus poli. \\
\hline
Les handicapés & The disabled & \eng{The handicapped} vieilli / péjoratif. \\
\hline
Les sans-abri & The homeless & \\
\hline
Les nécessiteux / les vulnérables & The needy / the vulnerable & Contexte social / humanitaire. \\
\hline
L'inconnu / le meilleur / le pire & The unknown / the best / the worst & Sens abstrait $\rightarrow$ verbe au \textbf{singulier}. \\
\hline
Les Anglais / les Français / les Chinois & The English / the French / the Chinese & Adjectifs en \eng{-sh, -ch, -ese, -ss} invariables, pluriel. \\
\hline
\end{tabularx}
\end{center}

\begin{attention}[Pièges fréquents des francophones]
\begin{itemize}
\item \textbf{Ajouter un \eng{s} :} \eng{The poors are many.} ✗ $\rightarrow$ \eng{The poor are many.} ✓
\item \textbf{Verbe au singulier pour un groupe :} \eng{The unemployed is increasing.} ✗ $\rightarrow$ \eng{The unemployed are increasing.} ✓
\item \textbf{Confondre nationalité adjectivale et nominale :} \eng{“A French”} ✗ (un Français) $\rightarrow$ \eng{A Frenchman} / \eng{A French woman} ; mais \eng{The French} (les Français, collectif) ✓.
\item \textbf{Traduire “les jeunes” par \eng{the youngs} :} ✗ $\rightarrow$ \eng{The young} ✓.
\end{itemize}
\end{attention}

\begin{exemplebox}[Exemples contextualisés (Cameroun \& International)]
\begin{itemize}
\item \eng{The government must create jobs for \textbf{the unemployed}.} Le gouvernement doit créer des emplois pour les chômeurs.
\item \eng{\textbf{The young} represent 60\% of the population in Cameroon.} Les jeunes représentent 60\% de la population au Cameroun.
\item \eng{\textbf{The disabled} need better access to public transport.} Les personnes handicapées ont besoin d'un meilleur accès aux transports publics.
\item \eng{We must protect \textbf{the vulnerable} during the crisis.} Nous devons protéger les vulnérables pendant la crise.
\item \eng{\textbf{The best} is yet to come.} Le meilleur reste à venir. (Abstrait, singulier)
\end{itemize}
\end{exemplebox}

\begin{exercice}[Entraînement : Adjectifs comme noms]
Traduisez en anglais :\par
1. Les riches devraient aider les pauvres.\par
2. Les chômeurs manifestent dans les rues de Douala.\par
3. Les personnes âgées méritent notre respect.\par
4. L'inconnu l'effraie.\par
5. Les Anglais boivent beaucoup de thé.\par
6. Les sans-abri dorment sous les ponts.
\end{exercice}

\begin{corrige}[Exercice « Adjectifs comme noms »]
1. \eng{The rich should help the poor.}\par
2. \eng{The unemployed are demonstrating in the streets of Douala.}\par
3. \eng{The elderly deserve our respect.}\par
4. \eng{The unknown frightens him / her.}\par
5. \eng{The English drink a lot of tea.}\par
6. \eng{The homeless sleep under bridges.}
\end{corrige}

\begin{aretenir}[L'essentiel : “The + adjective”]
\begin{itemize}
\item \eng{The + adjectif} = nom collectif (groupe de personnes) $\rightarrow$ verbe au \textbf{pluriel}.
\item Pas de \eng{s} à l'adjectif, pas de \eng{ones/people} ajouté.
\item Nationalités \eng{-ch, -sh, -ese, -ss} : \eng{the French, the English...} (pluriel).
\item Sens abstrait (\eng{the best, the unknown...}) $\rightarrow$ verbe au \textbf{singulier}.
\end{itemize}
\end{aretenir}

\courssub{Short answers : “So do I” / “Neither did he”}

\begin{textebox}[Observation : Dialogue au marché de Mokolo — Accords et désaccords]
\textbf{A:} “I sell plantains and cassava.”\\
\textbf{B:} “\textbf{So do I.} Business is good this season.”\\
\textbf{A:} “I don't like the new tax.”\\
\textbf{B:} “\textbf{Neither do I.} It reduces our profit.”\\
\textbf{C (passant):} “I can't read the price list.”\\
\textbf{A:} “\textbf{Neither can I.} Let's ask the secretary.”\\
\textbf{B:} “I have paid my licence fee.”\\
\textbf{C:} “\textbf{So have I.} We are in order.”
\end{textebox}

\begin{propriete}[Accord affirmatif : So + auxiliaire + sujet]
Pour dire « moi aussi », « lui aussi », « eux aussi » après une phrase \textbf{affirmative}, l'anglais utilise la structure inversée :
\begin{center}
\textbf{SO + auxiliaire + sujet}
\end{center}
\begin{itemize}
\item \eng{“I work in a bank.” — “So do I.”} « Je travaille dans une banque. — Moi aussi. »
\item \eng{“She lives in Douala.” — “So does her brother.”} « Elle habite à Douala. — Son frère aussi. »
\item \eng{“They visited Buea last year.” — “So did we.”} « Ils ont visité Buea l'an dernier. — Nous aussi. »
\item \eng{“He is tired.” — “So am I.”} « Il est fatigué. — Moi aussi. »
\item \eng{“We have finished.” — “So have they.”} « Nous avons fini. — Eux aussi. »
\item \eng{“I can swim.” — “So can my sister.”} « Je sais nager. — Ma sœur aussi. »
\item \eng{“You will pass the exam.” — “So will you!”} « Tu réussiras l'examen. — Toi aussi ! »
\end{itemize}
\textbf{Attention à l'ordre des mots} : l'auxiliaire se place \textbf{AVANT} le sujet (inversion), comme dans une question. On dit \eng{So do I} et jamais \eng{So I do} pour exprimer l'accord.
\end{propriete}

\begin{attention}[Ne confondez pas “So do I” et “So I do”]
\eng{So do I} (inversion) = « moi aussi ». Mais \eng{So I do} (ordre normal) existe avec un autre sens : il \textbf{CONFIRME} ce qu'on vient de dire, souvent avec surprise : \eng{“You look tired.” — “So I am!”} « Tu as l'air fatigué. — C'est vrai, je le suis ! ». Pour l'examen, retenez : accord = inversion.
\end{attention}

\begin{propriete}[Accord négatif : Neither / Nor + auxiliaire + sujet]
Pour dire « moi non plus », « lui non plus » après une phrase \textbf{négative}, on utilise :
\begin{center}
\textbf{NEITHER (ou NOR) + auxiliaire + sujet}
\end{center}
\begin{itemize}
\item \eng{“I don't like crowded markets.” — “Neither do I.”} « Je n'aime pas les marchés bondés. — Moi non plus. »
\item \eng{“She didn't pass the test.” — “Nor did her classmates.”} « Elle n'a pas réussi le test. — Ses camarades non plus. »
\item \eng{“He isn't ready.” — “Neither am I.”} « Il n'est pas prêt. — Moi non plus. »
\item \eng{“They haven't paid yet.” — “Neither have we.”} « Ils n'ont pas encore payé. — Nous non plus. »
\item \eng{“I can't drive.” — “Nor can my father.”} « Je ne sais pas conduire. — Mon père non plus. »
\end{itemize}
\cle{Neither} et \cle{nor} sont parfaitement interchangeables ici ; \eng{neither} est un peu plus courant à l'oral.
\end{propriete}

\begin{attention}[Jamais de double négation avec neither]
\eng{Neither} porte déjà la négation. Le verbe qui suit reste donc à la forme \textbf{AFFIRMATIVE} :
\begin{itemize}
\item \eng{Neither don't I.} ✗ (double négation, faute grave)
\item \eng{Neither do I.} ✓
\item \eng{Neither didn't he come.} ✗ $\rightarrow$ \eng{Neither did he come.} ✓
\end{itemize}
C'est l'inverse du français, où « non plus » s'ajoute à une phrase déjà négative (« Moi, je ne viens pas non plus »). En anglais, une seule négation par phrase.
\end{attention}

\begin{propriete}[L'auxiliaire de la réponse courte est TOUJOURS celui de la phrase de départ]
\begin{enumerate}
\item Si la phrase de départ contient déjà un auxiliaire (\eng{be}, \eng{have}, \eng{can}, \eng{will}, \eng{must}...), on le \textbf{REPREND} tel quel, accordé avec le nouveau sujet : \eng{“She is working.” — “So am I.”}
\item Si la phrase de départ contient un \textbf{verbe lexical seul} (present simple ou past simple), on utilise \cle{DO / DOES / DID} : \eng{“He works hard.” — “So do I.”} ; \eng{“They left early.” — “So did we.”}
\end{enumerate}
\end{propriete}

\begin{center}
\begin{tabularx}{\linewidth}{|X|X|X|}
\hline
\rowcolor{popblueL} Phrase de départ (temps / verbe) & Auxiliaire repris & Exemple \\
\hline
Present simple (\eng{work, works}) & do / does & \eng{She works. — So does he.} \\
\hline
Past simple (\eng{worked}) & did & \eng{She worked. — So did he.} \\
\hline
\eng{be} au présent (\eng{am/is/are}) & am / is / are & \eng{She is busy. — So is he.} \\
\hline
\eng{be} au passé (\eng{was/were}) & was / were & \eng{She was late. — So was he.} \\
\hline
Present perfect (\eng{has/have worked}) & have / has & \eng{She has left. — So has he.} \\
\hline
\eng{can} & can & \eng{She can type. — So can he.} \\
\hline
\eng{will} & will & \eng{She will come. — So will he.} \\
\hline
\eng{must} & must & \eng{She must go. — So must he.} \\
\hline
\eng{should} & should & \eng{She should study. — So should he.} \\
\hline
\end{tabularx}
\end{center}

\begin{popfigure}
\begin{tikzpicture}[font=\small]
\draw[fill=popblueL, draw=popblue, rounded corners] (0,5.0) rectangle (3.4,5.8);
\node[align=center] at (1.7,5.4) {\textbf{Phrase de départ}\\ \textbf{(verbe)}};
\draw[fill=popgreenL, draw=popgreen, rounded corners] (4.1,5.0) rectangle (7.5,5.8);
\node[align=center] at (5.8,5.4) {\textbf{Accord affirmatif (So...)}};
\draw[fill=poppinkL, draw=poppink, rounded corners] (8.2,5.0) rectangle (11.6,5.8);
\node[align=center] at (9.9,5.4) {\textbf{Accord négatif (Neither...)}};
\foreach \y/\a/\b/\c in {
  4.2/works/So do I/Neither do I,
  3.4/is working/So am I/Neither am I,
  2.6/worked/So did I/Neither did I,
  1.8/has worked/So have I/Neither have I,
  1.0/can work/So can I/Neither can I,
  0.2/will work/So will I/Neither will I}{
  \draw[draw=popline, fill=white, rounded corners] (0,\y-0.3) rectangle (3.4,\y+0.3);
  \node at (1.7,\y) {\eng{\a}};
  \draw[draw=popline, fill=white, rounded corners] (4.1,\y-0.3) rectangle (7.5,\y+0.3);
  \node at (5.8,\y) {\eng{\b}};
  \draw[draw=popline, fill=white, rounded corners] (8.2,\y-0.3) rectangle (11.6,\y+0.3);
  \node at (9.9,\y) {\eng{\c}};
  \draw[-{Stealth}, popmuted] (3.4,\y) -- (4.1,\y);
  \draw[-{Stealth}, popmuted] (7.5,\y) -- (8.2,\y);
}
\end{tikzpicture}
\legende{De la phrase de départ à la réponse courte : l'auxiliaire est repris, le sujet change.}
\end{popfigure}

\begin{exemplebox}[Exemples traduits — Contexte vie économique]
\begin{itemize}
\item \eng{“She has finished her accounting training.” — “So have I.”} « Elle a fini sa formation comptable. — Moi aussi. »
\item \eng{“They won't strike tomorrow.” — “Neither will we.”} « Ils ne feront pas grève demain. — Nous non plus. »
\item \eng{“My father works at the port of Douala.” — “So does mine.”} « Mon père travaille au port de Douala. — Le mien aussi. »
\item \eng{“I wasn't at the trade fair yesterday.” — “Neither was I.”} « Je n'étais pas à la foire commerciale hier. — Moi non plus. »
\item \eng{“The traders will open early tomorrow.” — “So will the taxi drivers.”} « Les commerçants ouvriront tôt demain. — Les chauffeurs de taxi aussi. »
\item \eng{“He must renew his licence.” — “So must she.”} « Il doit renouveler sa licence. — Elle aussi. »
\end{itemize}
\end{exemplebox}

\courssub{Comparaison avec le français}

\begin{center}
\begin{tabularx}{\linewidth}{|X|X|X|}
\hline
\rowcolor{popblueL} Français & English & Remarque \\
\hline
Moi aussi. & So do I. / Me too. & Inversion obligatoire dans \eng{So do I}. \\
\hline
Moi non plus. & Neither do I. / Me neither. & Une seule négation en anglais. \\
\hline
Lui aussi (il travaille). & So does he. & \eng{does} à la 3e personne. \\
\hline
Nous aussi (nous sommes prêts). & So are we. & On reprend \eng{be}, pas \eng{do}. \\
\hline
Eux non plus (ils ne sont pas venus). & Neither did they. & \eng{did} car past simple lexical. \\
\hline
Elle aussi (elle a réussi). & So has she. & Present perfect $\rightarrow$ \eng{has}. \\
\hline
\end{tabularx}
\end{center}

Le français recycle la même formule (« moi aussi / moi non plus ») quelle que soit la phrase ; l'anglais, lui, \cle{adapte l'auxiliaire} à chaque fois. C'est cette gymnastique qu'il faut automatiser.

\begin{methode}[Trois étapes pour construire la réponse courte]
\begin{enumerate}
\item \textbf{Phrase affirmative ou négative ?} Affirmative $\rightarrow$ \eng{SO...} ; négative $\rightarrow$ \eng{NEITHER/NOR...}
\item \textbf{Quel auxiliaire ?} Cherchez l'auxiliaire de la phrase de départ (\eng{is, has, can, will, did}...). S'il n'y en a pas (verbe lexical seul), prenez \eng{do/does} (présent) ou \eng{did} (passé).
\item \textbf{Inversion.} Placez l'auxiliaire \textbf{AVANT} le nouveau sujet : \eng{So + aux + sujet}. Vérifiez l'accord : \eng{So does she} (et non \eng{So do she}).
\end{enumerate}
\end{methode}

\begin{exoresolu}[Choix de l'auxiliaire — entraînement rapide]
Réagissez à chaque réplique par une réponse courte d'accord avec le sujet indiqué.\par
1. \eng{“I work in a cocoa factory.”} (I) \quad 2. \eng{“Mary is learning bookkeeping.”} (her brother)\par
3. \eng{“We visited the Bamenda market.”} (they) \quad 4. \eng{“He has opened a shop.”} (she)\par
5. \eng{“I can speak two languages.”} (my friend) \quad 6. \eng{“They won't strike tomorrow.”} (we)\par
7. \eng{“She doesn't earn much.”} (I) \quad 8. \eng{“The boss wasn't angry.”} (the workers)
\tcblower
1. Phrase affirmative, verbe lexical au présent $\rightarrow$ \eng{do} : \eng{So do I.}\par
2. Affirmative, auxiliaire \eng{is} $\rightarrow$ \eng{So is her brother.}\par
3. Affirmative, past simple $\rightarrow$ \eng{did} : \eng{So did they.}\par
4. Affirmative, present perfect, sujet \eng{she} $\rightarrow$ \eng{So has she.}\par
5. Affirmative, modal \eng{can} $\rightarrow$ \eng{So can my friend.}\par
6. Négative, auxiliaire \eng{will} $\rightarrow$ \eng{Neither will we.}\par
7. Négative, verbe lexical au présent $\rightarrow$ \eng{Neither do I.} (jamais \eng{Neither don't I}).\par
8. Négative, \eng{be} au passé, sujet pluriel $\rightarrow$ \eng{Neither were the workers.}
\end{exoresolu}

\courssub{Variantes familières et expression du désaccord}

\begin{definition}[“Me too” et “Me neither”]
À l'oral familier, on entend très souvent \eng{Me too} (accord affirmatif) et \eng{Me neither} (accord négatif) : \eng{“I love this song.” — “Me too!”}. C'est correct en conversation, mais à \textbf{éviter dans la rédaction formelle} et à l'examen écrit, où l'on attend \eng{So do I} / \eng{Neither do I}. Notez aussi que \eng{Me too} ne peut \textbf{JAMAIS} exprimer l'accord après une phrase négative : on dit \eng{Me neither}, jamais \eng{Me too} après \eng{“I don't like it.”}
\end{definition}

\begin{propriete}[Exprimer le désaccord : “Oh, I do!” / “I don't.”]
Pour \textbf{CONTREDIRE} poliment, on reprend l'auxiliaire \textbf{sans inversion}, souvent introduit par \eng{Oh} ou \eng{But} :
\begin{itemize}
\item \eng{“You don't like your job.” — “Oh, I do! I love it.”} « Tu n'aimes pas ton travail. — Mais si ! Je l'adore. »
\item \eng{“You passed the exam!” — “I didn't, actually.”} « Tu as réussi ! — Non, en fait, j'ai échoué. »
\item \eng{“She can't drive.” — “Oh yes, she can!”} « Elle ne sait pas conduire. — Mais si ! »
\end{itemize}
L'accent tonique porte sur l'auxiliaire : c'est lui qui porte la contradiction, comme le « si » français.
\end{propriete}

\begin{attention}[Pièges fréquents des francophones — Récapitulatif]
\begin{itemize}
\item Traduire « moi aussi » mot à mot : \eng{Me also} ✗ $\rightarrow$ \eng{So do I} ou \eng{Me too} ✓.
\item Oublier l'inversion : \eng{So I do} ✗ (pour l'accord) $\rightarrow$ \eng{So do I} ✓.
\item Garder \eng{do} avec \eng{be} : \eng{“I am tired.” — “So do I.”} ✗ $\rightarrow$ \eng{So am I} ✓.
\item Double négation : \eng{Neither didn't he} ✗ $\rightarrow$ \eng{Neither did he} ✓.
\item Oublier l'accord à la 3e personne : \eng{So do she} ✗ $\rightarrow$ \eng{So does she} ✓.
\item Utiliser \eng{Me too} après une phrase négative : \eng{“I hate taxes.” — “Me too.”} ✗ $\rightarrow$ \eng{Me neither} ✓.
\end{itemize}
\end{attention}

\begin{savaistu}[“So do I” : la marque d'un bon candidat]
Dans les examens oraux britanniques (Cambridge, IELTS), les examinateurs apprécient particulièrement les candidats qui réagissent avec \eng{So do I} ou \eng{Neither can I} au lieu de répéter toute la phrase (\eng{I also like football}). Cette réponse courte prouve que l'élève maîtrise les auxiliaires — exactement ce que ce chapitre vous entraîne à faire. Au Cameroun, c'est aussi un point fréquent du paper de grammaire au Probatoire et au Baccalauréat : une transformation de phrase avec \eng{so} ou \eng{neither} rapporte souvent un point facile... à condition de ne pas oublier l'inversion !
\end{savaistu}

\begin{exercice}[À vous de jouer]
Réagissez par une réponse courte d'accord, avec le sujet entre parenthèses :\par
1. \eng{“I enjoy working with customers.”} (I) \quad 2. \eng{“The shop closes at six.”} (the bank)\par
3. \eng{“We didn't buy anything.”} (she) \quad 4. \eng{“He has found a new job.”} (his wife)\par
5. \eng{“I will travel to Buea.”} (we) \quad 6. \eng{“She isn't satisfied with her salary.”} (I)
\end{exercice}

\begin{corrige}[Exercice « À vous de jouer »]
1. \eng{So do I.} 2. \eng{So does the bank.} 3. \eng{Neither did she.} 4. \eng{So has his wife.} 5. \eng{So will we.} 6. \eng{Neither am I.}
\end{corrige}

\begin{aretenir}[L'essentiel de la leçon 21]
\begin{itemize}
\item \textbf{The + adjectif} = nom collectif pluriel (\eng{the poor, the unemployed, the English}) $\rightarrow$ verbe pluriel. Sens abstrait (\eng{the best, the unknown}) $\rightarrow$ verbe singulier.
\item \textbf{Accord affirmatif :} \textbf{SO + auxiliaire + sujet} (inversion).
\item \textbf{Accord négatif :} \textbf{NEITHER/NOR + auxiliaire + sujet} (inversion, pas de double négation).
\item \textbf{Auxiliaire} = celui de la phrase de départ ; verbe lexical seul $\rightarrow$ \eng{do/does/did}.
\item \textbf{Oral familier :} \eng{Me too} / \eng{Me neither} ; \textbf{Désaccord :} \eng{Oh, I do!} / \eng{I don't.}
\end{itemize}
\end{aretenir}

\coursec{Comparaison français/anglais et entraînement aux réponses courtes}

Nous savons maintenant construire les réponses courtes : \eng{So + auxiliaire + sujet} pour « moi aussi » et \eng{Neither + auxiliaire + sujet} pour « moi non plus ». Dans cette dernière partie, nous allons confronter directement les deux langues, identifier les pièges qui attendent le francophone, puis nous entraîner jusqu'à ce que la bonne réponse devienne un réflexe. Le thème reste celui du chapitre : les métiers et la vie économique.

\courssub{Le grand écart : « moi aussi » face à « So do I »}

\begin{exemplebox}[Observer avant de raisonner]
Comparez attentivement ces paires de phrases :
\begin{itemize}
\item \eng{My father works in a bank. — So does mine.} « Mon père travaille dans une banque. — Le mien aussi. »
\item \eng{I don't work on Sundays. — Neither do I.} « Je ne travaille pas le dimanche. — Moi non plus. »
\item \eng{My sister can drive. — So can mine.} « Ma sœur sait conduire. — La mienne aussi. »
\item \eng{The traders arrived early. — So did the farmers.} « Les commerçants sont arrivés tôt. — Les agriculteurs aussi. »
\end{itemize}
\emph{Note :} \eng{mine} (le mien / la mienne) est un pronom possessif qui remplace \eng{my father}, \eng{my sister}, etc.
\end{exemplebox}

Qu'observe-t-on ? Le français se contente de deux petits mots : \emph{moi aussi} ou \emph{moi non plus}. Pas de verbe, pas d'auxiliaire, pas d'inversion. L'anglais, lui, exige une \cle{structure complète} : un mot de liaison (\eng{so} ou \eng{neither}), un \cle{auxiliaire} qui reprend le temps et le mode de la phrase de départ, puis le \cle{sujet}, le tout dans l'ordre inversé (auxiliaire avant sujet).

\begin{propriete}[La règle d'or de la comparaison]
\begin{itemize}
\item \textbf{Français} : « moi aussi » / « moi non plus » — formule invariable, identique dans tous les contextes.
\item \textbf{Anglais} : \eng{So / Neither + auxiliaire + sujet} — formule variable, qui dépend de la phrase qui précède (auxiliaire, temps, modalité).
\end{itemize}
Conséquence pratique : on ne peut pas traduire « moi aussi » une fois pour toutes. Il faut d'abord \emph{regarder la phrase de départ}, en extraire l'auxiliaire, puis construire la réponse.
\end{propriete}

\begin{center}
\begin{tabularx}{\linewidth}{|X|X|X|}
\hline
\rowcolor{popblueL} Français & English & Remarque \\
\hline
« Moi aussi. » (après une phrase affirmative) & \eng{So do I.} / \eng{So am I.} / \eng{So did I.} & L'auxiliaire change selon la phrase de départ. \\
\hline
« Moi non plus. » (après une phrase négative) & \eng{Neither do I.} / \eng{Neither can I.} & Jamais de \eng{not} : \eng{neither} porte déjà la négation. \\
\hline
« Moi aussi » (oral courant) & \eng{Me too.} & Accepté à l'oral, à éviter dans une rédaction formelle. \\
\hline
« Moi non plus » (oral courant) & \eng{Me neither.} & Oral seulement ; préférer \eng{Neither do I.} à l'écrit. \\
\hline
\end{tabularx}
\end{center}

\begin{attention}[L'erreur numéro un du francophone : « I also »]
Beaucoup d'élèves traduisent mot à mot « moi aussi » par \eng{I also} ou écrivent \eng{I also like trading}. Cette phrase ne signifie pas « moi aussi » en réponse à quelqu'un : elle signifie « j'aime \emph{aussi} le commerce (en plus d'autre chose) ». Pour réagir à la phrase d'un partenaire, il faut dire :
\begin{itemize}
\item \eng{So do I.} (structure inversée, la plus élégante) ;
\item ou \eng{I do too.} (structure non inversée, plus simple, correcte également).
\end{itemize}
De même, « moi non plus » ne se dit jamais \eng{I not either} ; on dit \eng{Neither do I.} ou \eng{I don't either.}
\end{attention}

\courssub{Les quatre pièges classiques}

\begin{attention}[Piège 1 : oublier l'inversion]
\eng{So I do} existe en anglais, mais il ne veut pas dire « moi aussi » ! Il exprime la \emph{confirmation insistante} : \eng{You never lock the shop. — So I do!} « Tu ne fermes jamais la boutique. — Mais si, je la ferme ! » Pour dire « moi aussi », l'ordre est obligatoirement \eng{So do I} (auxiliaire avant le sujet).
\end{attention}

\begin{attention}[Piège 2 : reprendre le verbe entier]
Après \eng{I like my job}, ne dites pas \eng{So I like} ni \eng{So like I}. L'anglais remplace le verbe lexical par l'auxiliaire \eng{do} : \eng{So do I.} Seul l'auxiliaire voyage dans la réponse courte ; le verbe principal reste derrière.
\end{attention}

\begin{attention}[Piège 3 : le mauvais auxiliaire au passé]
Après \eng{He came to Douala last year}, la réponse \eng{So do I} est fautive : la phrase de départ est au prétérit, l'auxiliaire doit donc être \eng{did} : \eng{So did I.} De même, après \eng{She has opened a shop}, on répond \eng{So have I}, et après \eng{They will strike tomorrow}, \eng{So will the teachers.}
\end{attention}

\begin{attention}[Piège 4 : la double négation]
\eng{Neither} est déjà négatif. Ne dites jamais \eng{Neither don't I} ni \eng{Neither I don't}. La bonne forme est \eng{Neither do I.} Rappelez-vous : un seul signe de négation par phrase anglaise.
\end{attention}

\courssub{La chaîne de décision : trois questions avant de répondre}

\begin{methode}[Souligner d'abord l'auxiliaire]
Avant de produire une réponse courte, suivez toujours ces trois étapes :
\begin{enumerate}
\item \textbf{Lire} la phrase de départ et décider si elle est \emph{affirmative} (\eng{so}) ou \emph{négative} (\eng{neither}).
\item \textbf{Souligner} l'auxiliaire (\eng{is, has, can, will...}) ; s'il n'y en a pas, repérer le temps du verbe et choisir \eng{do/does} (présent) ou \eng{did} (prétérit).
\item \textbf{Construire} : \eng{So/Neither + auxiliaire + sujet}, puis vérifier l'inversion.
\end{enumerate}
Exemple filé : \eng{My father is a trader.} → affirmative → auxiliaire \eng{is} → \eng{So is mine.}
\end{methode}

\begin{popfigure}
\begin{tikzpicture}[
node distance=0.55cm,
startstop/.style={rectangle, rounded corners, draw=popblue, fill=popblueL, text width=4.6cm, align=center, minimum height=0.9cm, font=\small},
decision/.style={diamond, draw=poporange, fill=poporangeL, aspect=2.2, align=center, font=\small, inner sep=1pt},
result/.style={rectangle, rounded corners, draw=popgreen, fill=popgreenL, text width=4.2cm, align=center, minimum height=0.9cm, font=\small},
arrow/.style={-{Stealth[length=3mm]}, thick, popdark}
]
\node[startstop, align=center] (start){Phrase de départ\\ \eng{My father is a trader.}};
\node[decision, below=of start] (aff) {Phrase affirmative ?};
\node[decision, below=of aff] (aux) {Quel auxiliaire ?};
\node[result, below=of aux, align=center] (resp){\eng{So is mine.}\\ « Le mien aussi. »};
\node[result, right=1.2cm of aff, text width=3.4cm] (neg) {Phrase négative :\\ \eng{Neither is mine.}};
\draw[arrow] (start) -- (aff);
\draw[arrow] (aff) -- node[right, font=\small]{oui} (aux);
\draw[arrow] (aff) -- node[above, font=\small]{non} (neg);
\draw[arrow] (aux) -- node[right, font=\small]{\eng{be} $\to$ \eng{is}} (resp);
\end{tikzpicture}
\legende{Chaîne de décision pour construire une réponse courte : polarité, auxiliaire, inversion.}
\end{popfigure}

\courssub{Tableau récapitulatif des formes}

\begin{center}
\begin{tabular}{|l|l|l|}
\hline
\rowcolor{popblueL} Phrase de départ & Réponse « aussi » & Réponse « non plus » \\
\hline
\eng{I am a tailor.} & \eng{So am I.} & --- \\
\hline
\eng{She works in Buea.} & \eng{So does he.} & --- \\
\hline
\eng{They sold cocoa.} & \eng{So did we.} & --- \\
\hline
\eng{He has retired.} & \eng{So has she.} & --- \\
\hline
\eng{We will strike.} & \eng{So will they.} & --- \\
\hline
\eng{I can sew.} & \eng{So can my brother.} & --- \\
\hline
\eng{I am not rich.} & --- & \eng{Neither am I.} \\
\hline
\eng{She doesn't drive.} & --- & \eng{Neither do I.} \\
\hline
\eng{He didn't come.} & --- & \eng{Neither did she.} \\
\hline
\eng{They haven't paid.} & --- & \eng{Neither have we.} \\
\hline
\eng{We can't wait.} & --- & \eng{Neither can they.} \\
\hline
\end{tabular}
\end{center}

\courssub{Exercices d'application}

\begin{exoresolu}[Compléter les réponses courtes]
Énoncé — Complétez chaque réponse avec l'auxiliaire correct.
\begin{enumerate}
\item \eng{My uncle grows cocoa. — So \ldots\ my aunt.}
\item \eng{I didn't attend the meeting. — Neither \ldots\ I.}
\item \eng{The bank has closed. — So \ldots\ the post office.}
\item \eng{She can speak English fluently. — So \ldots\ her husband.}
\item \eng{We are not paid well. — Neither \ldots\ we.} (à corriger : la phrase est déjà donnée, trouvez la bonne forme)
\end{enumerate}
\tcblower
Solution détaillée :
\begin{enumerate}
\item Phrase affirmative au présent simple, verbe \eng{grows} (3\textsuperscript{e} personne) : auxiliaire \eng{does}. → \eng{So does my aunt.}
\item Phrase négative au prétérit (\eng{didn't}) : \eng{Neither did I.}
\item Auxiliaire présent : \eng{has} ; sujet singulier : \eng{So has the post office.}
\item Auxiliaire modal \eng{can} repris tel quel : \eng{So can her husband.}
\item La phrase 5 telle quelle est incorrecte ; la bonne réponse à \eng{We are not paid well} est \eng{Neither are we.} (auxiliaire \eng{are}, jamais de double négation).
\end{enumerate}
\end{exoresolu}

\begin{exercice}[Transformer et corriger]
\textbf{A.} Transformez chaque phrase avec \eng{also} en réponse courte inversée.
\begin{enumerate}
\item \eng{I also work at the market.} (après \eng{My brother works at the market.})
\item \eng{I also went to Bamenda.} (après \eng{The traders went to Bamenda.})
\item \eng{I can also drive a bendskin.} (après \eng{My cousin can drive a bendskin.})
\end{enumerate}
\textbf{B.} Corrigez les erreurs.
\begin{enumerate}
\item \eng{He came late. — So do I.}
\item \eng{I don't like queues. — Neither don't I.}
\item \eng{She is a nurse. — So I am.}
\item \eng{They have finished. — So have finished we.}
\end{enumerate}
\end{exercice}

\begin{corrige}[Exercice : Transformer et corriger]
\textbf{A.} 1. \eng{So do I.} 2. \eng{So did I.} 3. \eng{So can I.}

\textbf{B.} 1. \eng{So did I.} (prétérit → \eng{did}). 2. \eng{Neither do I.} (une seule négation). 3. \eng{So am I.} (inversion obligatoire). 4. \eng{So have we.} (on ne reprend que l'auxiliaire, pas le participe).
\end{corrige}

\courssub{Expression orale : mini-dialogues sur les métiers}

\begin{experience}[Qui fait comme moi ?]
En binômes, chaque élève écrit trois phrases vraies ou inventées sur les métiers et la vie économique (deux affirmatives, une négative), par exemple \eng{My mother sells tomatoes at Mokolo market.} ou \eng{I don't want to be a taxi driver.} L'élève A lit sa phrase ; l'élève B doit réagir immédiatement avec une réponse courte correcte (\eng{So does mine.} / \eng{Neither do I.}), puis ajouter une phrase personnelle. On inverse ensuite les rôles. Le professeur circule et sanctionne uniquement deux fautes : l'absence d'inversion et le mauvais auxiliaire.
\end{experience}

\begin{textebox}[At the vocational training centre in Yaoundé]
\eng{— I want to become an electrician, says André.\\
— So do I, replies Chantal. My brother works for the electricity company, and so does my cousin.\\
— My father didn't finish secondary school, but he has a good job.\\
— Neither did mine, says André, yet he employs ten workers in his carpentry workshop.\\
— I can't save money at the end of the month, sighs Chantal.\\
— Neither can I! Everything is expensive nowadays. But the manager told us that salaries will increase next year.\\
— So will prices, I'm afraid, laughs André.}
\end{textebox}

Glossaire : \emph{electrician} (n., électricien) ; \emph{to employ} (v., employer) ; \emph{carpentry workshop} (n., atelier de menuiserie) ; \emph{to save money} (v., économiser de l'argent) ; \emph{nowadays} (adv., de nos jours) ; \emph{salary} (n., salaire) ; \emph{to increase} (v., augmenter).

Questions de compréhension : 1. What does Chantal's cousin do? 2. Why does André mention his father? 3. What good news does the manager announce?

\begin{aretenir}[L'essentiel de la section]
\begin{itemize}
\item « Moi aussi » = \eng{So + auxiliaire + sujet} ; « moi non plus » = \eng{Neither + auxiliaire + sujet}.
\item L'auxiliaire se \emph{recopie} depuis la phrase de départ ; sans auxiliaire, on utilise \eng{do/does/did}.
\item Inversion obligatoire ; une seule négation ; jamais le verbe entier dans la réponse.
\item À l'écrit, préférez toujours \eng{So do I} à \eng{Me too}.
\end{itemize}
\end{aretenir}

\newpage

\coursec{Activité d'intégration}

\coursec{Lesson 21 — Grammar: Adjectives as Nouns; Short Answers (\eng{So do I}, \eng{Neither did he})}

\courssub{Objectifs de la leçon}

À l'issue de cette leçon, tu seras capable de :

\begin{itemize}
\item identifier et utiliser correctement les \cle{adjectifs employés comme noms} (\eng{the poor}, \eng{the unemployed}, \eng{the young}, \eng{the disabled}, \eng{the rich}, \eng{the elderly}, \eng{the blind}) pour désigner des groupes de personnes ;
\item construire des \cle{réponses courtes} d'accord (\eng{So do I}, \eng{So am I}, \eng{So did he}) ou de désaccord (\eng{Neither do I}, \eng{Neither am I}, \eng{Neither did he}) en choisissant l'auxiliaire adéquat ;
\item éviter les erreurs classiques des francophones (pluriel en \emph{-s}, ordre des mots, auxiliaire inadapté, confusion \eng{So do I} / \eng{So I do}) ;
\item mobiliser ces structures dans une expression orale fluide (débat, discussion) et dans une production écrite d'opinion.
\end{itemize}

\courssub{Observation : les adjectifs employés comme noms (\eng{the + adjective})}

\begin{textebox}[Mountain Voice FM — Programme Announcement]
Good morning, dear listeners! This is Mountain Voice FM, your community radio in Buea. This Saturday, our programme \emph{Speak Out!} asks an important question: “Who needs help in our town?”

The unemployed are everywhere around us: young graduates sell phone credit in the streets because they cannot find jobs. The poor queue in front of microfinance banks every morning. The disabled often beg near the market because nobody employs them. And what about the elderly? Many of them have no pension at all.

Our guests this week are Mr. Tabi, a trader at Buea market; Mrs. Njie, a teacher at Government High School; Kevin, a young graduate; and our journalist, Comfort Eposi. They will debate and propose solutions. Do you agree with them? Call us and say “So do I!” — or tell us why you don't!

\emph{Speak Out!}, Saturday at ten o'clock, only on Mountain Voice FM.
\end{textebox}

\textbf{Glossaire :} \emph{listeners} (n., auditeurs) ; \emph{graduates} (n., diplômés) ; \emph{to queue} (v., faire la queue) ; \emph{to beg} (v., mendier) ; \emph{pension} (n., retraite, pension) ; \emph{trader} (n., commerçant) ; \emph{microfinance banks} (n., institutions de microfinance).

\begin{aretenir}[Règle : \eng{the + adjectif} = nom de groupe]
\begin{itemize}
\item \textbf{Forme :} \cle{the} + \textbf{adjectif} (sans \emph{-s}, sans article indéfini).
\item \textbf{Sens :} désigne une \textbf{catégorie de personnes} prise collectivement (groupe social, physique, moral).
\item \textbf{Accord :} le verbe qui suit est toujours au \textbf{pluriel}.
\item \textbf{Exemples courants :}
\end{itemize}

\begin{center}
\begin{tabularx}{\linewidth}{|X|X|X|}\hline
\rowcolor{popblueL} Structure & Traduction & Exemple \\ \hline
\eng{the poor} & les pauvres & \eng{The poor need support.} \\ \hline
\eng{the rich} & les riches & \eng{The rich should pay more taxes.} \\ \hline
\eng{the young} & les jeunes & \eng{The young are the future.} \\ \hline
\eng{the elderly} & les personnes âgées & \eng{The elderly deserve respect.} \\ \hline
\eng{the unemployed} & les chômeurs & \eng{The unemployed are increasing.} \\ \hline
\eng{the disabled} & les handicapés & \eng{The disabled face many barriers.} \\ \hline
\eng{the blind} / \eng{the deaf} & les aveugles / les sourds & \eng{The blind use Braille.} \\ \hline
\eng{the homeless} & les sans-abri & \eng{The homeless sleep rough.} \\ \hline
\end{tabularx}
\end{center}

\end{aretenir}

\begin{attention}[Pièges fréquents : francophones]
\begin{itemize}
\item \textbf{Interdit d'ajouter \emph{-s} :} \eng{the poors} $\times$, \eng{the unemployeds} $\times$, \eng{the elders} $\times$ (sauf nom \eng{elders} = chefs traditionnels).
\item \textbf{Verbe au pluriel obligatoire :} \eng{The poor \textbf{need} help} (pas \eng{needs}).
\item \textbf{Pas d'article indéfini :} on ne dit pas \eng{a poor} pour « un pauvre » (on dit \eng{a poor person}).
\item \textbf{Adjectifs de nationalité :} \eng{the English}, \eng{the French}, \eng{the Chinese} (invariables) mais \eng{the Americans}, \eng{the Germans} (pluriel en \emph{-s} car ce sont des noms propres adjectivés).
\end{itemize}
\end{attention}

\courssub{Comparaison français / anglais}

\begin{center}
\begin{tabularx}{\linewidth}{|X|X|X|}\hline
\rowcolor{popblueL} Français & English & Remarque \\ \hline
Les pauvres ont faim. & \eng{The poor are hungry.} & \eng{the + adj}, verbe pluriel \\ \hline
Les jeunes étudient. & \eng{The young study.} & Pas de \emph{-s} à \eng{young} \\ \hline
Un pauvre / une pauvre & \eng{A poor person / man / woman} & Il faut ajouter un nom \\ \hline
Les handicapés & \eng{The disabled} & Préféré à \eng{the handicapped} (daté) \\ \hline
\end{tabularx}
\end{center}

\courssub{Observation : les réponses courtes (\eng{So / Neither + auxiliaire + sujet})}

Dans l'annonce radio, l'animatrice dit : \eng{“Call us and say ‘So do I!’”}.
Regarde ces échanges tirés du débat (section 6) :
\begin{itemize}
\item \eng{A: The poor need help. \quad B: So do I.} (accord \emph{positif})
\item \eng{A: I cannot find a job. \quad B: Neither can I.} (accord \emph{négatif})
\item \eng{A: He lost his job. \quad B: So did she.} (accord \emph{passé})
\item \eng{A: We must act. \quad B: So must we.} (accord \emph{modal})
\end{itemize}

\begin{aretenir}[Règle : structure des réponses courtes]
\begin{itemize}
\item \textbf{Accord avec une idée positive :} \cle{So} + \textbf{auxiliaire} + \textbf{sujet}.
\item \textbf{Accord avec une idée négative :} \cle{Neither} (ou \eng{Nor} plus formel) + \textbf{auxiliaire} + \textbf{sujet}.
\item \textbf{L'auxiliaire est \emph{identique} à celui de la phrase de départ} (ou \eng{do/does/did} s'il n'y en a pas).
\item \textbf{Ordre :} Auxiliaire \textbf{avant} le sujet (inversion).
\end{itemize}
\end{aretenir}

\begin{propriete}[Tableau récapitulatif des auxiliaires]
\begin{center}
\begin{tabularx}{\linewidth}{|X|X|X|X|}\hline
\rowcolor{popblueL} Phrase de départ (exemple) & Auxiliaire requis & Accord (\eng{So...}) & Désaccord (\eng{Neither...}) \\ \hline
\eng{I work hard.} (\textsc{present simple}) & \eng{do / does} & \eng{So do I.} & \eng{Neither do I.} \\ \hline
\eng{She works hard.} & \eng{does} & \eng{So does he.} & \eng{Neither does he.} \\ \hline
\eng{They worked hard.} (\textsc{past simple}) & \eng{did} & \eng{So did we.} & \eng{Neither did we.} \\ \hline
\eng{I am tired.} (\eng{be}) & \eng{am / is / are} & \eng{So am I.} & \eng{Neither am I.} \\ \hline
\eng{He was late.} (\eng{be} passé) & \eng{was / were} & \eng{So was she.} & \eng{Neither was she.} \\ \hline
\eng{We can help.} (modal) & \eng{can / could / must...} & \eng{So can they.} & \eng{Neither can they.} \\ \hline
\eng{They have finished.} (parfait) & \eng{have / has / had} & \eng{So have I.} & \eng{Neither have I.} \\ \hline
\end{tabularx}
\end{center}
\end{propriete}

\begin{attention}[Erreurs fatales à éviter]
\begin{itemize}
\item \textbf{Ordre des mots :} \eng{So do I} $\checkmark$ \quad \eng{So I do} $\times$ (sauf confirmation : \eng{“You do work hard!” — “So I do.”} = « C'est vrai que je travaille dur »).
\item \textbf{Auxiliaire incohérent :} \eng{He didn't come. — Neither \textbf{doesn't} she} $\times$ $\rightarrow$ \eng{Neither \textbf{did} she} $\checkmark$.
\item \textbf{\eng{Me too} / \eng{Me neither} :} acceptables à l'oral très familier, mais \textbf{interdits à l'écrit formel} et dans un débat structuré (examen). Utilise \eng{So do I} / \eng{Neither do I}.
\item \textbf{Négation double :} \eng{I don't like it. — Neither do I.} (pas \eng{Neither don't I}).
\item \textbf{Sujet pronominal obligatoire :} \eng{So do I} (pas \eng{So do}).
\end{itemize}
\end{attention}

\courssub{Entraînement guidé : transformation}

\begin{exercice}[Mets à la forme courte d'accord (\eng{So...}) ou de désaccord (\eng{Neither...})]
\begin{enumerate}
\item \eng{The unemployed need jobs.} (Moi aussi) $\rightarrow$ \dots
\item \eng{The rich don't pay enough tax.} (Moi non plus) $\rightarrow$ \dots
\item \eng{My brother lost his job.} (Ma sœur aussi) $\rightarrow$ \dots
\item \eng{We can't solve this alone.} (Eux non plus) $\rightarrow$ \dots
\item \eng{I am worried about the future.} (Moi aussi) $\rightarrow$ \dots
\item \eng{The government didn't help the poor.} (Les ONG non plus) $\rightarrow$ \dots
\item \eng{She has found a job.} (Lui aussi) $\rightarrow$ \dots
\item \eng{You must vote.} (Nous aussi) $\rightarrow$ \dots
\end{enumerate}
\end{exercice}

\begin{corrige}[Exercice 4]
\begin{enumerate}
\item \eng{So do I.} (Présent simple, \eng{need} $\rightarrow$ \eng{do})
\item \eng{Neither do I.} (Négatif présent simple, \eng{don't} $\rightarrow$ \eng{do})
\item \eng{So did she.} (Prétérit, \eng{lost} $\rightarrow$ \eng{did})
\item \eng{Neither can they.} (Modal \eng{can't} $\rightarrow$ \eng{can})
\item \eng{So am I.} (Verbe \eng{be} $\rightarrow$ \eng{am})
\item \eng{Neither did they.} (Prétérit négatif \eng{didn't} $\rightarrow$ \eng{did})
\item \eng{So has he.} (Present perfect \eng{has found} $\rightarrow$ \eng{has})
\item \eng{So must we.} (Modal \eng{must} $\rightarrow$ \eng{must})
\end{enumerate}
\end{corrige}

\courssub{Production orale intégrée : Débat radio \eng{“Speak Out!”}}

\begin{methode}[Déroulement de l'activité (2 h)]
\begin{enumerate}
\item \textbf{Préparation (20 min) :} En groupes de 4. Lisez l'annonce (texte ci-dessus). Soulignez les 5 adjectifs-noms (\eng{the unemployed}, \eng{the poor}, \eng{the disabled}, \eng{the elderly}, \eng{the young} implicite). Identifiez la réponse courte (\eng{So do I}).
\item \textbf{Distribution des rôles (5 min) :} \eng{Host} (animateur), \eng{Trader} (commerçant), \eng{Teacher} (enseignant), \eng{Graduate} (jeune diplômé).
\item \textbf{Rédaction du brouillon (30 min) :} Chaque élève rédige ses répliques (6--8 phrases). Contraintes obligatoires :
      \begin{itemize}
      \item Minimum \textbf{4 adjectifs-noms différents} par groupe (répartis entre les rôles).
      \item Minimum \textbf{4 réponses courtes variées} par groupe (\eng{So do I}, \eng{Neither can I}, \eng{So did he}, \eng{So must we}, \eng{Neither did they}, \eng{So am I}...).
      \item Vocabulaire du chapitre : \eng{unemployment, salary, pension, loan, employer, charity, to earn, to support}.
      \end{itemize}
\item \textbf{Répétition / Auto-correction (20 min) :} Jouez le débat sans lire. Vérifiez : \eng{the + adj} (pas de \emph{-s}), verbe pluriel, auxiliaire correct dans la réponse courte, inversion (\eng{So do I}).
\item \textbf{Représentation \& Grille d'écoute (30 min) :} Chaque groupe joue (3 min). La classe remplit la grille.
\item \textbf{Bilan collectif (15 min) :} Vote du meilleur débat. Écriture au tableau d'une phrase-bilan synthétique.
\end{enumerate}
\end{methode}

\begin{methode}[Grille d'écoute de l'auditeur — À photocopier ou recopier]
Pendant le débat, coche chaque structure correcte entendue. Objectif : 8 structures minimum par groupe.

\begin{center}
\begin{tabularx}{\linewidth}{|X|X|X|}\hline
\rowcolor{popblueL} Structure attendue & Exemple relevé & Correct ? \\ \hline
Adjectif-nom 1 & \eng{the poor} & oui / non \\ \hline
Adjectif-nom 2 & \eng{the unemployed} & oui / non \\ \hline
Adjectif-nom 3 & \eng{the elderly} & oui / non \\ \hline
Adjectif-nom 4 & \eng{the disabled} & oui / non \\ \hline
Réponse courte 1 & \eng{So do I} & oui / non \\ \hline
Réponse courte 2 & \eng{Neither can I} & oui / non \\ \hline
Réponse courte 3 & \eng{So did he} & oui / non \\ \hline
Réponse courte 4 & \eng{So must we} & oui / non \\ \hline
\end{tabularx}
\end{center}
\end{methode}

\courssub{Production modèle et analyse (Corrigé commenté)}

\begin{exoresolu}[Script modèle du débat \eng{“Who needs help in our town?”}]
Énoncé : Rédige et joue un débat de 3 minutes à 4 voix, avec $\ge$ 4 adjectifs-noms et $\ge$ 4 réponses courtes variées.

\tcblower

\textbf{Script du débat (Mountain Voice FM — Speak Out!)}

\eng{\textbf{Host (Comfort):} Good morning, dear listeners, and welcome to Speak Out! Today we ask: who needs help in our town? Mr. Tabi, you are a trader at Buea market. What do you think?}

\eng{\textbf{Trader:} Thank you, Comfort. In my opinion, \textbf{the poor} need help first. Every morning I see them at the market: they cannot even buy a bag of rice. The government does not do enough for them.}

\eng{\textbf{Teacher:} \textbf{So do I}. I agree completely. But let me add \textbf{the young}. My students finish school, and then they find no jobs. \textbf{The unemployed} are more and more numerous in Buea.}

\eng{\textbf{Graduate:} \textbf{Neither can I} find a job! I finished university two years ago, and \textbf{so did} my friend Peter. \textbf{Neither of us has} work. We sell phone credit near the post office. \textbf{The young} are suffering, and \textbf{so are the elderly}. My grandmother has no pension at all.}

\eng{\textbf{Trader:} \textbf{Neither does} my mother. She worked all her life, and now she depends on me.}

\eng{\textbf{Host:} And what about \textbf{the disabled}? Nobody seems to speak about them.}

\eng{\textbf{Teacher:} You are right. \textbf{The disabled} cannot enter many buildings, and employers refuse to give them work. We must change this.}

\eng{\textbf{Graduate:} \textbf{So must we all}. I propose that the council creates small loans for young traders.}

\eng{\textbf{Trader:} \textbf{So do I}! And \textbf{the rich} should pay more taxes to help the poor.}

\eng{\textbf{Host:} Thank you, dear guests. Dear listeners, \textbf{the poor}, \textbf{the young}, \textbf{the elderly} and \textbf{the disabled} need us. Call us now and give your opinion!}

\textbf{Analyse linguistique du modèle :}

\begin{itemize}
\item \textbf{Adjectifs-noms (6) :} \eng{the poor}, \eng{the young}, \eng{the unemployed}, \eng{the elderly}, \eng{the disabled}, \eng{the rich}. Tous précédés de \eng{the}, sans \emph{-s}, verbe au pluriel (\eng{need}, \eng{are}, \eng{cannot}, \eng{should}).
\item \textbf{Réponses courtes (6) :}
      \begin{enumerate}
      \item \eng{So do I} (Teacher $\leftarrow$ Trader: \eng{does not do} $\rightarrow$ \eng{do}).
      \item \eng{Neither can I} (Graduate $\leftarrow$ idée négative \eng{cannot} / \eng{find no jobs} $\rightarrow$ \eng{can}).
      \item \eng{so did} (Graduate: \eng{finished} $\rightarrow$ \eng{did}).
      \item \eng{so are} (Graduate: \eng{are suffering} $\rightarrow$ \eng{are}).
      \item \eng{Neither does} (Trader $\leftarrow$ \eng{has no pension} sens négatif $\rightarrow$ \eng{does}).
      \item \eng{So must we all} (Graduate $\leftarrow$ \eng{must change} $\rightarrow$ \eng{must}).
      \item \eng{So do I} (Trader $\leftarrow$ \eng{I propose} $\rightarrow$ \eng{do}).
      \end{enumerate}
\item \textbf{Lexique économique :} \eng{job, pension, employers, loans, taxes, traders, salary, unemployment}.
\item \textbf{Naturalness :} L'enseignante utilise d'abord \eng{So do I} (court), puis développe. Le diplômé utilise \eng{Neither of us has} (accord de proximité / formel).
\end{itemize}
\end{exoresolu}

\courssub{Bilan de la leçon}

\begin{aretenir}[Ce qu'il faut retenir pour l'examen]
\begin{itemize}
\item \textbf{Adjectif-nom :} \cle{the + adjectif} = groupe de personnes $\rightarrow$ \textbf{verbe au pluriel}, \textbf{pas de \emph{-s}}. Ex. : \eng{The unemployed are desperate.}
\item \textbf{Réponse courte d'accord :} \cle{So + auxiliaire + sujet}. L'auxiliaire = celui de la phrase déclenchante (\eng{do/does/did} si verbe simple).
\item \textbf{Réponse courte de désaccord (accord au négatif) :} \cle{Neither + auxiliaire + sujet}.
\item \textbf{Ordre :} \textbf{Inversion} obligatoire (Auxiliaire -- Sujet). \eng{So do I} $\neq$ \eng{So I do}.
\item \textbf{Contexte :} Ces structures sont la marque d'un anglais \textbf{fluide, naturel, argumentatif} (épreuve orale, dissertation \eng{Opinion essay}).
\end{itemize}
\end{aretenir}

\begin{experience}[Prolongement : Courrier des auditeurs (Écriture individuelle)]
Rédige un message (60--80 mots) à \emph{Mountain Voice FM} pour réagir à l'émission. Contraintes :
\begin{itemize}
\item Au moins \textbf{2 adjectifs-noms}.
\item Au moins \textbf{2 réponses courtes} (\eng{So...} / \eng{Neither...}).
\item Vocabulaire du chapitre.
\end{itemize}
\textbf{Modèle :}
\eng{Dear Comfort, I listened to Speak Out! last Saturday. I agree with the teacher: \textbf{the young} and \textbf{the unemployed} need urgent help. \textbf{So do I} think the council should create loans. \textbf{Neither the government nor the rich do} enough. \textbf{So must we} act together. Yours, [Ton nom].}

\textbf{Correction du modèle :} \eng{Neither the government nor the rich do enough} (proximité \eng{rich} pluriel $\rightarrow$ \eng{do}). \eng{So must we} (court, pas \eng{So must we act}).
\end{experience}

\begin{savaistu}[Culture : \eng{The "Have-Nots"}]
En anglais journalistique et sociologique, on oppose souvent \eng{the haves} (les nantis, ceux qui ont) et \eng{the have-nots} (les démunis, ceux qui n'ont pas). Ces deux noms figés dérivent du verbe \eng{have}. On trouve aussi \eng{the haves and the have-nots} pour parler de la fracture sociale. Ex. : \eng{The gap between the haves and the have-nots is widening.}
\end{savaistu}

\newpage

\coursec{Méthodes et exercices résolus}

\courssub{Fiche méthode 1 : Identifier et construire « the + adjective »}

\begin{methode}[Construire un adjectif employé comme nom]
Pour parler d'une catégorie de personnes en général, l'anglais utilise la structure \cle{the + adjective}, sans ajouter de nom.

\begin{enumerate}
\item Repérer dans l'énoncé un adjectif qui désigne une qualité ou une situation humaine : \eng{rich}, \eng{poor}, \eng{young}, \eng{old}, \eng{unemployed}, \eng{disabled}, \eng{homeless}, \eng{blind}, \eng{dead}.
\item Placer \textbf{the} devant l'adjectif : \eng{the rich}, \eng{the unemployed}.
\item Ne jamais ajouter \eng{people} ni de pluriel en \textbf{-s} : on dit \eng{the poor}, jamais \eng{the poors} ni \eng{the poor people} quand on parle de la catégorie entière.
\item Accorder le verbe au \textbf{pluriel} : \eng{The unemployed \textbf{need} help.} (« Les chômeurs ont besoin d'aide. »)
\item Pour parler d'\textbf{une seule personne}, revenir à un nom ordinaire : \eng{a poor man}, \eng{an unemployed worker}.
\end{enumerate}

Cas particuliers : quelques adjectifs de nationalité en \textbf{-sh, -ch, -ese} fonctionnent de la même façon : \eng{the French}, \eng{the English}, \eng{the Chinese}. Et \eng{the + adjective} peut aussi désigner une chose abstraite, au \textbf{singulier} cette fois : \eng{the unknown} (« l'inconnu »), \eng{The best is yet to come.} (« Le meilleur est à venir. »)
\end{methode}

\begin{propriete}[Règle : sens et accord de « the + adjective »]
\begin{center}
\begin{tabular}{|l|l|l|}
\hline
\rowcolor{popblueL} Structure & Sens & Accord du verbe \\
\hline
\eng{the + adjective} (personnes) & une catégorie sociale : « les pauvres » & \textbf{pluriel} : \eng{The poor suffer.} \\
\hline
\eng{the + adjective} (abstrait) & une notion : « l'inconnu » & \textbf{singulier} : \eng{The unexpected happens.} \\
\hline
\eng{the + nationalité en -sh/-ch/-ese} & un peuple : « les Français » & \textbf{pluriel} : \eng{The French love food.} \\
\hline
\end{tabular}
\end{center}
\end{propriete}

\begin{exoresolu}[Transformation : adjectifs employés comme noms]
Énoncé — Rewrite each sentence using \emph{the + adjective}. Example : \eng{People who are rich should pay more taxes.} $\rightarrow$ \eng{The rich should pay more taxes.}

1. \eng{People who are poor suffer most during an economic crisis.}
2. \eng{Young people in Cameroon are full of energy and ideas.}
3. \eng{People without jobs are looking for opportunities in Douala and Yaoundé.}
4. \eng{Old people deserve our respect and care.}
5. \eng{People who cannot see learn to read with Braille.}
\tcblower
Solution détaillée.

\textbf{Raisonnement général :} chaque phrase contient un groupe « people who are / people without + adjectif ». On supprime le nom \eng{people} et la relative, on place \textbf{the} devant l'adjectif, et on garde le verbe au pluriel.

1. \eng{The poor suffer most during an economic crisis.} — \eng{poor} est l'adjectif ; \eng{the poor} = « les pauvres » ; le verbe \eng{suffer} reste au pluriel (sans \textbf{-s} à la 3\textsuperscript{e} personne).

2. \eng{The young in Cameroon are full of energy and ideas.} — « les jeunes » ; attention, pas de \eng{the youngs}.

3. \eng{The unemployed are looking for opportunities in Douala and Yaoundé.} — « les chômeurs » ; \eng{unemployed} est un participe passé employé comme adjectif, la structure reste identique.

4. \eng{The old deserve our respect and care.} — « les personnes âgées » ; on peut aussi dire \eng{the elderly}, plus poli et plus fréquent dans la presse.

5. \eng{The blind learn to read with Braille.} — « les aveugles ».

\textbf{Conclusion :} la structure \eng{the + adjective} remplace tout le groupe « les + nom de personnes + qui... » du français ; elle est toujours suivie d'un verbe au pluriel et ne prend jamais de \textbf{-s}.
\end{exoresolu}

\courssub{Fiche méthode 2 : Employer « the + adjective » dans un texte de compréhension}

\begin{methode}[Repérer et exploiter la structure dans un texte]
\begin{enumerate}
\item À la lecture, souligne chaque occurrence de \eng{the + adjective} et demande-toi : désigne-t-elle un \textbf{groupe de personnes} (verbe au pluriel) ou une \textbf{notion abstraite} (verbe au singulier : \eng{The unexpected happened.}) ?
\item Pour répondre aux questions de compréhension, reformule : \eng{the homeless} = \eng{people who have no home}.
\item Dans ta propre production (réponse rédigée, composition), utilise la structure pour éviter les répétitions de \eng{people} et montrer ton niveau de langue.
\item Vérifie toujours l'accord du verbe : c'est le piège préféré des correcteurs.
\end{enumerate}
\end{methode}

\begin{textebox}[A social worker in Bamenda]
Every Saturday, Mary volunteers at a centre in Bamenda. The centre feeds the hungry and gives clothes to the homeless. “The elderly come here for company, and the sick receive free medicine,” she explains. “The rich often forget that the poor are their neighbours. But the young are changing things: many students now help the disabled in their free time.”
\end{textebox}

\begin{exoresolu}[Compréhension de texte + grammaire]
Énoncé — Read the text above and answer the questions.

1. Who does the centre help? (Answer with three groups.)
2. Find in the text four examples of \emph{the + adjective} and translate them into French.
3. True or False? Justify : \eng{The sick receive expensive medicine.}
4. Grammar : why is there no \textbf{-s} on \eng{forget} in “\eng{The rich often forget}”?
\tcblower
Solution détaillée.

1. \eng{The centre helps the hungry, the homeless, the elderly, the sick and the disabled.} (Trois groupes suffisent ; on reprend les structures du texte.)

2. Quatre exemples et traductions :
\begin{itemize}
\item \eng{the hungry} = « les affamés / ceux qui ont faim » ;
\item \eng{the homeless} = « les sans-abri » ;
\item \eng{the elderly} = « les personnes âgées » ;
\item \eng{the sick} = « les malades » (on pouvait aussi citer \eng{the rich} = « les riches », \eng{the poor} = « les pauvres », \eng{the young} = « les jeunes », \eng{the disabled} = « les handicapés »).
\end{itemize}

3. \eng{False.} Le texte dit : « \eng{the sick receive free medicine} » ; \eng{free} signifie « gratuit », pas « cher ». Réponse rédigée : \eng{False. The text says the sick receive free medicine, not expensive medicine.}

4. \eng{The rich} désigne un \textbf{groupe de personnes} (« les riches ») : c'est un sujet \textbf{pluriel}. Au présent simple, la 3\textsuperscript{e} personne du pluriel ne prend pas de \textbf{-s} : \eng{The rich forget}, comme \eng{they forget}.

\textbf{Conclusion :} dans un texte, \eng{the + adjective} est un raccourci élégant pour désigner des catégories sociales ; le verbe qui suit se met au pluriel.
\end{exoresolu}

\courssub{Fiche méthode 3 : Construire « So do I » / « Neither did he »}

\begin{methode}[Répondre brièvement : exprimer l'accord]
Pour éviter de répéter toute une phrase, l'anglais utilise des réponses courtes avec \cle{so} (accord après phrase affirmative) ou \cle{neither/nor} (accord après phrase négative).

\begin{enumerate}
\item \textbf{Lire le signe de la phrase} : affirmative $\rightarrow$ \eng{So + auxiliaire + sujet} ; négative $\rightarrow$ \eng{Neither/Nor + auxiliaire + sujet}.
\item \textbf{Identifier l'auxiliaire de la phrase} et le reprendre :
\begin{itemize}
\item auxiliaire \eng{be} $\rightarrow$ \eng{So am I. / Neither is she.}
\item auxiliaire \eng{have} $\rightarrow$ \eng{So have we.}
\item modal (\eng{can, will, must...}) $\rightarrow$ \eng{So can I. / Neither will he.}
\item verbe ordinaire au présent $\rightarrow$ \eng{do/does} : \eng{So does she.}
\item verbe ordinaire au passé $\rightarrow$ \eng{did} : \eng{Neither did he.}
\end{itemize}
\item \textbf{Inverser} sujet et auxiliaire (comme dans une question) : \eng{So do I}, jamais \eng{So I do} pour exprimer l'accord.
\item Choisir le \textbf{pronom sujet} qui correspond à la personne qui répond : \eng{So do I} (« moi aussi »), \eng{Neither did he} (« lui non plus »).
\end{enumerate}

Remarque : \eng{So I do} (sans inversion) existe mais signifie « c'est vrai que je le fais ! » (confirmation insistante) — ce n'est pas la structure d'accord.
\end{methode}

\begin{propriete}[Tableau récapitulatif : le choix de l'auxiliaire]
\begin{center}
\begin{tabular}{|l|l|l|}
\hline
\rowcolor{popblueL} Phrase de départ & Accord (affirmatif) & Accord (négatif) \\
\hline
\eng{I am tired.} & \eng{So am I.} & \eng{I am not. $\rightarrow$ Neither am I.} \\
\hline
\eng{She likes tea.} & \eng{So do I. / So does he.} & \eng{She doesn't. $\rightarrow$ Neither do I.} \\
\hline
\eng{They went out.} & \eng{So did we.} & \eng{They didn't. $\rightarrow$ Neither did we.} \\
\hline
\eng{He has finished.} & \eng{So have I.} & \eng{He hasn't. $\rightarrow$ Neither have I.} \\
\hline
\eng{We can swim.} & \eng{So can she.} & \eng{We can't. $\rightarrow$ Neither can she.} \\
\hline
\eng{You will come.} & \eng{So will they.} & \eng{You won't. $\rightarrow$ Neither will they.} \\
\hline
\end{tabular}
\end{center}
\end{propriete}

\begin{exoresolu}[Réponses courtes : choix de l'auxiliaire]
Énoncé — Agree with each statement using \emph{so} or \emph{neither} + the correct auxiliary.

1. \eng{I like football.} (I)
2. \eng{She doesn't eat meat.} (her brother)
3. \eng{They went to the market yesterday.} (we)
4. \eng{He is a hard-working trader.} (his wife)
5. \eng{We haven't finished our homework.} (they)
6. \eng{My father can speak three languages.} (my mother)
\tcblower
Solution détaillée.

1. Phrase \textbf{affirmative}, verbe ordinaire \eng{like} au présent $\rightarrow$ auxiliaire \eng{do} : \eng{So do I.} (« Moi aussi. »)

2. Phrase \textbf{négative} (\eng{doesn't}) $\rightarrow$ \eng{neither} + \eng{does} : \eng{Neither does her brother.} (« Son frère non plus. »)

3. Affirmative, verbe ordinaire au \textbf{passé} (\eng{went}) $\rightarrow$ auxiliaire \eng{did} : \eng{So did we.} (« Nous aussi. »)

4. Affirmative avec \eng{be} au présent (\eng{is}) $\rightarrow$ on reprend \eng{be} : \eng{So is his wife.} (« Sa femme aussi. »)

5. Négative avec \eng{have} (\eng{haven't}) $\rightarrow$ \eng{Neither have they.} (« Eux non plus. »)

6. Affirmative avec le modal \eng{can} $\rightarrow$ \eng{So can my mother.} (« Ma mère aussi. »)

\textbf{Conclusion :} la clé est toujours la même : signe de la phrase (so/neither), puis reprise fidèle de l'auxiliaire (ou \eng{do/does/did} pour les verbes ordinaires), avec inversion du sujet.
\end{exoresolu}

\courssub{Fiche méthode 4 : Traduire « moi aussi / lui non plus » sans calquer le français}

\begin{methode}[Éviter le calque du français]
\begin{enumerate}
\item Repère la pensée française : « moi aussi » = accord avec une phrase \textbf{affirmative} $\rightarrow$ \eng{so} ; « moi non plus » = accord avec une phrase \textbf{négative} $\rightarrow$ \eng{neither}.
\item Ne traduis jamais « moi » par \eng{me} dans cette structure : l'anglais exige le \textbf{pronom sujet} : \eng{So do I}, pas \eng{So do me}.
\item Ne traduis jamais « aussi » par \eng{also} ou \eng{too} dans la réponse courte : \eng{Me too} est accepté à l'oral familier, mais à l'écrit scolaire on exige \eng{So do I}.
\item Pour le \textbf{désaccord}, utilise la structure opposée : \eng{— I love tea. — Oh, I don't! I prefer coffee.} (« — J'adore le thé. — Pas moi ! Je préfère le café. »)
\end{enumerate}

\begin{center}
\begin{tabularx}{\linewidth}{|X|X|X|}
\hline
\rowcolor{popblueL} Français & English & Remarque \\
\hline
— J'aime le café. — Moi aussi. & — I like coffee. — So do I. & \eng{do} reprend le verbe ordinaire \\
\hline
— Je ne fume pas. — Moi non plus. & — I don't smoke. — Neither do I. & \eng{neither} + auxiliaire positif \\
\hline
— Je suis fatigué. — Moi aussi. & — I am tired. — So am I. & on reprend \eng{be} \\
\hline
— Il est allé à Douala. — Elle aussi. & — He went to Douala. — So did she. & passé $\rightarrow$ \eng{did} \\
\hline
\end{tabularx}
\end{center}
\end{methode}

\begin{exoresolu}[Dialogue à compléter : réponses courtes dans une conversation]
Énoncé — Complete the dialogue between two students in Buea with \emph{so/neither + auxiliary + subject}.

\eng{— I enjoyed the match on Saturday. — \ldots\ldots\ldots\ldots\ldots\ldots\ (I). The Lions played well!}

\eng{— I didn't understand the economics lesson. — \ldots\ldots\ldots\ldots\ldots\ldots\ (Peter). It was too fast.}

\eng{— I am going to study commerce at university. — \ldots\ldots\ldots\ldots\ldots\ldots\ (we).}

\eng{— I can't swim. — \ldots\ldots\ldots\ldots\ldots\ldots\ (my sister).}

\eng{— I have never visited Bamenda. — \ldots\ldots\ldots\ldots\ldots\ldots\ (they).}
\tcblower
Solution détaillée.

1. \eng{So did I.} — phrase affirmative au passé (\eng{enjoyed}) $\rightarrow$ auxiliaire \eng{did}. « Moi aussi. »

2. \eng{Neither did Peter.} — phrase négative au passé (\eng{didn't understand}) $\rightarrow$ \eng{neither + did}. « Pierre non plus. »

3. \eng{So are we.} — phrase affirmative avec \eng{be going to} ; l'auxiliaire est \eng{be} : avec \eng{we}, on prend \eng{are}. « Nous aussi. »

4. \eng{Neither can my sister.} — phrase négative avec le modal \eng{can} $\rightarrow$ \eng{neither + can}. « Ma sœur non plus. »

5. \eng{Neither have they.} — phrase négative au present perfect (\eng{have never visited}) $\rightarrow$ \eng{neither + have}. « Eux non plus. »

\textbf{Conclusion :} dans un dialogue, la réponse courte doit toujours « photographier » l'auxiliaire exact de la phrase précédente : temps, modal et signe.
\end{exoresolu}

\courssub{Fiche méthode 5 : Produire un court texte avec les deux structures}

\begin{methode}[Rédiger un paragraphe qui utilise la leçon]
\begin{enumerate}
\item \textbf{Planifie} : choisis un thème du chapitre (vie économique, métiers, solidarité) : par exemple le travail des jeunes, le chômage, un marché.
\item Place au moins \textbf{deux structures} \eng{the + adjective} avec verbe au pluriel : \eng{The unemployed need training. The young want jobs.}
\item Ajoute un mini-échange avec \textbf{réponse courte} : \eng{— I want to help my community. — So do I.}
\item \textbf{Relis} en vérifiant : pas de \textbf{-s} sur l'adjectif-nom, verbe au pluriel après \eng{the + adjective}, bon auxiliaire dans les réponses courtes, inversion correcte.
\end{enumerate}
\end{methode}

\begin{exoresolu}[Composition guidée]
Énoncé — Write a short paragraph (5 to 7 lines) about young people and work in your town. Use at least two examples of \emph{the + adjective} and one short answer with \emph{so} or \emph{neither}.
\tcblower
Solution détaillée — modèle de rédaction.

\textbf{Raisonnement :} on choisit le thème « les jeunes et l'emploi à Douala ». On place \eng{the young} et \eng{the unemployed} (verbes au pluriel), puis un échange bref avec \eng{So do I}.

\textbf{Réponse rédigée :}

\eng{In Douala, the young are full of ambition: many of them dream of starting their own businesses. Unfortunately, the unemployed often lack training and capital. The government and private companies should work together to help them. My friend Sandra says, “I want to create jobs for my community.” So do I, because the future of our country belongs to those who work hard and help others.}

Traduction : « À Douala, les jeunes sont pleins d'ambition : beaucoup rêvent de créer leur propre entreprise. Malheureusement, les chômeurs manquent souvent de formation et de capital. Le gouvernement et les entreprises privées devraient travailler ensemble pour les aider. Mon amie Sandra dit : “Je veux créer des emplois pour ma communauté.” Moi aussi, car l'avenir de notre pays appartient à ceux qui travaillent dur et aident les autres. »

\textbf{Vérification :} \eng{the young are} (pluriel, correct) ; \eng{the unemployed lack} (pluriel, correct) ; \eng{So do I} reprend le verbe ordinaire \eng{want} au présent avec inversion (correct).

\textbf{Conclusion :} deux structures de la leçon suffisent à enrichir un paragraphe simple et à montrer la maîtrise du point de grammaire.
\end{exoresolu}

\begin{attention}[Les 5 erreurs qui coûtent des points au Bac]
\begin{enumerate}
\item \textbf{Le pluriel interdit} : écrire \eng{the poors}, \eng{the youngs}, \eng{the unemployeds}. L'adjectif-nom ne prend \textbf{jamais} de \textbf{-s} : \eng{the poor}, \eng{the unemployed}.
\item \textbf{L'accord du verbe} : écrire \eng{The rich \textbf{is}...}. Après \eng{the + adjective} désignant des personnes, le verbe est au \textbf{pluriel} : \eng{The rich \textbf{are}...}.
\item \textbf{Le calque « moi aussi »} : écrire \eng{Me too am} ou \eng{I also} en réponse courte. La structure exigée est \eng{So do I / So am I / So did I}, avec \textbf{inversion} et le \textbf{pronom sujet} (jamais \eng{me}).
\item \textbf{Le mauvais auxiliaire} : répondre \eng{So do I} à \eng{I went to Buea.} Au passé, il faut \eng{did} : \eng{So did I.} ; à \eng{I am tired.}, il faut \eng{So am I.} L'auxiliaire doit toujours correspondre à celui de la phrase.
\item \textbf{La double négation} : écrire \eng{Neither don't I.} Après \eng{neither/nor}, l'auxiliaire est à la forme \textbf{affirmative} : \eng{Neither do I.} (« Moi non plus. ») — la négation est déjà portée par \eng{neither}.
\end{enumerate}
\end{attention}

\newpage

\coursec{Exercices}

\courssub{Je vérifie mes connaissances}

\begin{exercice}[QCM : la bonne réponse courte]
Pour chaque phrase de départ, choisis la bonne réponse courte parmi les quatre options proposées.
\begin{enumerate}
\item \eng{“I work in a bank in Douala.” — \ldots}
\begin{enumerate}
\item \eng{So I work.} \item \eng{So do I.} \item \eng{So am I.} \item \eng{Neither do I.}
\end{enumerate}
\item \eng{“John didn't pass the entrance examination.” — \ldots}
\begin{enumerate}
\item \eng{So didn't Peter.} \item \eng{Neither Peter didn't.} \item \eng{Neither did Peter.} \item \eng{So did Peter.}
\end{enumerate}
\item \eng{“My sister is a nurse.” — \ldots}
\begin{enumerate}
\item \eng{So is mine.} \item \eng{So does mine.} \item \eng{Neither is mine.} \item \eng{So mine is.}
\end{enumerate}
\item \eng{“We have never visited Buea.” — \ldots}
\begin{enumerate}
\item \eng{So have we.} \item \eng{Neither we have.} \item \eng{Neither have we.} \item \eng{So we haven't.}
\end{enumerate}
\item \eng{“She can speak three languages.” — \ldots}
\begin{enumerate}
\item \eng{So can her brother.} \item \eng{So does her brother.} \item \eng{Neither can her brother.} \item \eng{So her brother can.}
\end{enumerate}
\end{enumerate}
\end{exercice}

\begin{exercice}[Vrai ou faux ? Justifie]
Indique si chaque phrase est correcte (true) ou incorrecte (false), puis justifie ta réponse en une phrase en français.
\begin{enumerate}
\item \eng{The poors need our help.}
\item \eng{The young often prefer living in big cities like Yaoundé.}
\item \eng{“I like cocoa farming.” — “So I do.”}
\item \eng{“He didn't come to work yesterday.” — “Neither did she.”}
\item \eng{The rich are not always happy.}
\item \eng{“My father is a teacher.” — “Neither is mine.”}
\end{enumerate}
\end{exercice}

\begin{exercice}[Vocabulaire : les adjectifs-noms]
Associe chaque adjectif-nom anglais à sa traduction française.
\begin{center}
\begin{tabularx}{\linewidth}{|X|X|}
\hline
\rowcolor{popblueL} \textbf{English} & \textbf{Français} \\
\hline
\eng{the unemployed} & a. les aveugles \\
\hline
\eng{the blind} & b. les chômeurs \\
\hline
\eng{the elderly} & c. les blessés \\
\hline
\eng{the injured} & d. les personnes âgées \\
\hline
\eng{the homeless} & e. les sans-abri \\
\hline
\eng{the deaf} & f. les sourds \\
\hline
\end{tabularx}
\end{center}
\end{exercice}

\begin{exercice}[Complète avec le bon auxiliaire]
Complète chaque réponse courte avec l'auxiliaire qui convient (\eng{do}, \eng{does}, \eng{did}, \eng{is}, \eng{are}, \eng{have}, \eng{has}, \eng{can}, \eng{will}, \eng{would}).
\begin{enumerate}
\item \eng{“I enjoy trading at the market.” — “So \ldots\ my mother.”}
\item \eng{“They didn't attend the meeting.” — “Neither \ldots\ we.”}
\item \eng{“She has worked in this company for ten years.” — “So \ldots\ her husband.”}
\item \eng{“I am interested in economics.” — “So \ldots\ I.”}
\item \eng{“We won't buy a new car this year.” — “Neither \ldots\ they.”}
\item \eng{“He can repair any engine.” — “So \ldots\ his apprentice.”}
\end{enumerate}
\end{exercice}

\courssub{Je m'entraîne}

\begin{exercice}[Traduction]
Traduis les phrases suivantes en anglais.
\begin{enumerate}
\item Les pauvres ont besoin de notre aide.
\item « J'habite à Bamenda. » — « Moi aussi. »
\item « Mon frère n'aime pas son métier. » — « Mon frère non plus. »
\item Les jeunes doivent respecter les personnes âgées.
\item « Elle a trouvé un emploi à Douala. » — « Ma sœur aussi. »
\item Les riches ne comprennent pas toujours les problèmes des pauvres.
\end{enumerate}
\end{exercice}

\begin{exercice}[Transformation : « the + adjective »]
Réécris chaque phrase en utilisant la structure \cle{the + adjective}, sans changer le sens.
\begin{enumerate}
\item \eng{Poor people need our help.} $\rightarrow$ \eng{The \ldots}
\item \eng{Young people like modern music.} $\rightarrow$ \eng{The \ldots}
\item \eng{Unemployed people are looking for jobs in the city.} $\rightarrow$ \eng{The \ldots}
\item \eng{Old people should receive a pension.} $\rightarrow$ \eng{The \ldots}
\item \eng{Blind people can learn to read with Braille.} $\rightarrow$ \eng{The \ldots}
\item \eng{Rich people often live in big houses.} $\rightarrow$ \eng{The \ldots}
\end{enumerate}
\end{exercice}

\begin{exercice}[Dialogue à trous]
Deux amis, Tabi and Ngong, parlent de leurs métiers. Complète le dialogue avec six réponses courtes utilisant \eng{so} ou \eng{neither} + auxiliaire.
\begin{textebox}[At the motor park]
\eng{Tabi: I work as a mechanic near the motor park.\\
Ngong: Really? (1) \ldots\ my cousin. He repairs lorries.\\
Tabi: I started this job five years ago.\\
Ngong: (2) \ldots\ my cousin. He finished his training in 2019.\\
Tabi: But I don't earn much money.\\
Ngong: (3) \ldots\ he. Life is difficult for mechanics.\\
Tabi: I am planning to open my own garage one day.\\
Ngong: (4) \ldots\ my cousin! He talks about it every day.\\
Tabi: I can't save enough money for the moment.\\
Ngong: (5) \ldots\ he. Everything is expensive now.\\
Tabi: Well, we will succeed one day.\\
Ngong: (6) \ldots\ we, my friend. We must be patient.}
\end{textebox}
\end{exercice}

\begin{exercice}[Correction d'erreurs]
Chaque phrase contient une erreur typique des francophones. Souligne l'erreur et corrige-la.
\begin{enumerate}
\item \eng{The poors live in this neighbourhood.}
\item \eng{“I like my job.” — “So I do.”}
\item \eng{“She didn't come.” — “Neither don't I.”}
\item \eng{The youngs prefer football to handball.}
\item \eng{“My father is a driver.” — “So does mine.”}
\item \eng{The richs should pay more taxes.}
\item \eng{“We haven't finished.” — “Neither haven't they.”}
\item \eng{“He works hard.” — “So works his brother.”}
\end{enumerate}
\end{exercice}

\courssub{Je me prépare au Bac}

\begin{exercice}[Type Bac — Compréhension et langue]
Lis le texte suivant, puis réponds aux questions en anglais, par des phrases complètes.
\begin{textebox}[Work and solidarity in a Cameroonian town]
\eng{In many towns of Cameroon, economic life is hard for a large part of the population. The unemployed spend their days looking for small jobs, while the young dream of travelling abroad to find better opportunities. The old, who have worked all their lives on farms or in offices, often receive very small pensions, and the sick cannot always pay for their treatment.}\\\eng{However, solidarity is still strong. The rich are sometimes generous: last year, a businessman from Douala offered free school fees to two hundred children. The disabled also receive help from associations which train them to sew, to make shoes or to repair telephones. “The strong must help the weak,” says Mama Rose, a trader at the central market. “When I earn money, I share with my neighbours. My sister does the same, and so do many women of this market.”}
\end{textebox}
\textbf{Comprehension questions}
\begin{enumerate}
\item According to the text, what do the unemployed do all day?
\item Why do the young dream of travelling abroad?
\item What did the businessman from Douala do last year?
\item How do associations help the disabled?
\item Find in the text one example of \eng{the + adjective} used as a noun and translate it into French.
\end{enumerate}
\textbf{Language work}
\begin{enumerate}
\item Transform: \eng{Old people often receive very small pensions.} (Use \eng{the + adjective}.)
\item Complete: \eng{“My sister shares with her neighbours.” — “So \ldots\ many women of this market.”}
\item Complete: \eng{“The sick cannot always pay for their treatment.” — “Neither \ldots\ the unemployed.”}
\end{enumerate}
\end{exercice}

\begin{exercice}[Type Bac — Traduction]
Traduis le passage suivant en anglais. (Environ 60 mots.)
\begin{quote}
Dans ma ville, les chômeurs sont de plus en plus nombreux. Les jeunes ne trouvent pas de travail après leurs études, et les vieux ne reçoivent pas toujours leur pension. « J'ai cherché un emploi pendant deux ans », dit mon frère. « Moi aussi », répond son ami. Heureusement, les riches de notre quartier aident parfois les pauvres avec de la nourriture et de l'argent.
\end{quote}
\end{exercice}

\begin{exercice}[Type Bac — Composition]
\textbf{Topic:} \eng{The economic problems of my town and how to solve them.}

Write a composition of \textbf{150 to 180 words}. In your essay, you must:
\begin{itemize}
\item present two or three economic problems of your town (unemployment, poverty, high prices, etc.);
\item use \textbf{at least three} adjectives used as nouns (\eng{the unemployed}, \eng{the poor}, \eng{the young}\ldots);
\item report \textbf{at least two} short answers with \eng{so} or \eng{neither} (for example: \eng{My friend thinks so, and so do I.});
\item conclude with your own suggestion or hope for the future.
\end{itemize}
\end{exercice}

\newpage

\coursec{Corrigés des exercices}

\begin{corrige}[Exercice 1 — QCM : la bonne réponse courte]
\begin{enumerate}
\item \textbf{b.} \eng{So do I.} — Phrase affirmative au présent simple avec un verbe ordinaire (\eng{work}) : on reprend l'auxiliaire \eng{do}. \eng{So I work} est un calque du français « moi je travaille aussi » ; \eng{am} ne convient pas car \eng{work} n'est pas le verbe \eng{be}.
\item \textbf{c.} \eng{Neither did Peter.} — Phrase négative au prétérit (\eng{didn't pass}) : on utilise \eng{neither} + \eng{did} + sujet. \eng{Neither Peter didn't} est une double négation impossible ; \eng{so} exige une phrase affirmative de départ.
\item \textbf{a.} \eng{So is mine.} — Le verbe de départ est \eng{be} (\eng{is}) : on reprend \eng{is}. \eng{Mine} remplace \eng{my sister}. Pas d'auxiliaire \eng{does} avec \eng{be}.
\item \textbf{c.} \eng{Neither have we.} — Phrase négative (\eng{never visited}) au present perfect : \eng{neither} + \eng{have} + sujet. L'inversion est obligatoire : pas de \eng{neither we have}.
\item \textbf{a.} \eng{So can her brother.} — Phrase affirmative avec le modal \eng{can} : on reprend \eng{can}. \eng{Does} est impossible après un modal.
\end{enumerate}
\textbf{Rappel de méthode :} (1) la phrase de départ est-elle affirmative ou négative ? (2) quel est l'auxiliaire (ou le verbe \eng{be}, ou le modal) ? (3) inversion obligatoire : \eng{so/neither + auxiliaire + sujet}.
\end{corrige}

\begin{corrige}[Exercice 2 — Vrai ou faux ? Justifie]
\begin{enumerate}
\item \textbf{False.} \eng{The poors} est incorrect : l'adjectif-nom reste invariable. Correction : \eng{The poor need our help.} (On n'ajoute jamais de \eng{-s} à \eng{the + adjective}.)
\item \textbf{True.} \eng{The young} = « les jeunes », suivi du verbe au pluriel (\eng{prefer}) : la phrase est correcte.
\item \textbf{False.} \eng{So I do} signifie « c'est vrai que je l'aime » (insistance, même sujet). Pour dire « moi aussi », il faut l'inversion : \eng{So do I.}
\item \textbf{True.} Phrase négative au prétérit (\eng{didn't come}) : \eng{Neither did she} est la réponse courte correcte.
\item \textbf{True.} \eng{The rich} + verbe au pluriel (\eng{are}) : correct. « Les riches ne sont pas toujours heureux. »
\item \textbf{False.} La phrase de départ est affirmative (\eng{is a teacher}) : on ne peut pas répondre par \eng{neither}, qui exige une phrase négative. Correction : \eng{So is mine.}
\end{enumerate}
\end{corrige}

\begin{corrige}[Exercice 3 — Vocabulaire : les adjectifs-noms]
\begin{center}
\begin{tabularx}{\linewidth}{|X|X|}
\hline
\rowcolor{popblueL} \textbf{English} & \textbf{Français} \\
\hline
\eng{the unemployed} & b. les chômeurs \\
\hline
\eng{the blind} & a. les aveugles \\
\hline
\eng{the elderly} & d. les personnes âgées \\
\hline
\eng{the injured} & c. les blessés \\
\hline
\eng{the homeless} & e. les sans-abri \\
\hline
\eng{the deaf} & f. les sourds \\
\hline
\end{tabularx}
\end{center}
\textbf{Point de vigilance :} tous ces adjectifs-noms désignent un \emph{groupe de personnes} et se construisent avec un verbe au \textbf{pluriel} : \eng{The unemployed are numerous.} — « Les chômeurs sont nombreux. »
\end{corrige}

\begin{corrige}[Exercice 4 — Complète avec le bon auxiliaire]
\begin{enumerate}
\item \eng{“I enjoy trading at the market.” — “So \textbf{does} my mother.”} — Présent simple, verbe ordinaire, sujet à la 3\textsuperscript{e} personne du singulier (\eng{my mother}) : \eng{does}.
\item \eng{“They didn't attend the meeting.” — “Neither \textbf{did} we.”} — Phrase négative au prétérit : on reprend \eng{did}.
\item \eng{“She has worked in this company for ten years.” — “So \textbf{has} her husband.”} — Present perfect avec \eng{has} (3\textsuperscript{e} personne) : on reprend \eng{has}.
\item \eng{“I am interested in economics.” — “So \textbf{am} I.”} — Verbe \eng{be} au présent, sujet \eng{I} : \eng{am}.
\item \eng{“We won't buy a new car this year.” — “Neither \textbf{will} they.”} — Phrase négative avec \eng{won't} : on reprend \eng{will} à la forme affirmative après \eng{neither}.
\item \eng{“He can repair any engine.” — “So \textbf{can} his apprentice.”} — Modal \eng{can} repris tel quel.
\end{enumerate}
\textbf{Règle :} l'auxiliaire de la réponse courte est toujours celui de la phrase de départ (\eng{be}, \eng{have}, modal), ou \eng{do/does/did} si le verbe de départ est un verbe ordinaire.
\end{corrige}

\begin{corrige}[Exercice 5 — Traduction]
\begin{enumerate}
\item \eng{The poor need our help.} — Pas de \eng{-s} à \eng{poor} ; verbe au pluriel.
\item \eng{“I live in Bamenda.” — “So do I.”} — Présent simple, verbe ordinaire : auxiliaire \eng{do}, inversion obligatoire.
\item \eng{“My brother doesn't like his job.” — “Neither does mine.”} — Phrase négative au présent : \eng{neither does} + \eng{mine} (= \eng{my brother}).
\item \eng{The young must respect the elderly.} (ou \eng{the old}) — Deux adjectifs-noms, invariables.
\item \eng{“She has found a job in Douala.” — “So has my sister.”} — Present perfect : on reprend \eng{has}.
\item \eng{The rich do not always understand the problems of the poor.} — Deux adjectifs-noms au pluriel ; négation avec \eng{do not}.
\end{enumerate}
\textbf{Points de vigilance :} ne jamais traduire « moi aussi » par \eng{me too} dans un registre écrit soutenu avec reprise de l'auxiliaire ; ne jamais écrire \eng{the poors} ou \eng{the youngs}.
\end{corrige}

\begin{corrige}[Exercice 6 — Transformation : « the + adjective »]
\begin{enumerate}
\item \eng{The poor need our help.}
\item \eng{The young like modern music.}
\item \eng{The unemployed are looking for jobs in the city.}
\item \eng{The old should receive a pension.} (ou \eng{The elderly})
\item \eng{The blind can learn to read with Braille.}
\item \eng{The rich often live in big houses.}
\end{enumerate}
\textbf{Méthode :} supprimer \eng{people}, mettre l'adjectif après \eng{the}, garder le verbe au pluriel. L'adjectif ne prend jamais de \eng{-s}.
\end{corrige}

\begin{corrige}[Exercice 7 — Dialogue à trous]
\begin{enumerate}
\item \eng{So does my cousin.} — Phrase affirmative au présent simple (\eng{work}), sujet 3\textsuperscript{e} personne.
\item \eng{So did my cousin.} — Phrase affirmative au prétérit (\eng{started}) : auxiliaire \eng{did}.
\item \eng{Neither does he.} — Phrase négative au présent (\eng{don't earn}) : \eng{neither} + \eng{does}.
\item \eng{So is my cousin!} — Verbe \eng{be} au présent (\eng{am planning} $\rightarrow$ auxiliaire \eng{be}) : \eng{so is}.
\item \eng{Neither can he.} — Phrase négative avec \eng{can't} : on reprend \eng{can}.
\item \eng{So will we, my friend.} — Phrase affirmative au futur (\eng{will succeed}) : \eng{so will we}.
\end{enumerate}
\textbf{Traduction du début :} « Tabi : Je travaille comme mécanicien près du parc de véhicules. — Ngong : Vraiment ? Mon cousin aussi. Il répare des camions. »
\end{corrige}

\begin{corrige}[Exercice 8 — Correction d'erreurs]
\begin{enumerate}
\item \eng{The \emph{poors}} $\rightarrow$ \eng{The \textbf{poor} live in this neighbourhood.} — L'adjectif-nom est invariable.
\item \eng{“So \emph{I do}.”} $\rightarrow$ \eng{“So \textbf{do I}.”} — Pour « moi aussi », inversion obligatoire ; \eng{So I do} exprime l'insistance.
\item \eng{“Neither \emph{don't} I.”} $\rightarrow$ \eng{“Neither \textbf{do} I.”} — Après \eng{neither}, l'auxiliaire est à la forme affirmative (pas de double négation).
\item \eng{The \emph{youngs}} $\rightarrow$ \eng{The \textbf{young} prefer football to handball.} — Jamais de \eng{-s}.
\item \eng{“So \emph{does} mine.”} $\rightarrow$ \eng{“So \textbf{is} mine.”} — Le verbe de départ est \eng{be} (\eng{is}) : on reprend \eng{is}, pas \eng{does}.
\item \eng{The \emph{richs}} $\rightarrow$ \eng{The \textbf{rich} should pay more taxes.} — Adjectif-nom invariable.
\item \eng{“Neither \emph{haven't} they.”} $\rightarrow$ \eng{“Neither \textbf{have} they.”} — Forme affirmative de l'auxiliaire après \eng{neither}.
\item \eng{“So \emph{works} his brother.”} $\rightarrow$ \eng{“So \textbf{does} his brother.”} — On ne répète pas le verbe ordinaire : on le remplace par l'auxiliaire \eng{do/does/did}.
\end{enumerate}
\end{corrige}

\begin{corrige}[Exercice 9 — Type Bac — Compréhension et langue]
\textbf{Comprehension questions} (réponses modèles)
\begin{enumerate}
\item \eng{The unemployed spend their days looking for small jobs.}
\item \eng{The young dream of travelling abroad because they want to find better opportunities.}
\item \eng{Last year, the businessman from Douala offered free school fees to two hundred children.}
\item \eng{Associations help the disabled by training them to sew, to make shoes or to repair telephones.}
\item \eng{Example: “the disabled”} (ou \eng{the unemployed}, \eng{the young}, \eng{the old}, \eng{the sick}, \eng{the rich}, \eng{the strong}, \eng{the weak}) — Traduction : « les handicapés ».
\end{enumerate}
\textbf{Language work}
\begin{enumerate}
\item \eng{The old often receive very small pensions.}
\item \eng{“My sister shares with her neighbours.” — “So \textbf{do} many women \textbf{in} this market.”} — Présent simple, sujet pluriel : \eng{do}.
\item \eng{“The sick cannot always pay for their treatment.” — “Neither \textbf{can} the unemployed.”} — Phrase négative avec \eng{cannot} : on reprend \eng{can}.
\end{enumerate}
\textbf{Barème indicatif (sur 10) :} compréhension 5 pts (1 pt par question, exigeant une phrase complète et correcte) ; language work 3 pts (1 pt par item) ; qualité générale de l'anglais 2 pts.
\textbf{Points de vigilance :} répondre par des phrases complètes (sujet + verbe) ; ne pas recopier le texte sans l'adapter à la question ; verbe au pluriel après \eng{the + adjective}.
\end{corrige}

\begin{corrige}[Exercice 10 — Type Bac — Traduction]
\textbf{Proposition de traduction :}
\begin{quote}
\eng{In my town, the unemployed are more and more numerous. The young do not find work after their studies, and the old do not always receive their pension. “I looked for a job for two years,” says my brother. “So did I,” his friend answers. Fortunately, the rich people of our neighbourhood — the rich — sometimes help the poor with food and money.}
\end{quote}
Version plus élégante pour la fin : \eng{Fortunately, the rich in our neighbourhood sometimes help the poor with food and money.}
\textbf{Barème indicatif (sur 10) :} respect du sens 4 pts ; correction grammaticale (adjectifs-noms invariables, pluriel, \eng{so did I}) 4 pts ; vocabulaire et orthographe 2 pts.
\textbf{Points de vigilance :}
\begin{itemize}
\item « les chômeurs » = \eng{the unemployed} (jamais \eng{the unemployeds}) ;
\item « Moi aussi » après une phrase au prétérit = \eng{So did I} (et non \eng{Me too} dans un texte soigné) ;
\item « pendant deux ans » avec un prétérit simple = \eng{for two years} ;
\item « de plus en plus nombreux » = \eng{more and more numerous}.
\end{itemize}
\end{corrige}

\begin{corrige}[Exercice 11 — Type Bac — Composition]
\textbf{Proposition de rédaction modèle (environ 165 mots) :}
\begin{quote}
\eng{My town faces serious economic problems. First, unemployment is very high: the unemployed spend their days in front of workshops, hoping for small jobs. The young suffer the most, because even after secondary school they cannot find work. My friend Paul finished his studies last year, but he has no job, and neither do many of his classmates.}\\\eng{Secondly, life is very expensive. The poor cannot buy enough food, and the sick often stay at home because they cannot pay for medicine. “Prices rise every month,” says my mother, and so do all the traders of our market.}\\\eng{To solve these problems, the government should build factories and help young people to create small businesses. The rich of our town could also support training centres for the young. I hope that one day, everyone will have a decent job and a better life.}
\end{quote}
\textbf{Barème indicatif (sur 20) :} respect du sujet et de la consigne de longueur 4 pts ; présence d'au moins trois adjectifs-noms corrects 4 pts ; présence d'au moins deux réponses courtes correctes (\eng{so/neither} + auxiliaire) 4 pts ; cohérence et organisation (introduction, problèmes, solutions, conclusion) 4 pts ; correction grammaticale et vocabulaire 4 pts.
\textbf{Points de vigilance :}
\begin{itemize}
\item Vérifier que chaque adjectif-nom est suivi d'un verbe au pluriel : \eng{the poor cannot}, \eng{the young suffer} ;
\item Vérifier l'inversion dans les réponses courtes : \eng{so do all the traders}, \eng{neither do many of his classmates} ;
\item Choisir le bon auxiliaire selon la phrase de départ ;
\item Structurer avec des connecteurs : \eng{first}, \eng{secondly}, \eng{to solve these problems}, \eng{I hope that}.
\end{itemize}
\end{corrige}

\newpage

\coursec{Fiche bilan}

\begin{popfigure}
\begin{tikzpicture}[mindmap, grow cyclic, every node/.style={concept, text=white, align=center, font=\small},
root concept/.append style={concept color=popdark, font=\bfseries},
level 1/.append style={level distance=3.6cm, sibling angle=51.5},
level 2/.append style={level distance=2.2cm}]
\node[root concept, align=center]{Lesson 21\\Grammar}
  child[concept color=popblue]{ node[align=center]{The + adjective\\= a group} 
    child[concept color=popblue!70]{ node[align=center]{the rich,\\the poor} }
  }
  child[concept color=poporange]{ node[align=center]{Nationalities\\-sh, -ch, -ese}
    child[concept color=poporange!70]{ node[align=center]{the French,\\the Chinese} }
  }
  child[concept color=popgreen]{ node[align=center]{Plural verb\\always}
    child[concept color=popgreen!70]{ node[align=center]{The old\\are...} }
  }
  child[concept color=poppurple]{ node[align=center]{So + aux + I\\(agreement +)}
    child[concept color=poppurple!70]{ node[align=center]{So do I,\\So can she} }
  }
  child[concept color=popgold]{ node[align=center]{Neither + aux\\(agreement -)}
    child[concept color=popgold!70]{ node[align=center]{Neither did he,\\Nor am I} }
  }
  child[concept color=poppink]{ node[align=center]{Right auxiliary\\= tense match}
    child[concept color=poppink!70]{ node[align=center]{do / did /\\have / will} }
  }
  child[concept color=popmuted]{ node[align=center]{Francophone\\traps}
    child[concept color=popmuted!70]{ node[align=center]{no -s, no\\“me too”} }
  };
\end{tikzpicture}
\legende{Carte mentale de la leçon : les deux structures clés — \eng{the + adjective} et les réponses courtes \eng{so do I / neither did he}.}
\end{popfigure}

\begin{bilan}
\textbf{1. Lexique clé (English / Français)}
\begin{itemize}
  \item \eng{the rich} : les riches ; \eng{the poor} : les pauvres ; \eng{the young} : les jeunes ; \eng{the old / the elderly} : les personnes âgées
  \item \eng{the unemployed} : les chômeurs ; \eng{the homeless} : les sans-abri ; \eng{the disabled} : les handicapés ; \eng{the blind} : les aveugles
  \item \eng{the French} : les Français ; \eng{the English} : les Anglais ; \eng{the Chinese} : les Chinois ; \eng{the Swiss} : les Suisses
  \item \eng{to agree} : être d'accord ; \eng{so do I} : moi aussi ; \eng{neither do I} : moi non plus
\end{itemize}

\textbf{2. Règles à connaître par cœur}
\begin{center}
\begin{tabularx}{\linewidth}{|X|X|}\hline
\rowcolor{popblueL} \textbf{Règle} & \textbf{Exemple} \\ \hline
\eng{the + adjective} désigne un \textbf{groupe de personnes} (jamais un individu seul) & \eng{The government should help the poor.} « Le gouvernement devrait aider les pauvres. » \\ \hline
L'adjectif reste \textbf{invariable} : jamais de -s, jamais de \eng{ones} & \eng{the richs} $\rightarrow$ FAUX ; \eng{the rich} $\rightarrow$ CORRECT \\ \hline
Le verbe qui suit est toujours au \textbf{pluriel} & \eng{The unemployed are protesting.} « Les chômeurs manifestent. » \\ \hline
Nationalités en \textbf{-sh, -ch, -ese, -iss} : \eng{the + adjective} ; autres : nom pluriel en -s & \eng{the French}, \eng{the Japanese} mais \eng{Cameroonians}, \eng{Americans} \\ \hline
Accord avec une phrase \textbf{affirmative} : \eng{So + auxiliaire + sujet} & \eng{I like cocoa. — So do I.} « Moi aussi. » \\ \hline
Accord avec une phrase \textbf{négative} : \eng{Neither/Nor + auxiliaire + sujet} & \eng{She didn't come. — Neither did he.} « Lui non plus. » \\ \hline
L'auxiliaire \textbf{reprend le temps} de la phrase : do/does (présent), did (prétérit), have/has (present perfect), will, can, am/is/are & \eng{I have finished. — So has Peter.} ; \eng{I am tired. — So am I.} \\ \hline
\end{tabularx}
\end{center}

\textbf{3. Méthode express}
\begin{itemize}
  \item \textbf{Réponse courte en 3 étapes} : (1) repérer affirmatif ou négatif ; (2) repérer l'auxiliaire ou le temps (be, have, modal, sinon do/does/did) ; (3) construire \eng{So/Neither + auxiliaire + nouveau sujet}.
  \item \textbf{Traduction} : « les + adjectif » (les jeunes, les riches) $\rightarrow$ \eng{the + adjective}, sans ajouter \eng{people} sauf pour une personne précise : \eng{a rich man} « un homme riche ».
\end{itemize}
\end{bilan}

\begin{aretenir}[Checklist avant le Bac]
$\square$ Je sais que \eng{the + adjective} désigne un groupe entier de personnes.\\
$\square$ Je n'ajoute jamais de -s à l'adjectif : \eng{the poor}, pas \eng{the poors}.\\
$\square$ Je mets le verbe au pluriel après \eng{the + adjective} : \eng{The blind are...}\\
$\square$ Je connais les nationalités à adjectif (\eng{the French, the English, the Chinese}) et celles à nom en -s (\eng{Cameroonians, Germans}).\\
$\square$ Je sais dire « moi aussi » : \eng{So do I} (jamais \eng{me too} dans une réponse grammaticale complète).\\
$\square$ Je sais dire « moi non plus » : \eng{Neither do I / Nor do I}.\\
$\square$ Je choisis le bon auxiliaire selon le temps : do/does, did, have/has, will, can, be.\\
$\square$ Je respecte l'inversion : \eng{So does she}, jamais \eng{So she does} pour exprimer l'accord.\\
$\square$ Je sais que \eng{neither} est déjà négatif : pas de \eng{not} en plus.\\
$\square$ Je peux transformer un dialogue complet avec \eng{so} et \eng{neither} sans erreur.\\
$\square$ Je peux employer \eng{the + adjective} dans une phrase sur la vie économique : \eng{The government must create jobs for the unemployed.}\\
$\square$ Je vérifie toujours : affirmatif $\rightarrow$ \eng{so} ; négatif $\rightarrow$ \eng{neither/nor}.
\end{aretenir}

\findecours

\end{document}
